% Incognita paper P000, a paper of the owner (AGENTS.md 5.1; CONTRIBUTING.md, section 8). Rules:
% AGENTS.md, section 5 (template preamble, except the font size and the margins of the author's version:
% 10pt and 2.3cm, for which its displays were written; at most 100 pages). Working notes go in notebook/.
% Converted on 2026-11-14 from the author's readable version of 4 November 2026 (statements and
% proofs unchanged; form adapted to the library: sections, blueprint, packages).
% Compile: latexmk -pdf -interaction=nonstopmode -halt-on-error -cd papers/<P>/paper.tex (as CI),
% after `python3 scripts/certify.py bib` (papers/library.bib, the library's papers; AGENTS.md 5.6).
% Blueprint conventions: every definition/lemma/proposition/theorem/corollary has a \label;
% statements start with \uses{def:...} when they use definitions; every proof starts with an
% exhaustive \uses{...}; a proof away from its statement has \proves{...}. Never write \lean,
% \leanok or \mathlibok in this file. After acceptance this file is never edited (AGENTS.md 6.6).
\documentclass[10pt]{article}
\usepackage[T1]{fontenc}
\usepackage[utf8]{inputenc}
\usepackage{lmodern}
\usepackage[a4paper,margin=2.3cm]{geometry}
\usepackage{amsmath,amssymb,mathtools}
\usepackage{xcolor}
\usepackage[colorlinks=true,linkcolor=blue!50!black,citecolor=blue!50!black,urlcolor=blue!50!black]{hyperref}
\usepackage{../../templates/paper/blueprint-macros}

% Macros of this paper: \newcommand only.
% Macros of the paper on HY. Every symbol is a link to its definition: anchor \dt{s:xxx}, macro \sy{s:xxx}{symbol}.
\newcommand{\dt}[1]{\hypertarget{#1}{}}
\newcommand{\sy}[2]{{\hyperlink{#1}{#2}}}
\newcommand{\E}{\mathbb{E}}
\newcommand{\PP}{\mathbb{P}}
\newcommand{\1}{\mathbf{1}}
\newcommand{\KL}{\mathrm{KL}}
\newcommand{\kl}{\mathrm{kl}}
\newcommand{\Ber}{\mathrm{Ber}}
\newcommand{\cA}{\mathcal A}
%---------------------------------------------------------------- model
\newcommand{\m}{\sy{s:model}{\mathbf{m}}}
\newcommand{\mtwo}{\sy{s:model}{\mathbf m_{(2)}}}
\newcommand{\cP}{\sy{s:model}{\mathcal P}}
\newcommand{\Dg}{\sy{s:model}{\Delta}}
\newcommand{\Ki}{\sy{s:kinf}{\mathcal K_{\inf}}}
\newcommand{\kap}{\sy{s:kinf}{\kappa}}
\newcommand{\Lset}[1]{\sy{s:kinf}{\Lambda}(#1)}
\newcommand{\da}{\sy{s:kinf}{d}_a}
\newcommand{\Ng}{\sy{s:count}{N}}
\newcommand{\Reg}{\sy{s:count}{\mathfrak R}}
\newcommand{\hnu}{\sy{s:stack}{\widehat\nu}}
\newcommand{\hm}{\sy{s:stack}{\widehat{\mathbf m}}}
\newcommand{\Xs}{\sy{s:stack}{X}}
\newcommand{\Mz}{\sy{s:M0}{\mathcal M_0}}
%---------------------------------------------------------------- grid
\newcommand{\tl}{\sy{s:grid}{t}}
\newcommand{\Lg}{\sy{s:grid}{\mathsf L}}
\newcommand{\xg}{\sy{s:grid}{x}}
\newcommand{\Bg}{\sy{s:grid}{B}}
%---------------------------------------------------------------- algorithm HY
\newcommand{\HY}{\hyperref[def:hy]{\textup{HY}}}
\newcommand{\wz}{\sy{s:alg}{w_0}}
\newcommand{\yg}{\sy{s:alg}{y}}
\newcommand{\Ug}{\sy{s:alg}{U}}
\newcommand{\tauA}{\sy{s:alg}{\tau}}
\newcommand{\sA}{\sy{s:alg}{s}}
\newcommand{\tnu}{\sy{s:alg}{\widetilde\nu}}
\newcommand{\tm}{\sy{s:alg}{\widetilde{\mathbf m}}}
\newcommand{\tms}{\sy{s:alg}{\widetilde{\mathbf m}^\star}}
\newcommand{\Ld}{\sy{s:alg}{\widehat c}}
\newcommand{\Kt}{\sy{s:alg}{\widetilde K}}
\newcommand{\Wg}{\sy{s:alg}{W}}
\newcommand{\rr}{\sy{s:alg}{r}}
\newcommand{\mh}{\sy{s:alg}{\widehat\mu}}
\newcommand{\Kh}{\sy{s:alg}{\widehat K}}
\newcommand{\Qp}{\sy{s:alg}{\mathcal Q}}
\newcommand{\Lh}{\sy{s:alg}{\widehat\Lambda}}
\newcommand{\hg}{\sy{s:alg}{h}}
\newcommand{\Zg}{\sy{s:alg}{Z}}
\newcommand{\nt}{\sy{s:alg}{n}}
\newcommand{\om}{\sy{s:alg}{\omega}}
\newcommand{\ody}{\sy{s:alg}{\omega^{\mathrm{dy}}}}
\newcommand{\thv}{\sy{s:alg}{\vartheta}}
\newcommand{\gam}{\sy{s:alg}{\gamma}}
\newcommand{\Jg}{\sy{s:alg}{J}}
\newcommand{\Act}{\sy{s:act}{\mathfrak A}}
\newcommand{\HH}{\sy{s:H}{H}}
%---------------------------------------------------------------- tools
\newcommand{\Phis}{\sy{s:tilt}{\Phi}}
\newcommand{\Psis}{\sy{s:tilt}{\Psi}}
\newcommand{\taun}{\sy{s:gridl}{\tau}}
\newcommand{\lamk}{\sy{s:gridl}{\lambda}}
\newcommand{\kn}{\sy{s:gridl}{k}}
%---------------------------------------------------------------- lower bound
\newcommand{\rhom}{\sy{s:tiltb}{\rho^-_s}}
\newcommand{\rhop}{\sy{s:tiltb}{\rho^+_s}}
\newcommand{\QQ}{\sy{s:Q}{\mathbb Q}}
\newcommand{\nzero}{\sy{s:Q}{n_0}}
\newcommand{\Mb}{\sy{s:Q}{M}}
\newcommand{\Sl}{\sy{s:Q}{S}}
\newcommand{\cB}{\sy{s:lbproof}{c_B}}
%---------------------------------------------------------------- x > 1: constants, ranks, events
\newcommand{\eb}{\sy{s:constA}{\bar\varepsilon}}
\newcommand{\eg}{\sy{s:constA}{\eta}}
\newcommand{\Ceta}{\sy{s:constA}{C_\eta}}
\newcommand{\cst}{\sy{s:constA}{c_\star}}
\newcommand{\ew}{\sy{s:constA}{\varepsilon_\omega}}
\newcommand{\De}{\sy{s:constA}{D}}
\newcommand{\Kp}{\sy{s:constA}{K^+}}
\newcommand{\kape}{\sy{s:constA}{\kappa}}
\newcommand{\Rt}{\sy{s:constA}{\widetilde R}}
\newcommand{\rt}{\sy{s:constA}{\widetilde\rho}}
\newcommand{\ngr}{\sy{s:constA}{n}}
\newcommand{\cAc}{\sy{s:constA}{c_A}}
\newcommand{\CLU}{\sy{s:constA}{C_{LU}}}
\newcommand{\hb}{\sy{s:constA}{\bar h}}
\newcommand{\xp}{\sy{s:constA}{x_+}}
\newcommand{\sig}{\sy{s:sums}{\sigma}}
\newcommand{\SE}{\sy{s:SE}{S'}}
\newcommand{\lA}{\sy{s:ranksA}{\ell_A}}
\newcommand{\lone}{\sy{s:ranksA}{\ell_1}}
\newcommand{\ug}{\sy{s:ranksA}{u}}
\newcommand{\Ag}{\sy{s:evA}{A}}
\newcommand{\Vg}{\sy{s:evA}{V}}
\newcommand{\Pig}{\sy{s:evA}{\Pi}}
\newcommand{\Lo}{\sy{s:evA}{\mathrm{Lo}}}
\newcommand{\Up}{\sy{s:evA}{\mathrm{Up}}}
\newcommand{\Gg}{\sy{s:evA}{\Gamma}}
\newcommand{\Fg}{\sy{s:evA}{F}}
\newcommand{\mt}{\sy{s:decA}{m}}
\newcommand{\Yg}{\sy{s:decA}{Y}}
\newcommand{\epl}{\sy{s:dy}{\varepsilon}}
\newcommand{\cone}{\sy{s:dy}{c_1}}
\newcommand{\etwo}{\sy{s:dy}{e_2}}
\newcommand{\lonep}{\sy{s:dy}{\ell'_1}}
%---------------------------------------------------------------- x = 1: constants, events, water level
\newcommand{\ez}{\sy{s:constC}{\varepsilon_0}}
\newcommand{\ezz}{\sy{s:constC}{\eta_0}}
\newcommand{\dl}{\sy{s:constC}{\delta}}
\newcommand{\Cz}{\sy{s:constC}{C_{0}}}
\newcommand{\Co}{\sy{s:constC}{C_1}}
\newcommand{\Dp}{\sy{s:constC}{D'}}
\newcommand{\Me}{\sy{s:constC}{M}}
\newcommand{\Cg}{\sy{s:constC}{C'}}
\newcommand{\Rz}{\sy{s:constC}{R}}
\newcommand{\rz}{\sy{s:constC}{\rho}}
\newcommand{\up}{\sy{s:constC}{u'}}
\newcommand{\yb}{\sy{s:constC}{\bar y}}
\newcommand{\thc}{\sy{s:constC}{\theta_a}}
\newcommand{\ctl}{\sy{s:constC}{c_{\theta}}}
\newcommand{\epth}{\sy{s:constC}{\varepsilon_{\theta}}}
\newcommand{\lzp}{\sy{s:ranksC}{\ell'_0}}
\newcommand{\ly}{\sy{s:ranksC}{\ell_y}}
\newcommand{\lonez}{\sy{s:ranksC}{\ell_L}}
\newcommand{\lz}{\sy{s:ranksC}{\ell_0}}
\newcommand{\Lam}{\sy{s:water}{\Lambda}}
\newcommand{\lW}{\sy{s:water}{\ell_W}}
\newcommand{\Dev}{\sy{s:evC}{\mathcal D}}
\newcommand{\Vs}{\sy{s:evC}{\mathcal V^\star}}
\newcommand{\Va}{\sy{s:evC}{\mathcal V^a}}
\newcommand{\cU}{\sy{s:evC}{\mathcal U}}
\newcommand{\cG}{\sy{s:evC}{\mathcal G}}
\newcommand{\cC}{\sy{s:evC}{\mathcal C}}
\newcommand{\cO}{\sy{s:evC}{\mathcal O}}
\newcommand{\Shi}{\sy{s:good}{S}}
\newcommand{\Sgo}{\sy{s:good}{S^{\mathrm g}}}
\newcommand{\Eamb}{\sy{s:amb}{\mathcal E}}
\newcommand{\lfour}{\sy{s:amb}{\ell_4}}
\newcommand{\zet}{\sy{s:amb}{\zeta}}
\newcommand{\zetp}{\sy{s:amb}{\zeta'}}
\newcommand{\nl}{\sy{s:amb}{n}}
\newcommand{\Acc}{\sy{s:amb}{\mathrm{Acc}}}
\newcommand{\KC}{\sy{s:amb}{\mathrm{KC}}}
\newcommand{\cell}{\sy{s:amb}{\Xi}}
\newcommand{\Cfour}{\sy{s:amb}{C_4}}
\newcommand{\Cpi}{\sy{s:sumC}{C_\pi}}
\newcommand{\RHo}{\sy{s:RHo}{\mathcal R^{1}}}
\newcommand{\lfive}{\sy{s:decC}{\ell_5}}


%---------------------------------------------------------------- readability (rules of redaction-preuves.md)
% setgoal: what is fixed / assumed, then what is to be shown, before any computation.
\newcommand{\Setting}{\noindent\textbf{Setting.}\ }
\newcommand{\Goal}{\par\smallskip\noindent\textbf{Goal.}\ }
\newcommand{\Plan}{\par\smallskip\noindent\textbf{Plan.}\ }
% named steps of a proof
\newcommand{\stp}[2]{\par\medskip\noindent\textbf{Step #1: #2.}\ }
% vocabulary of the paper: a word that is a link to its definition (anchor \dt{t:xxx} or \dt{s:xxx})
\newcommand{\term}[2]{\hyperlink{#1}{#2}}
% role of a lemma in the proof that uses it
\newcommand{\role}[1]{\par\noindent{\small\textbf{Role.}\ #1}\par}
% justification written on the line of an aligned chain
\newcommand{\why}[1]{\text{\small\color{black!70}(#1)}}
%---------------------------------------------------------------- statuses and Lean (hidden by \statutfalse)
\setlength{\emergencystretch}{3em}

\title{Batched bandits on an unknown imposed grid:\\ the optimal constant $x^2/\mathcal K_{\inf}$ for every $x\ge1$}
\author{Adrien Prevost\thanks{Submitted by its author to the Incognita library (roadmap BatchedBandits, item F4) on
  2026-11-14. The proofs were checked by AI agents only, not by a human referee. The conversion to the format of the
  library was done with AI assistance (Claude Opus 5.5). The author verified the main statements in a private Lean~4
  project (Subsection~\ref{sec:lean}).}}
\date{2026-11-14}

\begin{document}
\sloppy
\maketitle

\begin{abstract}
We study $K$-armed stochastic bandits with rewards in $[0,1]$ whose feedback arrives in batches whose ends $t_1<t_2<\cdots$ are imposed by the environment and are
\emph{not known in advance}: at a round of batch $\ell+1$ the learner knows $t_1,\dots,t_\ell$ and the current round, but not when the batch will end. We measure the
regret at the ends of the batches, and assume only that $x_\ell:=\log t_\ell/\log t_{\ell-1}$ converges to some $x\in[1,\infty)$. We show that the optimal first-order
constant is $x^2/\mathcal K_{\inf}(\nu_a,\mathbf m^\star)$ for every $x\ge1$. The lower bound holds for every consistent batched learner, even one that knows the grid
(for the $\limsup$; for the $\liminf$ when $x=1$ or when the counts are logarithmic): the exploration needed at the end of batch $\ell+2$ must be done before the end of
batch $\ell$. The upper bound is attained by a single algorithm, HY, which knows nothing about the grid nor about the instance and is optimal \emph{simultaneously} on every
grid with $x_\ell\to x$, including $x=1$ with batches of arbitrary sizes (its remainder terms depend on the grid). HY combines three ingredients: a paced exploration plan that
follows $x_\ell^2\log t$ at every round, an IMED phase at the start of every long batch, and a consistency threshold that compares the evidence against an arm
with its own evidence one test earlier. HY is a proof of existence: its
remainder terms are astronomical.
\end{abstract}

\section{Main result}\label{sec:main}\label{sec:intro}

\paragraph{Motivation.} In many sequential decision problems the feedback is not available after every decision but by blocks: clinical trials run in phases, agronomic experiments give
their yields once per season, recommender systems are updated once per day. The \emph{batched bandit} model captures this: the learner chooses arms at every round, but
observes the rewards only at the end of each batch. Most of the literature lets the learner choose the batches, and asks how few batches suffice
\cite{PRCS16,GHRZ19,JTXHXG21}, \cite{EKMM21} (to be checked). We consider the opposite situation, where the batches are \emph{imposed} by the environment
and are not even known in advance: the learner discovers the end of a batch when it happens.

\paragraph{Setting in a nutshell.} The batch ends $t_1<t_2<\cdots$ are deterministic and fixed by the environment. In the information model $\Mz$ that we consider, at a round
$t\in(t_\ell,t_{\ell+1}]$ of batch $\ell+1$ the learner knows $t_1,\dots,t_\ell$, the current round $t$, its past actions and the rewards of the rounds $1,\dots,t_\ell$; it
knows neither $t_{\ell+1}$ nor any later end, nor any horizon. The rewards are in $[0,1]$, with no other assumption on their distributions than $\m^\star<1$. Since the
learner cannot react inside a batch, the natural criterion is the regret \emph{at the ends of the batches}, $\Reg(t_\ell)$. The growth of the grid is measured by
\[
x_\ell:=\frac{\log t_\ell}{\log t_{\ell-1}},
\]
and we assume that $x_\ell\to x\in[1,\infty)$. Geometric, polynomial and bounded batch sizes give $x=1$; doubly exponential grids $t_\ell\approx t_1^{x^\ell}$, as used
by Perchet et al.~\cite{PRCS16} and Gao et al.~\cite{GHRZ19}, give $x>1$.


\begin{theorem}[the optimal constant is $x^2/\Ki$, attained without knowing the grid]\label{thm:main}
\uses{def:model, def:grid, def:learner, def:kinf, def:hy, def:M0}
Let $x\in[1,\infty)$, $\nu\in\cP$ and $a$ a sub-optimal arm. On every grid with $\xg_\ell\to x$:
\begin{enumerate}
\item[(a)] every consistent batched learner, even one that knows the grid, satisfies $\limsup_\ell\E_\nu[\Ng_a(\tl_\ell)]/\log\tl_\ell\ge x^2/\Ki(\nu_a,\m^\star)$, and the same
for the $\liminf$ if $x=1$ or if its counts are logarithmic at the ends of the batches;
\item[(b)] the single learner \HY$(\wz,\beta,\ody,\thv,\gam)$ of $\Mz$, which knows neither the grid nor the instance, is consistent and has
$\lim_\ell\E_\nu[\Ng_a(\tl_\ell)]/\log\tl_\ell=x^2/\Ki(\nu_a,\m^\star)$.
\end{enumerate}
Hence $x^2/\Ki(\nu_a,\m^\star)$ is the optimal constant at the ends of the batches, for every $x\ge1$, and one learner attains it simultaneously on every grid with
$\xg_\ell\to x$, $x=1$ included (the remainder terms of Theorem~\ref{thm:upper} depend on the grid). The same holds for the regret, with the constant
$C_x:=x^2\sum_{c:\Dg_c>0}\Dg_c/\Ki(\nu_c,\m^\star)$: every consistent batched learner has $\limsup_\ell\Reg(\tl_\ell)/\log\tl_\ell\ge C_x$, and \HY{} has
$\lim_\ell\Reg(\tl_\ell)/\log\tl_\ell=C_x$.
\end{theorem}


\paragraph{Constants.} The first-order constant $x^2/\Ki(\nu_a,\m^\star)$ is explicit. The remainder terms behind the limits
(Propositions~\ref{prop:fixedC} and~\ref{prop:dyA}) do not depend on the rank $\ell$; they depend on the instance, on the
grid and on an accuracy $\varepsilon$, and they are astronomically large: the theorem describes the first-order behaviour
only. Since they depend on the grid, the convergence is not uniform over the grids with $\xg_\ell\to x$.
The remainders are \emph{existential}: they are written as formulas, but these formulas go through ranks defined as
``the smallest rank such that'' a condition holds from that rank on (for example $\lfive(\varepsilon)$ in
Proposition~\ref{prop:fixedC}, and $\lW$ of Lemma~\ref{lem:W}), whose existence is proved but which are not computed;
in the author's Lean statements they are existential too.

\paragraph{Roadmap items.} This paper serves item F4 of the roadmap BatchedBandits, whose Appendix~A states (verbatim,
items~1 and~2 of ``Statements''; there $L_\ell=\log t_\ell$, $d_a=K_{\inf}(\nu_a,m^\star)$, and $y_k$, $\omega$, $\varepsilon_0$
are the objects of \HY{} and of its proof):
\begin{quote}\small
On every grid with $x_\ell \to x \in [1, \infty)$, every $\nu$ with rewards in $[0,1]$ and
$m^\star < 1$, every suboptimal arm $a$:
\begin{enumerate}
\item \textbf{Lower bound, ``explore two batches ahead''.} Every batched learner that is consistent at batch
   ends, even one that knows the grid, has
   \begin{itemize}
   \item if $x = 1$: $\liminf_\ell \mathbb E N_a(t_\ell)/L_\ell \ge 1/d_a$ (this holds on every grid);
   \item if $x > 1$: $\limsup_\ell \mathbb E N_a(t_\ell)/L_\ell \ge x^2/d_a$, and the same with
     $\liminf$ if $\mathbb E N_a(t_\ell) = O(L_\ell)$;
   \item for the regret: $\liminf_\ell \mathrm{Reg}(t_\ell)/L_\ell \ge x^2 \sum_{c:\Delta_c>0} \Delta_c / K_{\inf}(\nu_c, m^\star)$
     when $x = 1$, or when $x > 1$ and every suboptimal count is $O(L_\ell)$.
   \end{itemize}
\item \textbf{Upper bound.} HY with its default parameters (below) is a deterministic learner of $M_0$,
   consistent, with $\lim_\ell \mathbb E N_a(t_\ell)/L_\ell = x^2/d_a$ and the matching regret limit.
   Non-asymptotic forms: for $x = 1$,
   $\mathbb E N_a(t_L) \le L_L/d_a + (2/d_a + 4c^\star/d_a^2)\,\varepsilon L_L + R_1(\varepsilon)$
   for every $\varepsilon \in (0, \varepsilon_0]$ and every $L$, with $c^\star := (2-m^\star)/(1-m^\star)$;
   for $x > 1$, with a constant margin $\omega$,
   $\mathbb E N_a(t_\ell) \le (1+\omega)(1/d_a + 3c^\star\varepsilon/d_a^2)\,H_\ell + R_x(\varepsilon)$
   where $H_\ell := \max_{k \le \ell} y_{k-1} L_k$ and $H_\ell/L_\ell \to x^2$. The remainders do
   not depend on $L$ or $\ell$ and are astronomically large.
\end{enumerate}
\end{quote}
What is proved here, item by item:
\begin{itemize}
\item Item~1 is Theorem~\ref{thm:lower}, (i), (ii) and (iii) respectively; Theorem~\ref{thm:main}(a) and its regret part
restate it.
\item Item~2, first sentence: Theorem~\ref{thm:upper} (\HY{} is deterministic and in $\Mz$ by Definition~\ref{def:hy}
and Properties~(i) of Subsection~\ref{sec:algo}); Theorem~\ref{thm:main}(b) restates it.
\item Item~2, the non-asymptotic form for $x=1$: Proposition~\ref{prop:fixedC}, with $\Co=2/\da+4\cst/\da^2$
(Definition~\ref{def:constC}) and $R_1(\varepsilon)=\RHo(\varepsilon)$.
\item Item~2, the non-asymptotic form for $x>1$ with a constant margin $\omega$ \emph{is not proved in this paper}.
For $x>1$ the paper proves only Proposition~\ref{prop:dyA}, a bound on the $\limsup$ for \HY{} with the default
(dyadic) margins $\ody_\ell\to0$, which is what Theorem~\ref{thm:upper} needs.
\item Item~3 of Appendix~A (the limit over all rounds) \emph{is not proved in this paper}. Proposition~\ref{prop:dyA}
bounds the $\limsup$ over all rounds when $x>1$, but neither the matching $\liminf$ over all rounds nor the case
$x=1$ over all rounds is proved here.
\end{itemize}
The paper also serves item F1 (the one-batch lower bound on every grid): Lemma~\ref{lem:onebatch} with the remark of
Theorem~\ref{thm:lower}(i).


\paragraph{Discussion.} \emph{On the assumption.} The only assumption on the grid is the convergence of $\xg_\ell$; its limit $x$ is not known to \HY, and bounded batch
sizes, batches of size one and short batches alternating with long ones are allowed. Grids for which $\xg_\ell$ does not converge are not covered. \emph{On the distributions.} No parametric assumption is made: the results hold for every instance of $\cP$, that is, for all distributions on $[0,1]$
with $\m^\star<1$, with any number of optimal arms. \emph{On the remainders.} The remainder terms are astronomically large: the theorems describe the first-order behaviour only. \emph{On the information.} The lower bound holds even for
learners that know the grid, and \HY{} knows nothing about it: knowing the grid does not improve the first-order constant at the ends of the batches. \emph{Formal verification.}
The author verified these statements in a private Lean project (Section~\ref{sec:lean}).


\paragraph{Contributions.} We show that the optimal first-order constant of $\E[\Ng_a(t_\ell)]/\log t_\ell$ for a sub-optimal arm $a$ is
\[
\frac{x^2}{\Ki(\nu_a,\m^\star)}\qquad\text{for every }x\in[1,\infty),
\]
where $\Ki$ is the non-parametric divergence of Honda and Takemura~\cite{HT10}, and that it is attained by one algorithm that needs no knowledge of the grid.
\begin{itemize}
\item \emph{Lower bound} (Theorem~\ref{thm:lower}). For every consistent batched learner, even one that knows the whole grid, $\limsup_\ell\E[\Ng_a(t_\ell)]/\log t_\ell\ge x^2/\Ki(\nu_a,\m^\star)$,
and the same holds for the $\liminf$ when $x=1$ or when the counts are logarithmic. The factor $x^2$ means that the learner must explore \emph{two batches ahead}: at
the end of batch $\ell$ it must already have explored enough for the horizon $t_{\ell+2}$, because the allocation of batch $\ell+1$ is decided without the rewards of
batch $\ell$.
\item \emph{Upper bound} (Theorem~\ref{thm:upper}). The algorithm \HY{} (Definition~\ref{def:hy}) satisfies $\lim_\ell\E[\Ng_a(t_\ell)]/\log t_\ell=x^2/\Ki(\nu_a,\m^\star)$ on
every instance with rewards in $[0,1]$ and $\m^\star<1$, for \emph{every} grid with $x_\ell\to x\in[1,\infty)$: bounded batch sizes, batches of size one, short batches
alternating with long ones are all allowed.
\item \emph{Optimality} (Theorem~\ref{thm:main}). The two theorems give the same constant, so that $x^2/\Ki$ is the optimal constant, attained by one learner of $\Mz$
simultaneously on every convergent grid.
\item \emph{Formal verification} (Section~\ref{sec:lean}). The author verified Theorems~\ref{thm:lower} and~\ref{thm:upper} and Theorem~\ref{thm:main} in a
private Lean~4 project, for \HY{} with any tie rules; formalising this paper in the library ports that verification.
\end{itemize}


\paragraph{The algorithm in a nutshell.} \HY{} must explore two batches ahead without knowing where the current batch ends, and must stay reliable on every grid with
$x_\ell\to x$, from bounded batch sizes to doubly exponential grids. It does so with three ingredients (Section~\ref{sec:algo}). \emph{Pacing}: an arm rejected by the test at
$t_\ell$ is kept at $(1+\omega)x_\ell^2\log t/\widehat K$ pulls at every round $t$, so that, whenever the batch ends, its exploration aims at the test that will be made then;
the arms that are not clearly rejected receive a polylogarithmic net. \emph{An IMED phase}: every batch starts with IMED rounds \cite{HT15}, computed with the
statistics frozen at $t_\ell$ and the current counts; a batch that lasts beyond $t_\ell^{1+\vartheta_\ell}$ then switches to the paced plan. The IMED rounds certify a
deterministic ``water level'' that makes the leader reliable when $x=1$, whatever the batch sizes. \emph{A consistency threshold}: an arm is declared ambiguous only if its
evidence has dropped, since the previous test, below the evidence it had then; this keeps the test robust when the IMED rounds lag behind $\log t$.



\subsection{Formal verification}\label{sec:lean}

The author verified Theorems~\ref{thm:lower}, \ref{thm:upper} and~\ref{thm:main} in a private Lean~4 project (Mathlib and
LeanMachineLearning, no \texttt{sorry}, axioms \texttt{propext}, \texttt{Classical.choice}, \texttt{Quot.sound}), for \HY{}
with any fixed tie rules, the remainders being existential. Those Lean statements write $\liminf\ge L$ as ``for every
$c<L$, eventually $c\le u_\ell$'', $\limsup\ge L$ as ``for every $c<L$, $c<u_\ell$ for infinitely many $\ell$'', and limits
as convergences. The Lean proof follows the architecture of the proofs below, but many lemmas are formalized in a different
form, so the proofs are not claimed verified as written. That project is not part of this library: formalising this paper
here (frozen statements, then proofs) is the next step, and only the library's audit can mark a node as verified.


\setcounter{tocdepth}{1}
\tableofcontents

%======================================================================

\section{Setting}\label{sec:setting}

The conventions of the roadmap BatchedBandits, made precise. Every definition below is a node; the algorithm \HY{} is defined in Subsection~\ref{sec:algo}.

\begin{definition}[grid]\label{def:grid}\dt{s:grid}
The batch sizes are deterministic integers $\Bg_\ell\ge1$, $\ell\ge1$, with $\Bg_1\ge\max(2,K)$. The end of batch $\ell$ is
$\tl_\ell:=\Bg_1+\dots+\Bg_\ell$ ($\tl_0:=0$), and batch $\ell$ is the set of rounds $(\tl_{\ell-1},\tl_\ell]$. Let $\Lg_\ell:=\log\tl_\ell>0$ for $\ell\ge1$, and
\[
\xg_\ell:=\frac{\Lg_\ell}{\Lg_{\ell-1}}\ \ge1\quad(\ell\ge2),\qquad\xg_1:=1,\qquad\text{so that}\qquad\tl_\ell=\tl_{\ell-1}^{\,\xg_\ell}\quad(\ell\ge2).
\]
\end{definition}
\begin{definition}[assumption on the grid]\label{def:ass-x}
\uses{def:grid}
$\xg_\ell\to x$ as $\ell\to\infty$, for some $x\in[1,\infty)$. Nothing else is assumed on the grid.
\end{definition}

\begin{definition}[arms and the class $\cP$]\label{def:model}\dt{s:model}
There are $K\ge2$ arms $c\in\cA$; arm $c$ has a reward distribution $\nu_c$ on $[0,1]$ with mean $\m_c$. Let $\m^\star:=\max_c\m_c$,
$\Dg_c:=\m^\star-\m_c$, and, when some arm is sub-optimal, $\mtwo:=\max\{\m_c:\m_c<\m^\star\}$ and $\Dg:=\m^\star-\mtwo$. The class $\cP$ is the set of instances
$\nu=(\nu_c)_{c\in\cA}$ of distributions on $[0,1]$ with $\m^\star<1$. Throughout, $a$ is a fixed sub-optimal arm ($\m_a<\m^\star$) and $a^\star$ a fixed optimal arm.
\end{definition}

\begin{definition}[the divergence $\Ki$]\label{def:kinf}\dt{s:kinf}
\uses{def:model}
For $\mu\in[0,1)$ let $\kap(\mu):=\log\frac1{1-\mu}$ and $\Lset{\mu}:=[0,\frac1{1-\mu})$. For a distribution $G$ on $[0,1]$,
\[
\Ki(G,\mu):=\sup_{\lambda\in\Lset{\mu}}\E_{X\sim G}\big[\log\big(1-\lambda(X-\mu)\big)\big]\ \in[0,\kap(\mu)],
\]
and $\Ki(G,1):=0$ if $\E_G[X]=1$, $+\infty$ otherwise. By Honda and Takemura~\cite[Theorem~8]{HT10} (Lemma~\ref{cite:honda2010asymptotically-primal} and the remark after it), $\Ki(G,\mu)$ is the infimum of $\KL(G,H)$ over the distributions $H$ on
$[0,1]$ with $\E_H[X]>\mu$; the upper bound only uses the dual formula above. For an empirical distribution, the supremum can be taken over rational $\lambda$, so $\Ki$
is a measurable function of the sample. We write $\da:=\Ki(\nu_a,\m^\star)$; it is positive (Lemma~\ref{lem:kinf}(iv)).
\end{definition}

\begin{definition}[batched learners, counts, regret, consistency]\label{def:learner}\dt{s:count}
\uses{def:grid, def:model}
A \emph{batched learner} chooses at round $t$ an arm $A_t$ drawn from a probability kernel (its policy) that may depend on all the
past actions and on the rewards of the rounds $s\le\tl_\ell$ when $t\in(\tl_\ell,\tl_{\ell+1}]$, but not on the rewards of the current batch; the reward $R_t$ has
conditional law $\nu_{A_t}$ given the past and $A_t$. A batched learner may know the whole grid. For an arm $c$ and a round $t$, the number of pulls up to $t$
\emph{included} is
\[
\Ng_c(t):=\sum_{u\le t}\1(A_u=c),\qquad\text{and}\qquad\Reg(t):=\sum_c\Dg_c\,\E[\Ng_c(t)]
\]
is the (pseudo-)regret. A batched learner is \emph{consistent} on $\cP$ (for the given grid) if $\E_\nu[\Ng_c(\tl_\ell)]=o(\tl_\ell^\alpha)$ as $\ell\to\infty$, for every
$\nu\in\cP$, every sub-optimal arm $c$ of $\nu$ and every $\alpha>0$: consistency is only required at the ends of the batches, which is the criterion of the paper.
\end{definition}

\begin{definition}[the information model $\Mz$]\label{def:M0}\dt{s:M0}
\uses{def:grid}
A learner of $\Mz$ is a batched learner whose policy at a round $t\in(\tl_\ell,\tl_{\ell+1}]$ depends on the grid only
through $\tl_1,\dots,\tl_\ell$ and $t$: it does not know $\tl_{\ell+1}$, and it does not know any horizon. The lower bound of Section~\ref{sec:results} holds for every
batched learner, and the upper bound is attained by a deterministic learner of $\Mz$.
\end{definition}

\begin{definition}[reward stacks]\label{def:stacks}\dt{s:stack}
\uses{def:learner}
As usual \cite[Section~4.6]{LS} (to be checked), we realize the rewards with \emph{stacks}: $\Xs_{c,k}$ is the reward of the
$k$-th pull of arm $c$, the variables $\Xs_{c,k}$ are independent with $\Xs_{c,k}\sim\nu_c$, and $\hnu_{c,n}$, $\hm_{c,n}$ are the empirical distribution and the mean of
$\Xs_{c,1},\dots,\Xs_{c,n}$. At a round $t$ the learner has observed $\hnu_c(t):=\hnu_{c,\Ng_c(t)}$ and $\hm_c(t):=\hm_{c,\Ng_c(t)}$ only when $t$ is a batch end; at the end of batch
$\ell$ these are the statistics of all the rewards of $c$ observed so far (Lemma~\ref{lem:stacks}).
\end{definition}


\subsection{The algorithm HY}\label{sec:algo}

\HY{} works batch by batch. At each end $\tl_\ell$ it computes, from $\tl_1,\dots,\tl_\ell$ and the rewards observed so far, an IMED leader and IMED divergences, and a
test of every non-leader against the best well-sampled arm. In batch $\ell+1$ it first plays IMED rounds, with the frozen statistics and the current counts, up to the round
$\tl_\ell^{1+\thv_\ell}$ (the \emph{IMED phase}), and, if the batch is still running then, it follows the paced plan given by the test (the \emph{paced phase}).

\begin{definition}[the algorithm \HY$(\wz,\beta,\om,\thv,\gam)$]\label{def:hy}\dt{s:alg}
\uses{def:grid, def:learner, def:stacks, def:kinf}
\emph{Parameters.} A pseudo-count $\wz\ge0$, a net exponent $\beta\ge2$, non-increasing margins $\om_\ell\in(0,1]$, activation exponents $\thv_\ell\in(0,1]$ and
consistency margins $\gam_\ell\in(0,\frac12]$, all deterministic functions of $\tl_1,\dots,\tl_\ell$.

\smallskip\noindent\emph{Default parameters} (used in Theorem~\ref{thm:upper} and Theorem~\ref{thm:main}(b)):
\begin{center}
\fbox{\begin{tabular}{lll}
pseudo-count & $\wz:=1$ & (any $\wz\ge0$ is covered by the theorems)\\[2pt]
net exponent & $\beta:=2$ & (any $\beta\ge2$)\\[2pt]
margins & $\om_\ell:=\ody_\ell:=2^{-k_\ell}$, & $k_\ell:=\big\lfloor\tfrac14\log_2\max(\Lg_\ell,1)\big\rfloor$\\[2pt]
activation exponents & $\thv_\ell:=\begin{cases}1/\log\log\tl_\ell&\text{if }\tl_\ell\ge e^e,\\ 1&\text{if }\tl_\ell<e^e\end{cases}$ & \\[8pt]
consistency margins & $\gam_\ell:=\min\big(\tfrac12,\max(\Lg_\ell,1)^{-1/4}\big)$ &
\end{tabular}}
\end{center}

\emph{Batch $1$} pulls the arms in turn. At the end $\tl_\ell$ of every batch $\ell\ge1$, the learner computes, from $\tl_1,\dots,\tl_\ell$ only,
\[
\yg_\ell:=\xg_\ell^{\,2},\qquad\Ug_{\ell+1}:=\big\lceil\Lg_\ell^\beta\big\rceil,\qquad\tauA_\ell:=\tl_\ell^{\,1+\thv_\ell},\qquad\sA_\ell:=\lfloor\tauA_\ell\rfloor
\]
($\Ug_1:=1$, $\yg_0:=1$), and, writing $N_c:=\Ng_c(\tl_\ell)$ for every arm $c$, the following statistics (all computed at every end $\tl_\ell$):
\begin{enumerate}
\item[(1)] \emph{Pseudo-observation.} $\tnu_c(\tl_\ell):=\dfrac{N_c\,\hnu_c(\tl_\ell)+\wz\,\delta_0}{N_c+\wz}$, and its mean $\tm_c(\tl_\ell):=\dfrac{N_c\,\hm_c(\tl_\ell)}{N_c+\wz}$.
\item[(2)] \dt{t:leader}\emph{Leader and IMED divergences.} The \emph{leader} $\Ld_\ell$ is an arm of largest $\tm_c(\tl_\ell)$ (fixed tie rule), $\tms_\ell:=\tm_{\Ld_\ell}(\tl_\ell)$, and $\Kt_\ell(c):=\Ki(\tnu_c(\tl_\ell),\tms_\ell)$.
\item[(3)] \dt{t:ref}\emph{Reference and test divergences.} $\Wg_\ell:=\{c:N_c\ge\frac12\Ug_\ell\}$ (the \emph{well-sampled} arms); the \emph{reference} $\rr_\ell$ is an arm of largest $\tm_c(\tl_\ell)$ in $\Wg_\ell$
($\rr_\ell:=\Ld_\ell$ if $\Wg_\ell=\emptyset$); the \emph{test level} is $\mh_\ell:=\tm_{\rr_\ell}(\tl_\ell)$, and $\Kh_\ell(c):=\Ki(\tnu_c(\tl_\ell),\mh_\ell)$ for every arm $c$.
These divergences are only used by the paced phase, but they are computed at every end: the threshold of the next test uses $\Kh_\ell$.
\end{enumerate}
\emph{Batch $\ell+1$.} At a round $t\in(\tl_\ell,\tl_{\ell+1}]$ (the learner knows $t$, not $\tl_{\ell+1}$):

\smallskip\noindent\textbf{IMED phase, $t\le\sA_\ell$.} Pull an arm minimizing $\Jg_\ell(c,t):=\Ng_c(t-1)\Kt_\ell(c)+\yg_\ell\log\Ng_c(t-1)$ (with $\infty\cdot n:=\infty$ for $n\ge1$),
the leader $\Ld_\ell$ whenever it is a minimizer (other ties: fixed rule).

\smallskip\noindent\textbf{Paced phase, $t>\sA_\ell$.} The test below uses only the data up to $\tl_\ell$ (it can be made at $\tl_\ell$ or at the first paced round).
\begin{enumerate}
\item[(4)] \emph{Consistency threshold.} $\Lh_1:=\Lg_1$; for $\ell\ge2$, with $\Qp_\ell:=\{c\notin\Wg_\ell:\ c\ne\Ld_{\ell-1}\}$ (poorly sampled arms, except the previous leader),
\[
\Lh_\ell:=\min\Big(\Lg_\ell,\ \min_{c\in\Qp_\ell}\Big[\frac{N_c\,\Kh_{\ell-1}(c)}{\yg_{\ell-1}}+\log N_c\Big]\Big)\quad(\min\emptyset:=+\infty),\qquad\hg_\ell:=(1-\gam_\ell)\,\yg_{\ell-1}\,\Lh_\ell .
\]
\item[(5)] \dt{t:verdict}\emph{Verdicts.} An arm $c\ne\Ld_\ell$ is \emph{rejected} if $\Kh_\ell(c)>0$ and $\Zg_\ell(c):=N_c\Kh_\ell(c)\ge\hg_\ell$; it is \emph{ambiguous} otherwise.
\item[(6)] \dt{t:net}\emph{Paced targets.} For $c\ne\Ld_\ell$ and every integer $t\ge\tl_\ell$, $\nt_{\ell,c}(t):=K-1+\Ug_{\ell+1}$ if $c$ is ambiguous (the \emph{net}), and
$\nt_{\ell,c}(t):=K-1+\min\big(\Ug_{\ell+1},\lceil(1+\om_\ell)\yg_\ell\log t/\Kh_\ell(c)\rceil\big)$ if $c$ is rejected ($h/\infty:=0$).
\item[(7)] \emph{Rule.} Pull the first arm $c\ne\Ld_\ell$, in a fixed order, with $\Ng_c(t-1)<\nt_{\ell,c}(t)$; if there is none, pull $\Ld_\ell$.
\end{enumerate}
\end{definition}

\begin{definition}[vocabulary and first facts]\label{def:hy-vocab}\dt{s:act}
\uses{def:hy, def:grid}
The following words and facts about \HY{} are used throughout.
\begin{itemize}
\item \emph{Phases.} The \emph{IMED phase} of batch $\ell+1$ is $(\tl_\ell,\min(\sA_\ell,\tl_{\ell+1})]$, its \emph{paced phase} is $(\sA_\ell,\tl_{\ell+1}]$.
\item \emph{Activation.} Batch $\ell+1$ is \emph{activated} if it has paced rounds, i.e.\ $\tl_{\ell+1}>\sA_\ell$. Let $\Act:=\{\ell\ge1:\tl_{\ell+1}>\sA_\ell\}$. Activation is a
property of the grid alone.
\item \emph{Every arm is observed.} Batch $1$ has $\Bg_1\ge K$ rounds, so $N_c\ge1$ for every arm at every end $\tl_\ell$. Hence $\log N_c\ge0$, $\ 0\le\Lh_\ell\le\Lg_\ell\ $ and
$\hg_\ell\ge0$.
\item \emph{First test.} At $\ell=1$, $\Ug_1=1$, so $\Wg_1=\cA$ and $\Qp_1=\emptyset$.
\end{itemize}
\end{definition}

\begin{center}
\fbox{\begin{minipage}{0.96\textwidth}\small
\noindent\textbf{Pseudo-code of \HY$(w_0,\beta,\omega,\vartheta,\gamma)$}; at a round $t$ of batch $\ell+1$ only $t_1,\dots,t_\ell$ and $t$ are known.
\begin{tabbing}
\hspace*{1.4em}\=\hspace*{1.4em}\=\hspace*{1.4em}\=\hspace*{1.4em}\=\hspace*{1.4em}\=\kill
\textbf{batch $1$:} pull the arms in turn; $U_1\gets1$, $y_0\gets1$\\
\textbf{for} $\ell=1,2,\dots$ \hfill(at time $t_\ell$)\\
\>$y_\ell\gets x_\ell^2$, \ $U_{\ell+1}\gets\lceil\mathsf L_\ell^\beta\rceil$, \ $s_\ell\gets\lfloor t_\ell^{1+\vartheta_\ell}\rfloor$\\
\>$\widetilde\nu_c,\widetilde{\mathbf m}_c\gets$ statistics with the pseudo-count $w_0$; \ $\widehat c_\ell\gets\arg\max_c\widetilde{\mathbf m}_c$\\
\>$\widetilde K_\ell(c)\gets\mathcal K_{\inf}(\widetilde\nu_c,\widetilde{\mathbf m}_{\widehat c_\ell})$ for every arm $c$\\
\>$W_\ell\gets\{c:N_c(t_\ell)\ge\frac12U_\ell\}$; \ $r_\ell\gets\arg\max_{c\in W_\ell}\widetilde{\mathbf m}_c$ \ ($r_\ell\gets\widehat c_\ell$ if $W_\ell=\emptyset$)\\
\>$\widehat\mu_\ell\gets\widetilde{\mathbf m}_{r_\ell}$; \ $\widehat K_\ell(c)\gets\mathcal K_{\inf}(\widetilde\nu_c,\widehat\mu_\ell)$ for every arm $c$ \hfill(used by the next test)\\
\>\textbf{for} each round $t=t_\ell+1,t_\ell+2,\dots$ until the batch ends \hfill($t_{\ell+1}$ is not known)\\
\>\>\textbf{if} $t\le s_\ell$ \hfill(\textbf{IMED phase})\\
\>\>\>pull $\arg\min_c\big[N_c(t-1)\widetilde K_\ell(c)+y_\ell\log N_c(t-1)\big]$, and $\widehat c_\ell$ if it is a minimizer\\
\>\>\textbf{else} \hfill(\textbf{paced phase})\\
\>\>\>\textbf{if} $t=s_\ell+1$ \hfill(test, with the data up to $t_\ell$)\\
\>\>\>\>$\widehat\Lambda_1\gets\mathsf L_1$ if $\ell=1$; otherwise\\
\>\>\>\>$\widehat\Lambda_\ell\gets\min\big(\mathsf L_\ell,\min_{c\notin W_\ell,\,c\ne\widehat c_{\ell-1}}[N_c(t_\ell)\widehat K_{\ell-1}(c)/y_{\ell-1}+\log N_c(t_\ell)]\big)$\\
\>\>\>\>$h_\ell\gets(1-\gamma_\ell)\,y_{\ell-1}\,\widehat\Lambda_\ell$\\
\>\>\>\>\textbf{for} each $c\ne\widehat c_\ell$: \textbf{if} $\widehat K_\ell(c)>0$ and $N_c(t_\ell)\widehat K_\ell(c)\ge h_\ell$ \hfill(rejected)\\
\>\>\>\>\>$n_{\ell,c}(\cdot):=K-1+\min\big(U_{\ell+1},\lceil(1+\omega_\ell)y_\ell\log(\cdot)/\widehat K_\ell(c)\rceil\big)$\\
\>\>\>\>\textbf{else} $n_{\ell,c}(\cdot):=K-1+U_{\ell+1}$ \hfill(ambiguous: net)\\
\>\>\>\textbf{if} some $c\ne\widehat c_\ell$ has $N_c(t-1)<n_{\ell,c}(t)$: pull the first such $c$; \textbf{else} pull $\widehat c_\ell$
\end{tabbing}
\end{minipage}}
\end{center}

\paragraph{Comments: the three ideas.}
\begin{enumerate}
\item[(1)] \emph{Pacing.} A rejected arm is not pulled at the start of the batch up to a fixed target: its target $(1+\om_\ell)\yg_\ell\log t/\Kh_\ell(c)$ grows with the
current round $t$. If the batch ends at $\tl_{\ell+1}$, the plan evaluated there is $(1+\om_\ell)\xg_\ell^2\Lg_{\ell+1}/\Kh_\ell(c)$, and the next test is made at the level
$\approx\xg_\ell^2\Lg_{\ell+1}\approx\Lg_{\ell+3}$: whatever the batch turns out to be, the plan aims at the test that will actually be made, two batches ahead. The coefficient
$\yg_\ell=\xg_\ell^2$ is fixed at $\tl_\ell$, from the past ends only. A queue lemma (Lemma~\ref{lem:service}) shows that the demands are served at every round, up to
$K-1$ pulls, after an initial burst of size $O(K\Ug_{\ell+1})$.
\item[(2)] \emph{IMED phase.} The IMED phase uses the index of IMED \cite{HT15} with the statistics frozen at $\tl_\ell$: the logarithm in the index uses the current count, the weight
$\yg_\ell$ multiplies it, and a pseudo-observation $0$ of weight $\wz$ lowers the means of the poorly sampled arms. If the batch is long, the IMED phase lasts about
$\tl_\ell^{1+\thv_\ell}$ rounds; the arm that receives most of them certifies that every IMED index is at least $\log\tl_\ell$ at the end of the batch
(Lemma~\ref{lem:W}, \emph{the IMED phase lifts the water level}). This ``water level'' is what makes the leader of a long batch reliable when $x=1$, even after many short batches.
\item[(3)] \emph{Consistency threshold.} The test does not use the level $\yg_{\ell-1}\Lg_\ell$ directly but $(1-\gam_\ell)\yg_{\ell-1}\Lh_\ell$, where $\Lh_\ell$ is a water level
computed with the divergences $\Kh_{\ell-1}$ of the \emph{previous} test. An arm $c\in\Qp_\ell$ is then declared ambiguous only if
$\Ng_c(\tl_\ell)\Kh_\ell(c)<(1-\gam_\ell)\big[\Ng_c(\tl_\ell)\Kh_{\ell-1}(c)+\yg_{\ell-1}\log\Ng_c(\tl_\ell)\big]$: its divergence has dropped by a factor about $1-\gam_\ell$ between two
tests. A sub-optimal arm with $\gtrsim\log t$ observations almost never does this, however late IMED is with respect to $\log t$ (Lemma~\ref{lem:amb}, \emph{consistency test}). When $x>1$,
$\Lh_\ell=\Lg_\ell$ from some rank on, on every run (Lemma~\ref{lem:thrA}).
\end{enumerate}
\paragraph{Properties.} (i) \emph{Information.} Every quantity used in batch $\ell+1$ is a function of $\tl_1,\dots,\tl_\ell$, of the observations up to $\tl_\ell$, of the current counts and of
$t$: \HY{} is a deterministic learner of $\Mz$. It never uses $\Bg_{\ell+1}$, a horizon or the instance, so it can be run when the end of the current batch is only revealed
when it happens. (ii) \emph{Tie rules.} The leader, the reference $\rr_\ell$, the arm pulled among several minimizers of $\Jg_\ell$ and the order of the paced rule are given by
fixed rules; the proofs do not use these rules, and the Lean theorems hold for any rules satisfying their specification, with any order of batch~$1$ in which every arm
is pulled (Section~\ref{sec:lean}). (iii) \emph{Cost.} Each test costs $K$ evaluations of $\Ki$ (one-dimensional concave maximizations), and each round $O(K)$ operations.
(iv) A batch that is not activated consists of IMED rounds only; the leader is always the IMED leader, with no eligibility condition.


%======================================================================
\section{Closest known}\label{sec:closest}

The KNOWN.md of the roadmap BatchedBandits has no reviewed entry yet. The roadmap states the result of this paper in its
Appendix~A, from the author's manuscript, and its known-result item F1 (the one-batch lower bound, $x=1$) is
Lemma~\ref{lem:onebatch} below. The related literature, as described in the author's manuscript and completed in the
revision of 2026-11-14 (static grids, batches chosen for computation, multi-stage tests, doubling tricks):

\paragraph{Related work.} \emph{Batched bandits with a chosen grid}: Perchet et al.~\cite{PRCS16} for two arms, Gao et al.~\cite{GHRZ19} and Esfandiari et al.~\cite{EKMM21} (to be checked) for
$K$ arms, Jin et al.~\cite{JTXHXG21} and \cite{JXXG21} (to be checked) for asymptotic optimality with few batches, Karbasi et al.~\cite{KMS21} (to be checked) for Thompson sampling. These algorithms design
their grid and do not apply to imposed batches. \emph{Batches chosen for computation}: the KLinf-UCB of Agrawal, Juneja and Koolen \cite[Section~3, Algorithm~1]{AJK21}
pulls the arm of largest index for a batch of $\max\{1,\tilde\varepsilon N\}$ rounds, $N$ being its current count, so its batches are chosen by the algorithm (and random); on
heavy-tailed classes its constant is optimal up to the factor $1+\tilde\varepsilon$ \cite[Theorem~1 and Corollary~2]{AJK21}. \emph{Imposed batches}: Kalkanli and {\"O}zg{\"u}r
\cite[Theorem~1]{KO21b} show that Gaussian Thompson sampling has a regret equivalent to
$\sum_a(2/\Delta_a)\log T$ almost surely when $\limsup_\ell\log B_{\ell+1}/\log t_\ell<1$ (their condition~(6)), which implies $x_\ell\to1$; their constant is optimal for Gaussian rewards of
unit variance, not for bounded rewards. Their batch ends may be chosen by an adversary at the previous end, as functions of the history \cite[Section~II.A]{KO21b}: random
ends of this kind are the subject of Open question~\ref{oq:random}. A companion work studies imposed batch sizes with a batched version of IMED and with an algorithm that plans its exploration from a
\emph{known} grid \cite{P26} (to be checked); here the grid is unknown, and we reprove everything that we use from that work. \emph{Delayed feedback}: batching is a form of delay; the reduction of
Joulani et al.~\cite[Theorem~6]{JGS13} adds to the regret without delays a term proportional to the largest number of missing feedbacks, here the largest batch, which preserves optimality only for batches of size $o(\log t)$. \emph{Non-parametric
optimality}: for distributions on $[0,1]$ the relevant divergence is $\Ki$ \cite{BK96} (to be checked), \cite[(2) and Theorem~8]{HT10}; it is attained sequentially by IMED
\cite{HT15}, KL-UCB \cite{CGMMS13,GHMS22} and non-parametric Thompson sampling \cite{RH20}. \emph{Multi-stage tests}: the mechanism
of our lower bound is that of Lorden \cite[Theorem~3 and its proof]{L83}: under a density $g$ intermediate between the two hypotheses, the sample size of an asymptotically
optimal $k$-stage test after $k-1$ stages is forced to its full asymptotic value, so that the last stage is negligible; evaluating the expected sample size at a parameter
between the two hypotheses goes back to Hoeffding \cite{H60}. Here the stages are
imposed, and the price of not being able to react inside a stage is the factor $x^2$. \emph{Doubling tricks}: Besson and Kaufmann \cite[Section~6]{BK18} conjecture
that the exponential doubling trick, which restarts a base algorithm at times $T_i\approx T_0^{b^i}$, never brings the right constant; they show that it multiplies a bound
$c\log T$ by at most $b^2/(b-1)$ \cite[Remark~8]{BK18}.

\paragraph{Why Theorem~\ref{thm:main} does not follow from these works.}
\begin{itemize}
\item \emph{The classical change of measure} \cite{LR85} (to be checked), \cite[(6)]{GMS19}, applied at the ends of the
batches, gives Lemma~\ref{lem:onebatch}: the count at $\tl_\ell$ must already pay for the horizon $\tl_{\ell+1}$, hence the
constant $x/\Ki(\nu_a,\m^\star)$ when $x>1$, not $x^2/\Ki(\nu_a,\m^\star)$. The attempt to reach $x^2$ this way fails because
this change of measure tilts every reward of $a$; it does not use that the allocation of batch $\ell+1$ is decided without
the rewards of batch $\ell$. Lemma~\ref{lem:twobatches} tilts only those rewards, towards a distribution of mean exactly
$\m^\star$.
\item \emph{Sequential index algorithms} (IMED \cite{HT15}, KL-UCB \cite{CGMMS13,GHMS22}) attain
$1/\Ki$ when the feedback is immediate. On imposed batches their statistics are frozen at $\tl_\ell$; when $x>1$ a batch
lasts about $\tl_\ell^{\,x}$ rounds, so a wrong decision taken at $\tl_\ell$ (for instance an optimal arm judged
sub-optimal) costs a power of $\tl_\ell$. The test of \HY{} makes such mistakes summable after multiplication by the
batch size, thanks to a strict gain in the deviation inequalities (Lemmas~\ref{lem:single} and~\ref{lem:pair}); its paced
plan aims at the test that will be made two batches later; its IMED phase and its consistency threshold keep the leader
and the test reliable when $x=1$ (Subsection~\ref{sec:algo}).
\item \emph{Batched bandits with a chosen grid} \cite{PRCS16,GHRZ19}, \cite{EKMM21} (to be checked)
design the batch ends. Their algorithms do not apply when the ends are imposed and revealed only when they happen, and
their lower bounds (\cite[Theorem~4]{PRCS16}, \cite[Section~3]{GHRZ19}) are minimax bounds in terms of the number of
batches; they do not give the instance-dependent constant at imposed batch ends.
\item \emph{Static grids and anytime learners} \cite[Theorem~4]{JTXHXG21}: for one-parameter exponential families (a unique
best arm, bounded variance), ``for any static grid and constant $c>0$, no anytime algorithm can achieve the asymptotically
optimal regret [\dots] within $c\log\log T$ batches'', asymptotic optimality meaning there that $\lim_T R_T/\log T$ equals the
constant $\sum_i\Delta_i/\kl(\mu_i,\mu_1)$ along all rounds $T$. A grid with $\xg_\ell\to x>1$ has about
$\log\log T/\log x$ batch ends below $T$, so Theorem~4 implies, for these families, that the sequential constant is not
attained on such a grid: this is a qualitative special case of Theorem~\ref{thm:lower}(ii). It is presented by its authors
as a confirmation of a conjecture of Besson and Kaufmann \cite[Section~6]{BK18} (stated there for the exponential doubling
trick). Its proof \cite[Appendix~C]{JTXHXG21} is the one-batch change of measure: inside a gap between two batch ends
$t_j<t_{j+1}$ it takes two times $N_n<N_{n+1}=N_n^{x'}$ with $x'=e^{1/(4c)}$, changes the mean of one arm only, on the
observations up to $t_j$, and applies the Bretagnolle--Huber inequality; this gives a count at $t_j$ at least
$x'(1-\epsilon')/(\kl(\mu_i,\mu_1)+\epsilon')$ times $\log t_j$ along a subsequence of ends, and no quantitative factor is
stated. With $N_{n+1}$ just below $t_{j+1}$, the same argument gives at most the factor $\log t_{j+1}/\log t_j\to x$, that
is the bound of Lemma~\ref{lem:onebatch}, not the factor $x^2$. Theorem~\ref{thm:lower}(ii) does not follow, for three
reasons: the factor $x^2$ needs the change of measure of Lemma~\ref{lem:twobatches}, which tilts only the rewards of $a$ in
one batch towards a law of mean exactly $\m^\star$, whereas moving the mean of every reward of $a$ cannot give more than
$x$; Theorem~4 is about exponential families and $\kl$, not about all distributions on $[0,1]$ and $\Ki$ (with any number of
optimal arms); and it does not give the constant $x^2/\Ki$, which is the content of Theorem~\ref{thm:lower}(ii). It is a
lower bound only, so Theorem~\ref{thm:upper} does not follow from it either.
\item \emph{Imposed batches with Thompson sampling} \cite[Theorem~1]{KO21b}: only grids with $x=1$, and the constant
$2/\Delta_a$ of Gaussian rewards of unit variance, which is not the optimal constant for rewards in $[0,1]$.
\item \emph{Delayed feedback} \cite[Theorem~6]{JGS13}: the reduction preserves optimality only for batches of size
$o(\log t)$, so not when $x>1$.
\item \emph{The companion manuscript} \cite{P26} (to be checked; unpublished; imposed batch sizes, a batched IMED and a plan
computed from a known grid): every result used from it is proved again here.
\end{itemize}

%======================================================================
\section{Proofs}\label{sec:proofs}\label{sec:results}

This section states the two theorems behind Theorem~\ref{thm:main}, explains how their proofs work, and proves Theorem~\ref{thm:main}; Sections~\ref{sec:app-tools}--\ref{sec:app-end} prove the theorems.

\subsection{The two theorems}


\begin{theorem}[lower bound: explore two batches ahead]\label{thm:lower}
\uses{def:ass-x, def:model, def:learner, def:grid, def:kinf}
Let the grid satisfy Assumption~\ref{def:ass-x} with $x\in[1,\infty)$, let a batched learner be consistent on $\cP$, let $\nu\in\cP$ and let $a$ be a sub-optimal arm.
\begin{enumerate}
\item[(i)] If $x=1$: $\displaystyle\liminf_{\ell\to\infty}\frac{\E_\nu[\Ng_a(\tl_\ell)]}{\log\tl_\ell}\ \ge\ \frac1{\Ki(\nu_a,\m^\star)}$. (This holds on every grid, convergent or not,
since $\log\tl_{\ell+1}\ge\log\tl_\ell$ in Lemma~\ref{lem:onebatch}.)
\item[(ii)] If $x>1$: $\displaystyle\limsup_{\ell\to\infty}\frac{\E_\nu[\Ng_a(\tl_\ell)]}{\log\tl_\ell}\ \ge\ \frac{x^2}{\Ki(\nu_a,\m^\star)}$; and if moreover
$\E_\nu[\Ng_a(\tl_\ell)]=O(\log\tl_\ell)$, the same holds with a $\liminf$.
\item[(iii)] (Regret.) If $x=1$, or if $x>1$ and $\E_\nu[\Ng_c(\tl_\ell)]=O(\log\tl_\ell)$ for every sub-optimal arm $c$,
\[
\liminf_{\ell\to\infty}\frac{\Reg(\tl_\ell)}{\log\tl_\ell}\ \ge\ x^2\sum_{c:\Dg_c>0}\frac{\Dg_c}{\Ki(\nu_c,\m^\star)} .
\]
\end{enumerate}
\end{theorem}
\noindent(see proof in Section~\ref{sec:app-lower}.)

\noindent For $x>1$, the bound $x^2/\Ki$ is $x$ times larger than the bound $x/\Ki$ that the classical change of measure gives at the ends of the batches (the
count at $\tl_\ell$ must cover the horizon $\tl_{\ell+1}$, Lemma~\ref{lem:onebatch}). The factor $x^2$ comes from a sharper change of measure, in which only the rewards of
$a$ in batch $\ell$ are tilted towards a distribution of mean exactly $\m^\star$; the learner cannot tell whether $a$ is barely sub-optimal or barely optimal, and has
to allocate batch $\ell+1$, of size $\approx\tl_{\ell+1}$, without being able to react inside it.


\begin{theorem}[upper bound: HY attains $x^2/\Ki$ without knowing the grid]\label{thm:upper}
\uses{def:ass-x, def:model, def:hy, def:learner, def:grid, def:kinf}
Let the grid satisfy Assumption~\ref{def:ass-x} with $x\in[1,\infty)$, let $\nu\in\cP$ and let $a$ be a sub-optimal arm. Then \HY{} with the default parameters,
\HY$(\wz,\beta,\ody,\thv,\gam)$, is consistent on $\cP$ and
\[
\lim_{\ell\to\infty}\frac{\E[\Ng_a(\tl_\ell)]}{\log\tl_\ell}=\frac{x^2}{\Ki(\nu_a,\m^\star)},\qquad
\lim_{\ell\to\infty}\frac{\Reg(\tl_\ell)}{\log\tl_\ell}=x^2\sum_{c:\Dg_c>0}\frac{\Dg_c}{\Ki(\nu_c,\m^\star)} .
\]
The same holds for every $\wz\ge0$ and every $\beta\ge2$, the other parameters being the default ones.
\end{theorem}
\noindent(see proof in Subsection~\ref{sec:app-proofupper}; it combines a non-asymptotic bound for $x>1$, Proposition~\ref{prop:dyA}, and one for $x=1$,
Proposition~\ref{prop:fixedC}.)


\subsection{Proof of the main result}

\begin{proof}[Proof of Theorem~\ref{thm:main}]
\proves{thm:main}
\uses{thm:lower, thm:upper, def:learner, def:grid, def:model, def:hy}
\stp{(a), (b)}{counts}(a) is Theorem~\ref{thm:lower}, and (b) is Theorem~\ref{thm:upper}.

\stp2{regret of a consistent learner}Two cases.
\begin{itemize}
\item If some sub-optimal arm $c$ has $\limsup_\ell\E[\Ng_c(\tl_\ell)]/\log\tl_\ell=\infty$, then
$\limsup_\ell\Reg(\tl_\ell)/\log\tl_\ell\ge\Dg_c\limsup_\ell\E[\Ng_c(\tl_\ell)]/\log\tl_\ell=\infty$.
\item Otherwise every sub-optimal arm has $\E[\Ng_c(\tl_\ell)]=O(\log\tl_\ell)$, and Theorem~\ref{thm:lower}(iii) gives $\liminf_\ell\Reg(\tl_\ell)/\log\tl_\ell\ge C_x$.
\end{itemize}
\stp3{regret of \HY}The limit is Theorem~\ref{thm:upper}.
\end{proof}


\subsection{How the proofs work}\label{sec:how}

\paragraph{The lower bound.} For $x=1$ it is the classical change of measure \cite{LR85} (to be checked), \cite[(6)]{GMS19}, applied to the observations available at the
end of batch $\ell+1$: they contain the actions of batch $\ell+1$ but only the rewards of the first $\ell$ batches, so that the count $\Ng_a(\tl_\ell)$ must already pay for
the horizon $\tl_{\ell+1}$ (Lemma~\ref{lem:onebatch}). For $x>1$, consider the histories in which the rewards of $a$ in batch $\ell$ look like a distribution $G$ of mean
exactly $\m^\star$ with $\KL(\nu_a,G)\approx\da$ (a \emph{boundary tilt}, Lemma~\ref{lem:tilts}). The instances in which $\nu_a$ is replaced by $(1-s)G+s\delta_0$ ($a$ barely
sub-optimal) or by $(1-s)G+s\delta_1$ ($a$ barely optimal) explain these histories equally well, at a cost $O(s)$ per reward of batch $\ell$ and $\approx\da$ per earlier
reward. Giving the bulk of batch $\ell+1$ to $a$ is a polynomial loss under the first, not giving it is a polynomial loss under the second, so consistency forces
$\Ng_a(\tl_{\ell-1})\da\gtrsim\log\tl_{\ell+1}$, that is $\Ng_a(\tl_\ell)\gtrsim\log\tl_{\ell+2}/\da\approx x^2\log\tl_\ell/\da$. Only the rewards of batch $\ell$ are tilted, so the
number of pulls of $a$ in that batch, decided at $\tl_{\ell-1}$, has the same law under the tilted measure; the tightness of this number is where logarithmic counts are used.

\paragraph{The upper bound when $x>1$: a paced plan with an IMED prelude.} Every batch is eventually activated, the IMED phase covers the logarithmic window
$[\Lg_\ell,(1+\thv_\ell)\Lg_\ell]$, which is negligible in front of $[\Lg_\ell,\xg_{\ell+1}\Lg_\ell]$, and the paced phase is long. The proof is a decomposition of $\Ng_a(t)$ over
the batches. (i) \emph{Structure} (Lemma~\ref{lem:structA}): by the queue lemma started at $\sA_\ell$, every demand is served,
the leader receives at least $\Ug_{\ell+1}$ pulls, and an arm that is not well sampled at $\tl_{\ell+1}$ was rejected at $\tl_\ell$ and planned up to
$(1+\om)\yg_\ell\Lg_{\ell+1}/\Kh_\ell$; consequently $\Lh_\ell=\Lg_\ell$ and the threshold is deterministic (Lemma~\ref{lem:thrA}). (ii) \emph{The leader}
(Lemma~\ref{lem:leaderA}): if the leader is sub-optimal, either no optimal arm is well sampled, and then an optimal arm was rejected at the previous test against a level
below $\mtwo+\varepsilon$ (event $\Vg_\ell$), or a poorly sampled sub-optimal arm has a mean above $\m^\star-3\varepsilon$ although it was rejected at the previous test
against a level above $\m^\star$ (event $\Pig_\ell$). Both have probability $\lesssim\tl_{\ell+1}^{-(1+\eg)}$ thanks to a strict gain $\eg>0$ in the deviation inequalities of
$n\Ki$ (Lemmas~\ref{lem:single} and~\ref{lem:pair}), which makes them summable after multiplication by the batch size. (iii) \emph{Good batches} (Lemma~\ref{lem:goodA}):
the plan of batch $\ell$, evaluated at the actual end $\tl_\ell$, aimed above the present threshold, so $a$ is rejected and planned at $(1+\om)\yg_\ell\log t/\Ki(\nu_a,\m^\star)$
up to $\varepsilon$; the IMED phase pulls $a$ fewer times. Every other mistake costs at most the net $2\Ug_{\ell+1}$, which is polylogarithmic; a layer-cake argument sums
these nets (Lemma~\ref{lem:sums}).

\paragraph{The upper bound when $x=1$: the water level goes through the paced phases.} A batch made of IMED rounds satisfies a deterministic property (W2): every arm
pulled in a batch had, at the start of the batch, a logarithmic count below the final water level $\Lam$ of the batch. With the growth of the grid ($\xg_\ell\to1$) this
gives many ``high'' batches, whose leaders have $\tl^{1-\dl}$ observations and are therefore accurate; a chain argument then shows that an arm with few observations
was explored at the level $\approx\log\tl_j$ in each high batch, which makes the bad events of the decomposition rare. A paced phase does not satisfy (W2) in general.
\HY{} does, from some rank on, because an activated batch first runs $\approx\tl_\ell^{1+\thv_\ell}$ IMED rounds: the arm that gets most of them certifies a water level above
$\log\tl_\ell$, above the logarithm of every count at the start of the batch (Lemma~\ref{lem:W}, \emph{the IMED phase lifts the water level}). So the whole water-level machinery holds for \HY{}
(Section~\ref{sec:app-water}), and gives an accurate leader \emph{and} an optimal arm with at least $\frac12\Ug$ observations (Lemma~\ref{lem:optW}), hence an accurate test
level. In a paced phase a rejected sub-optimal arm then stays below $K+(1+\om_\ell)\yg_\ell\log t/\Ki(\hnu_a,\m^\star-2\varepsilon)$, the bound of IMED up to $(1+\om_\ell)\yg_\ell\to1$.

\paragraph{The consistency threshold.} The remaining cost is a net $\approx\Lg_\ell^\beta$ when $a$ is ambiguous at the start of an activated batch. With a threshold
$\approx\yg\Lg_\ell$, this would need the evidence of $a$ to be $\approx\yg\Lg_\ell$, i.e.\ IMED to be up to date with $\log t$; after short batches it may lag by a fixed
fraction of $\log t$. The consistency threshold avoids the question: if $a\in\Qp_\ell$, then $\Ng_a\Kh_{\ell-1}(a)\ge\yg_{\ell-1}(\Lh_\ell-\log\Ng_a)$ by definition of $\Lh_\ell$, and
$a$ is rejected as soon as $\Kh_\ell(a)\ge(1-\gam_\ell/2)\Kh_{\ell-1}(a)$. The two divergences are close to $\Ki(\nu_a,\m^\star)$ when the two test levels and $\Ki(\hnu_{a,n},\m^\star)$
are $\gam_\ell$-accurate, which fails with probability $e^{-c\sqrt{\Lg_\ell}/\log^2\Lg_\ell}$ for $\gam_\ell=\Lg_\ell^{-1/4}$ (Lemma~\ref{lem:amb}, \emph{consistency test}). Finally, activated batches are
sparse on the logarithmic scale (Lemma~\ref{lem:sparse}, \emph{sparse activations}), which makes the nets summable.


%======================================================================
\section{Tools}\label{sec:app-tools}

\par\smallskip
\Plan this section gathers the tools used by the proofs.
\begin{itemize}
\item Properties of $\Ki$ (Lemma~\ref{lem:kinf}, classical \cite{HT10,GHMS22}; we prove them from the dual formula, the only formula used in the
upper bound) and of the pseudo-observation (Lemma~\ref{lem:pseudo}).
\item Deviation inequalities: Hoeffding, deviations of $n\Ki$ \emph{with a gain} (Lemmas~\ref{lem:single}, \ref{lem:pair}), and explicit deviations of the empirical $\Ki$
(Lemmas~\ref{lem:lowerdev}, \ref{lem:upperdev}). They are variants of known results \cite{CGMMS13,GHMS22,HT15} with the constants that we need; we
prove them.
\item Reward stacks (Lemma~\ref{lem:stacks}), and elementary sums and sequences (Subsection~\ref{sec:app-sums}).
\end{itemize}
\par\smallskip

\subsection{Properties of \texorpdfstring{$\mathcal K_{\inf}$}{Kinf}}

\begin{lemma}[properties of $\Ki$]\label{lem:kinf}
\uses{def:kinf}
Let $G$ be a distribution on $[0,1]$ with mean $p$, and $\mu\in[0,1)$.
\begin{enumerate}
\item[(i)] $0\le\Ki(G,\mu)\le\kap(\mu)$, and $\kap$ is non-decreasing on $[0,1)$.
\item[(ii)] If $p\ge\mu$, then $\Ki(G,\mu)=0$.
\item[(iii)] $\mu\mapsto\Ki(G,\mu)$ is non-decreasing on $[0,1)$.
\item[(iv)] If $p\le q<\mu$, then $\Ki(G,\mu)\ge(\mu-q)^2/8$.
\item[(v)] If $0<\varepsilon\le\mu$, then $\Ki(G,\mu-\varepsilon)\ge\Ki(G,\mu)-\varepsilon/(1-\mu)$.
\item[(vi)] (scaling) If $\lambda\in\Lset{\mu}$, $\tau\in(0,1]$ and $y\in[0,1]$, then $\log\big(1-\tau\lambda(y-\mu)\big)\ge\log\tau+\log\big(1-\lambda(y-\mu)\big)$.
\end{enumerate}
\end{lemma}
\role{used throughout: monotonicity in the level (iii), quadratic lower bound (iv), shift of the level (v), scaling (vi).}
\noindent(see proof in Subsection~\ref{sec:pf-lem:kinf}.)


\begin{lemma}[pseudo-observation]\label{lem:pseudo}
\uses{def:hy, def:kinf, def:learner, def:grid, def:stacks}
Let $G$ be a distribution on $[0,1]$ of mean $p$, $N\ge1$ an integer, $\tilde G:=(NG+\wz\delta_0)/(N+\wz)$ and $\tilde p:=Np/(N+\wz)$ its mean. For every $\mu\in[0,1)$:
(i) $\Ki(\tilde G,\mu)\ge\Ki(G,\mu)$; (ii) $N\Ki(\tilde G,\mu)\le N\Ki(G,\mu)+\wz\kap(\mu)$; (iii) $p-\frac{\wz}{N+\wz}\le\tilde p\le p$.
In particular, for $c$ with $\Ng_c(\tl_\ell)\ge\frac12\Ug_\ell$, $\tm_c(\tl_\ell)\ge\hm_c(\tl_\ell)-2\wz/\Ug_\ell$.
\end{lemma}
\role{compares the statistics of \HY{} (with the pseudo-observation) with the empirical ones, in Sections~\ref{sec:app-xsup} and~\ref{sec:app-xone}.}
\noindent(see proof in Subsection~\ref{sec:pf-lem:pseudo}.)


\noindent We apply Lemma~\ref{lem:pseudo} with $G=\hnu_c(\tl_\ell)$ and $N=\Ng_c(\tl_\ell)\ge1$; then $\tilde G=\tnu_c(\tl_\ell)$ and $\tilde p=\tm_c(\tl_\ell)$.

\subsection{Deviations}

\begin{lemma}[Hoeffding]\label{lem:hoeffding}
Let $Z_1,Z_2,\dots$ be i.i.d.\ in $[0,1]$ with mean $\mu_Z$ and $S_n:=Z_1+\dots+Z_n$. (a) For $n\ge1$ and $y\ge0$, $\PP(S_n\ge n\mu_Z+y)$ and $\PP(S_n\le n\mu_Z-y)$ are at most
$e^{-2y^2/n}$. (b) For $\rho>0$ and $u>0$, $\PP(\exists n\ge u:|S_n/n-\mu_Z|>\rho)\le2e^{-2\rho^2u}/(1-e^{-2\rho^2})$ ($n$ ranges over the integers).
\end{lemma}
\role{accurate means: Lemmas~\ref{lem:probasA}(a), \ref{lem:controlC}(a), and the event $\Acc$ of Lemma~\ref{lem:amb}; (a) also in Lemma~\ref{lem:typ}
and in the proofs of Lemmas~\ref{lem:lowerdev}, \ref{lem:upperdev}.}
\noindent(see proof in Subsection~\ref{sec:pf-lem:hoeffding}.)


\begin{definition}[the empirical $\Ki$ and a grid of the dual variable]\label{def:empkinf}\dt{s:tilt}
\uses{def:kinf}
In the rest of this subsection $Y_1,Y_2,\dots$ are i.i.d.\ with law $\nu$ on $[0,1]$ and mean $m$, $S_n:=Y_1+\dots+Y_n$, and
$\hnu_n$ is the empirical distribution of $Y_1,\dots,Y_n$. For $\theta\ge0$ and $q\in\mathbb R$ let $\Phis_\nu(\theta,q):=\E[e^{\theta(Y_1-q)}]$ and
$\Psis_\nu(\theta,q):=\E[(1-Y_1)e^{\theta(Y_1-q)}]$. For reals $y_1,\dots,y_n\in[0,1]$ with empirical distribution $\widehat G$ and $r\in[0,1)$, let
$F(\lambda):=\sum_{i\le n}\log(1-\lambda(y_i-r))$, so that $n\Ki(\widehat G,r)=\sup_{\lambda\in\Lset{r}}F(\lambda)$.\dt{s:gridl} For $n\ge1$ let $\taun_n:=1-\frac1{2n}$,
$\lamk^{(k)}:=\taun_n^k/(1-r)$ and $\kn_n:=\lceil2n\log\frac n{1-r}\rceil+1$.
\end{definition}

\begin{lemma}[a grid for the dual formula]\label{lem:grid}
\uses{def:kinf, def:empkinf}
Let $n\ge1$ and $z\ge2$. If $n\Ki(\widehat G,r)\ge z$, then $F(\lamk^{(k)})\ge z-2$ for some $k\in\{1,\dots,\kn_n\}$. Moreover $\lamk^{(k)}\in\Lset{r}$ for $k\ge1$, and
$\kn_n\le(4+2\kap(r))n^2$.
\end{lemma}
\role{reduces the supremum of the dual formula to a finite grid, in the proofs of Lemmas~\ref{lem:single} and~\ref{lem:pair}.}
\noindent(see proof in Subsection~\ref{sec:pf-lem:grid}.)


\begin{lemma}[single, with a gain]\label{lem:single}
\uses{def:kinf, def:stacks, def:empkinf}
Let $0\le r<m$ and $\eta\ge0$ with $2\eta r\le m-r$. Then for every $n\ge1$ and $z\in\mathbb R$,
\[
\PP\big(n\Ki(\hnu_n,r)\ge z\big)\ \le\ \kn_ne^{2(1+2\eta)}e^{-(1+2\eta)z}\ \le\ \big(4+2\kap(r)\big)e^{2(1+2\eta)}\,n^2e^{-(1+2\eta)z}.
\]
\end{lemma}
\role{the probability that an optimal arm is wrongly rejected: Lemmas~\ref{lem:constA}(b), \ref{lem:controlC}(b).}
\noindent(see proof in Subsection~\ref{sec:pf-lem:single}.)


\begin{lemma}[pairs, with a gain]\label{lem:pair}
\uses{def:kinf, def:stacks, def:empkinf}
Let $0\le r<1$, $q\in\mathbb R$, $\theta\ge0$ and $\eta\ge0$ with $e^{2\eta\kap(r)}\Phis_\nu(\theta,q)\le1$ and $e^{2\eta\kap(r)}\Psis_\nu(\theta,q)\le1-r$. Then for all integers
$1\le n\le n'$ and every $z\in\mathbb R$,
\[
\PP\big(n\Ki(\hnu_n,r)\ge z,\ S_{n'}\ge n'q\big)\ \le\ \kn_ne^{2(1+2\eta)}e^{-(1+2\eta)z}\ \le\ \big(4+2\kap(r)\big)e^{2(1+2\eta)}\,n^2e^{-(1+2\eta)z}.
\]
\end{lemma}
\role{the probability that a rejected sub-optimal arm later looks good: Lemmas~\ref{lem:constA}(c), \ref{lem:controlC}(c).}
\noindent(see proof in Subsection~\ref{sec:pf-lem:pair}.)


\begin{lemma}[a good tilt exists]\label{lem:theta}
\uses{def:model, def:empkinf}
Assume $m<\m^\star<1$.
\begin{enumerate}
\item[(i)] There is $\theta>0$ with $\ \Phi^\star:=\Phis_\nu(\theta,\m^\star)<1\ $ and $\ \Psi^\star:=\Psis_\nu(\theta,\m^\star)<1-\m^\star$.
\item[(ii)] For such a $\theta$, let
\[
c_\theta:=\frac12\min\Big(\log\frac1{\Phi^\star},\ \log\frac{1-\m^\star}{\Psi^\star}\Big)>0,\qquad\varepsilon_\theta:=\min\Big(\frac{c_\theta}{6\theta},\ (1-\m^\star)\big(1-e^{-c_\theta/2}\big)\Big)>0 .
\]
Then for every $\varepsilon\in(0,\varepsilon_\theta]$,
\[
e^{c_\theta}\Phis_\nu(\theta,\m^\star-3\varepsilon)\ \le\ 1\qquad\text{and}\qquad e^{c_\theta}\Psis_\nu(\theta,\m^\star-3\varepsilon)\ \le\ 1-(\m^\star+\varepsilon).
\]
\end{enumerate}
\end{lemma}
\role{provides the tilt $\theta$ needed by Lemma~\ref{lem:pair}; used in Definitions~\ref{def:constA} and~\ref{def:constC}.}
\noindent(see proof in Subsection~\ref{sec:pf-lem:theta}.)


\begin{lemma}[lower deviation of the empirical $\Ki$]\label{lem:lowerdev}
\uses{def:kinf, def:stacks}
Let $\mu\in[0,1)$, $\varepsilon'>0$, $R':=\kap(\mu)-\log(1-e^{-\varepsilon'/4})$ and $\rho':=\varepsilon'^2/(2R'^2)$. Then for every $N\ge1$,
$\PP\big(\Ki(\hnu_N,\mu)\le\Ki(\nu,\mu)-\varepsilon'\big)\le e^{-\rho'N}$.
\end{lemma}
\role{the empirical divergence of $a$ is not too small: Lemmas~\ref{lem:probasA}(c), \ref{lem:controlC}(d), and the event $\KC$ of Lemma~\ref{lem:amb}.}
\noindent(see proof in Subsection~\ref{sec:pf-lem:lowerdev}.)


\begin{lemma}[upper deviation of the empirical $\Ki$]\label{lem:upperdev}
\uses{def:kinf, def:stacks}
With the notation of Lemma~\ref{lem:lowerdev}, let $\tau:=e^{-\varepsilon'/4}$ and $n_{\mathrm g}:=\lceil\frac8{\varepsilon'(1-\tau)(1-\mu)}\rceil+1$. Then for every $n\ge1$,
$\PP\big(\Ki(\hnu_n,\mu)\ge\Ki(\nu,\mu)+\varepsilon'\big)\le n_{\mathrm g}e^{-\rho'n}$.
\end{lemma}
\role{the empirical divergence of $a$ is not too large: Lemma~\ref{lem:probasA}(c) and the event $\KC$ of Lemma~\ref{lem:amb}.}
\noindent(see proof in Subsection~\ref{sec:pf-lem:upperdev}.)


\subsection{Reward stacks}

\begin{lemma}[stacks]\label{lem:stacks}
\uses{def:stacks, def:hy, def:grid, def:learner}
(i) We may assume that the bandit is built from the stacks: the variables $\Xs_{c,k}$ are independent, $\Xs_{c,k}\sim\nu_c$, $\Xs_{c,k}\in[0,1]$, and the $k$-th pull of $c$
receives $\Xs_{c,k}$. (ii) The run of \HY{} is then a deterministic function of the stacks; at every end $\tl_\ell$, $\hnu_c(\tl_\ell)=\hnu_{c,\Ng_c(\tl_\ell)}$ and
$\hm_c(\tl_\ell)=\hm_{c,\Ng_c(\tl_\ell)}$ are the statistics of the rewards of $c$ observed up to $\tl_\ell$. (iii) Every batched learner can be realized on the stacks and on an
independent internal randomization $U$, uniform on $[0,1]^{\mathbb N}$, so that, for every arm $a$ and every $\ell$, the actions up to round $\tl_{\ell+1}$ are a
measurable function of $U$, of the stacks of the arms $c\ne a$ and of $\Xs_{a,1},\dots,\Xs_{a,\Ng_a(\tl_\ell)}$, and the event $\{\Ng_a(\tl_\ell)=n\}$ only depends on $U$, on the
stacks of the arms $c\ne a$ and on $\Xs_{a,1},\dots,\Xs_{a,n}$; the counts have the same law as in the original run.
\end{lemma}
\role{lets every probability be computed on independent i.i.d.\ stacks; (iii) carries the change of measure of Section~\ref{sec:app-lower}.}
\noindent(see proof in Subsection~\ref{sec:pf-lem:stacks}.)


\noindent From now on, in the proofs of the upper bound, we fix an outcome of the stacks with all rewards in $[0,1]$ and work with the run of \HY{} on it.

\subsection{Elementary sums and sequences}\label{sec:app-sums}

\dt{s:sums}For $p>0$ and $\theta>0$ let $\sig_p(\theta):=\frac2\theta e^{\theta/2}\big(\frac{2p}{e\theta}\big)^p$.

\begin{lemma}[sums]\label{lem:sums}
\uses{def:grid}
\begin{enumerate}
\item[(i)] For $\theta>0$, $\sum_{\ell\ge2}\Bg_\ell\tl_\ell^{-(1+\theta)}\le\tl_1^{-\theta}/\theta\le1/\theta$.
\item[(ii)] For $p>0$ and $\theta>0$, $\sum_{\ell\ge2}\Bg_\ell(1+\Lg_\ell)^p\tl_\ell^{-(1+\theta)}\le\sig_p(\theta)$.
\item[(iii)] If $f\ge0$ is non-increasing on $[0,\infty)$, $\sum_{n\ge0}f(n)\le f(0)+\int_0^\infty f$; and for $c,p>0$, $\int_0^\infty e^{-cs^{1/p}}ds=\Gamma(p+1)c^{-p}$.
\item[(iv)] (counting) Let $i_0\le j$, $0=\lambda_{i_0-1}\le\lambda_{i_0}\le\dots\le\lambda_j$ reals, $v_i:=\lceil\lambda_i\rceil$, $\beta>0$ and $g\ge0$ non-increasing on $[0,\infty)$. Then
$\sum_{i=i_0}^j(v_i-v_{i-1})g(\lambda_i^{1/\beta})\le\sum_{n\ge0}g(n^{1/\beta})$.
\item[(v)] (layer cake) Let $0=v_{i_0-1}\le v_{i_0}\le\dots\le v_j$ and events $E_{i_0},\dots,E_j$. Then $\max_{i_0\le i\le j}v_i\1(E_i)\le\sum_{i=i_0}^j(v_i-v_{i-1})\1\big(\bigcup_{i'=i}^jE_{i'}\big)$.
\item[(vi)] For $\theta>0$, $\sum_{\ell\ge2}\Bg_\ell(\log\tl_\ell)^4\tl_\ell^{-(1+\theta)}\le\big(\frac8{e\theta}\big)^4\big(1+\frac2\theta\big)$ and $\sum_{\ell\ge2}\Bg_\ell\tl_\ell^{-(1+\theta)}\le1+\frac2\theta$.
\end{enumerate}
\end{lemma}
 
\role{sums of the probabilities weighted by the batch sizes: Lemmas~\ref{lem:sumsA}, \ref{lem:sumC} and Proposition~\ref{prop:dyA}.}
\noindent(see proof in Subsection~\ref{sec:pf-lem:sums}.)


\begin{lemma}[a sequence $\varepsilon_L\to0$]\label{lem:diag}
\uses{def:grid}
Let $\bar\varepsilon>0$, $f:(0,\bar\varepsilon]\to[0,\infty)$ and $(\tl_L)$ the grid. There is $(\varepsilon_L)_{L\ge1}$ in $(0,\bar\varepsilon]$ with $\varepsilon_L\to0$ and
$f(\varepsilon_L)/\log\tl_L\to0$.
\end{lemma}
\role{lets the accuracy $\varepsilon$ tend to $0$ with $L$, at the end of the proof of Theorem~\ref{thm:upper} ($x=1$).}
\noindent(see proof in Subsection~\ref{sec:pf-lem:diag}.)


\begin{lemma}[distinct counts]\label{lem:distinct}
\uses{def:learner, def:grid}
For a batch $j\ge2$ let $n_j(a):=\Ng_a(\tl_j)-\Ng_a(\tl_{j-1})$. (i) If $\Ng_a(\tl_1)\ge1$, then for $g\ge0$ and a property $Q$ of integers,
$\sum_{j=2}^Lg(\Ng_a(\tl_{j-1}))\1\big(Q(\Ng_a(\tl_{j-1})),\ n_j(a)>0\big)\le\sum_{n=1}^{\tl_L}g(n)\1(Q(n))$. (ii) If $\Ng_a(\tl_j)\le M_x$ for every $j\in[2,L]$ such that $P(j)$
holds and $n_j(a)>0$, then $\sum_{j=2}^Ln_j(a)\1(P(j))\le M_x$.
\end{lemma}
\role{counts the pulls of $a$ once over the batches, in Lemma~\ref{lem:decompC}.}
\noindent(see proof in Subsection~\ref{sec:pf-lem:distinct}.)


\begin{lemma}[chain]\label{lem:chain}
Let $Y\ge0$, $\delta'<1$, $D_{\max}>0$, $(w_j)$ non-decreasing reals and $S$ a finite non-empty set of integers with $0<D_j\le D_{\max}$ and $w_{j+1}D_j\ge Y$ for $j\in S$. Then either
some $j\in S$ has $w_jD_j\ge(1-\delta')Y$, or $w_{j'+1}\ge(1-\delta')^{-(|S|-1)}Y/D_{\max}$ for every $j'\ge\max S$.
\end{lemma}
\role{turns evidence at the end of high batches into evidence at their start, in Lemmas~\ref{lem:fewC} and~\ref{lem:optW}.}
\noindent(see proof in Subsection~\ref{sec:pf-lem:chain}.)

\subsection{Proof of Lemma~\ref{lem:kinf}}\label{sec:pf-lem:kinf}
\par\smallskip
\Setting $G$ on $[0,1]$ with mean $p$, $\mu\in[0,1)$; every item uses only the dual formula $\Ki(G,\mu)=\sup_{\lambda\in\Lset{\mu}}\E_G[\log(1-\lambda(X-\mu))]$.
\par\smallskip
\begin{proof}
\proves{lem:kinf}
\uses{def:kinf}
\stp{(i)}{range}$\lambda=0$ gives the value $0$, so $\Ki\ge0$. For $\lambda\in\Lset{\mu}$ and $y\in[0,1]$, $1-\lambda(y-\mu)\le1+\lambda\mu\le1+\frac\mu{1-\mu}=\frac1{1-\mu}$, so every
value is at most $\kap(\mu)$. And $\kap(\mu)=\log\frac1{1-\mu}$ is non-decreasing.

\stp{(ii)}{level below the mean}Since $\log(1+u)\le u$, for every $\lambda\ge0$,
\[
\E\log\big(1-\lambda(X-\mu)\big)\ \le\ -\lambda\,\E[X-\mu]\ =\ \lambda(\mu-p)\ \le\ 0 .
\]

\stp{(iii)}{monotonicity in the level}Let $0\le\mu\le\mu'<1$, $\lambda\in\Lset{\mu}$ and $\lambda':=\lambda\frac{1-\mu}{1-\mu'}\in\Lset{\mu'}$. The affine functions $y\mapsto1-\lambda(y-\mu)$
and $y\mapsto1-\lambda'(y-\mu')$:
\begin{itemize}
\item coincide at $y=1$ (both equal $1-\lambda(1-\mu)$);
\item at $y=0$ satisfy $1+\lambda\mu\le1+\lambda'\mu'$ (as $\lambda\le\lambda'$ and $\mu\le\mu'$).
\end{itemize}
So the first is below the second on $[0,1]$, and $\E\log(1-\lambda(X-\mu))\le\Ki(G,\mu')$. Take the supremum over $\lambda$.

\stp{(iv)}{a quadratic lower bound}For $u\ge-\frac12$, $\log(1+u)\ge\frac u{1+u}=u-\frac{u^2}{1+u}\ge u-2u^2$. Let $\lambda_0:=(\mu-q)/4\le\frac14$, so $\lambda_0\in\Lset{\mu}$, and
$u:=-\lambda_0(X-\mu)$, so $|u|\le\frac14$. Then
\begin{align*}
\Ki(G,\mu)\ \ge\ \E\big[\log(1+u)\big]\ &\ge\ \lambda_0(\mu-p)-2\lambda_0^2\qquad\why{$\E u=\lambda_0(\mu-p)$, $u^2\le\lambda_0^2$}\\
&\ge\ \lambda_0(\mu-q)-2\lambda_0^2\ =\ \frac{(\mu-q)^2}4-\frac{(\mu-q)^2}8\ =\ \frac{(\mu-q)^2}8\qquad\why{$p\le q$}.
\end{align*}

\stp{(v)}{shifting the level down}Let $\lambda\in\Lset{\mu}$, $\alpha:=\frac{1-\mu}{1-\mu+\varepsilon}\in(0,1]$ and $\lambda':=\alpha\lambda\in\Lset{\mu-\varepsilon}$. For $y\in[0,1]$,
\[
1-\lambda'(y-\mu+\varepsilon)-\alpha\big(1-\lambda(y-\mu)\big)\ =\ 1-\alpha-\alpha\lambda\varepsilon\ \ge\ 0,
\]
since $\alpha\lambda\varepsilon\le\frac\varepsilon{1-\mu+\varepsilon}=1-\alpha$. Hence $\Ki(G,\mu-\varepsilon)\ge\E\log(1-\lambda(X-\mu))+\log\alpha$, and
$\log\alpha=-\log(1+\frac\varepsilon{1-\mu})\ge-\frac\varepsilon{1-\mu}$. Take the supremum over $\lambda$.

\stp{(vi)}{scaling}$1-\tau\lambda(y-\mu)-\tau\big(1-\lambda(y-\mu)\big)=1-\tau\ge0$, and both sides are positive; take logarithms.
\end{proof}

\subsection{Proof of Lemma~\ref{lem:pseudo}}\label{sec:pf-lem:pseudo}
\begin{proof}
\proves{lem:pseudo}
\uses{def:kinf, def:hy}
\stp1{the test functions}For $\lambda\in\Lset{\mu}$ let $g_\lambda(y):=\log(1-\lambda(y-\mu))$ on $[0,1]$. It is bounded and non-increasing, and $g_\lambda(0)=\log(1+\lambda\mu)\le\kap(\mu)$. By
linearity,
\[
\E_{\tilde G}[g_\lambda]\ =\ \frac{N\,\E_G[g_\lambda]+\wz\,g_\lambda(0)}{N+\wz}.
\]
\stp{(i)}{}$g_\lambda(0)\ge\E_G[g_\lambda]$ ($g_\lambda$ is non-increasing), so $\E_{\tilde G}[g_\lambda]\ge\E_G[g_\lambda]$; take the supremum.

\stp{(ii)}{}$(N+\wz)\E_{\tilde G}[g_\lambda]\le N\Ki(G,\mu)+\wz\kap(\mu)$; take the supremum, and use $N\Ki(\tilde G,\mu)\le(N+\wz)\Ki(\tilde G,\mu)$.

\stp{(iii)}{}$p-\tilde p=\wz p/(N+\wz)\in[0,\wz/(N+\wz)]$.

\stp4{last claim}(iii) with $\wz/(N+\wz)\le\wz/N\le2\wz/\Ug_\ell$ when $N\ge\frac12\Ug_\ell$.
\end{proof}

\subsection{Proof of Lemma~\ref{lem:hoeffding}}\label{sec:pf-lem:hoeffding}
\begin{proof}
\proves{lem:hoeffding}
(a) is Hoeffding~\cite{H63} (to be checked). (b) By (a) with $y=n\rho$, $\PP(|S_n/n-\mu_Z|>\rho)\le2e^{-2\rho^2n}$. A union bound over the integers $n\ge\lceil u\rceil$ gives
\[
2\sum_{n\ge\lceil u\rceil}e^{-2\rho^2n}\ =\ \frac{2e^{-2\rho^2\lceil u\rceil}}{1-e^{-2\rho^2}}\ \le\ \frac{2e^{-2\rho^2u}}{1-e^{-2\rho^2}} .
\]
\end{proof}

\subsection{Proof of Lemma~\ref{lem:grid}}\label{sec:pf-lem:grid}
\par\smallskip
\Setting $n\ge1$, $z\ge2$, and $n\Ki(\widehat G,r)\ge z$.
\Goal a grid point $\lamk^{(k)}$, $1\le k\le\kn_n$, with $F(\lamk^{(k)})\ge z-2$.
\par\smallskip
\begin{proof}
\proves{lem:grid}
\uses{lem:kinf, def:kinf, def:empkinf}
\stp1{the grid}$\lamk^{(k)}(1-r)=\taun_n^k<1$ for $k\ge1$, so $\lamk^{(k)}\in\Lset{r}$. And $\kn_n\le2n(\log n+\kap(r))+2\le(4+2\kap(r))n^2$.

\stp2{a good $\lambda$ is not too small}Take $\lambda\in\Lset{r}$ with $F(\lambda)\ge z-1\ge1$. Since $\log(1+u)\le u$, $1\le F(\lambda)\le n\lambda r<n\lambda$, so $\lambda>\frac1n$.

\stp3{it lies between two grid points}As $\taun_n^{\kn_n}\le e^{-\kn_n/(2n)}\le\frac{1-r}n$, we have $\lamk^{(\kn_n)}\le\frac1n<\lambda<\lamk^{(0)}$. So
$\lambda\in(\lamk^{(k)},\lamk^{(k-1)}]$ for some $1\le k\le\kn_n$, and $\tau:=\lamk^{(k)}/\lambda\in[\taun_n,1)$.

\stp4{rounding down costs at most $1$}By Lemma~\ref{lem:kinf}(vi) and $\log(1-u)\ge-\frac u{1-u}$,
\[
F(\lamk^{(k)})\ \ge\ F(\lambda)+n\log\tau\ \ge\ z-1+n\log\Big(1-\frac1{2n}\Big)\ \ge\ z-1-\frac{1}{2-1/n}\ \ge\ z-2 .
\]
\end{proof}

\subsection{Proof of Lemma~\ref{lem:single}}\label{sec:pf-lem:single}
\noindent\emph{Idea.} A Chernoff bound at each point of the grid, with the power $1+2\eta$ instead of $1$; the extra power costs $e^{2\eta\lambda r}$, which the mean pays
when $2\eta r\le m-r$.
\begin{proof}
\proves{lem:single}
\uses{lem:grid, def:kinf, def:empkinf}
\stp1{reduction}Let $s:=1+2\eta$. If $z<2$, the bound is at least $1$. Let $z\ge2$. By Lemma~\ref{lem:grid} and a union bound over the $\kn_n$ grid points, it suffices to
show $\PP(F(\lambda)\ge z-2)\le e^{-s(z-2)}$ for each $\lambda\in\Lset{r}$, where $F(\lambda):=\sum_{k\le n}\log y_k$ and $y_k:=1-\lambda(Y_k-r)\in(0,1+\lambda r]$.

\stp2{one factor}Since $y_k^{2\eta}\le(1+\lambda r)^{2\eta}\le e^{2\eta\lambda r}$,
\[
\E[y_k^s]\ \le\ e^{2\eta\lambda r}\,\E[y_k]\ =\ e^{2\eta\lambda r}\big(1-\lambda(m-r)\big)\ \le\ e^{2\eta\lambda r-\lambda(m-r)}\ \le\ 1\qquad\why{$2\eta r\le m-r$}.
\]

\stp3{Markov}By Markov's inequality and independence, $\PP(F(\lambda)\ge z-2)\le e^{-s(z-2)}\prod_k\E[y_k^s]\le e^{-s(z-2)}$.
\end{proof}

\subsection{Proof of Lemma~\ref{lem:pair}}\label{sec:pf-lem:pair}
\noindent\emph{Idea.} The event that the empirical distribution is far from the level $r$ at count $n$ and that the mean climbs above $q$ at a later count $n'$ has
probability about $e^{-(1+2\eta)z}$, even when $r$ is above $m$; the factor $e^{\theta(Y_k-q)}$ pays for the climb.
\begin{proof}
\proves{lem:pair}
\uses{lem:single, lem:grid, def:kinf, def:empkinf}
\stp1{reduction}Let $s:=1+2\eta$ and $z\ge2$ (otherwise the bound is trivial). As in Lemma~\ref{lem:single}, by Lemma~\ref{lem:grid} it suffices to bound, for each
$\lambda\in\Lset{r}$, with $y_k:=1-\lambda(Y_k-r)$ and $F(\lambda):=\sum_{k\le n}\log y_k$, the probability of $E_\lambda:=\{F(\lambda)\ge z-2,\ S_{n'}\ge n'q\}$.

\stp2{Markov}On $E_\lambda$, $sF(\lambda)+\theta(S_{n'}-n'q)\ge s(z-2)$ ($\theta\ge0$). By Markov's inequality and independence,
\[
\PP(E_\lambda)\ \le\ e^{-s(z-2)}\,\underbrace{\E\big[y_1^se^{\theta(Y_1-q)}\big]}_{A}{}^{n}\,\underbrace{\Phis_\nu(\theta,q)}_{B}{}^{n'-n}.
\]

\stp3{the factor $A$}Since $y_1\le\frac1{1-r}$, $y_1^{2\eta}\le e^{2\eta\kap(r)}$. Moreover $y_1=(1-\lambda(1-r))+\lambda(1-Y_1)$, with both coefficients non-negative. So
\begin{align*}
A\ &\le\ e^{2\eta\kap(r)}\,\E\big[y_1e^{\theta(Y_1-q)}\big]\\
&=\ \big(1-\lambda(1-r)\big)\,e^{2\eta\kap(r)}\Phis_\nu(\theta,q)+\lambda\,e^{2\eta\kap(r)}\Psis_\nu(\theta,q)&&\why{definitions of $\Phis$, $\Psis$}\\
&\le\ \big(1-\lambda(1-r)\big)+\lambda(1-r)\ =\ 1&&\why{assumptions}.
\end{align*}

\stp4{the factor $B$}$B\le e^{-2\eta\kap(r)}\le1$. Hence $\PP(E_\lambda)\le e^{-s(z-2)}$.
\end{proof}

\subsection{Proof of Lemma~\ref{lem:theta}}\label{sec:pf-lem:theta}
\noindent\emph{Idea.} $\theta\mapsto\Phis_\nu(\theta,\m^\star)$ equals $1$ at $0$ with derivative $m-\m^\star<0$; at its minimizer the derivative vanishes, which also
makes $\Psis_\nu$ small; if $\nu$ has no mass above $\m^\star$, let $\theta\to\infty$. A margin $c_\theta$ survives the shift of the level to $\m^\star-3\varepsilon$ for
$\varepsilon$ small.
\begin{proof}
\proves{lem:theta}
\uses{def:model, def:empkinf}
\stp0{notation}Let $Z:=Y_1-\m^\star\in[-1,1]$ and $\varphi(\theta):=\E e^{\theta Z}$. Then $\E Z=m-\m^\star<0$, $\Phis_\nu(\theta,\m^\star)=\varphi(\theta)$ and
$\Psis_\nu(\theta,\m^\star)=\E[(1-\m^\star-Z)e^{\theta Z}]$.

\stp1{case $\PP(Y_1>\m^\star)=0$}Then $Z\le0$, and by dominated convergence, as $\theta\to\infty$, $\varphi(\theta)\to p_0:=\PP(Z=0)$ and $\Psis_\nu(\theta,\m^\star)\to(1-\m^\star)p_0$.
Since $Z\ge-\1(Z\ne0)$, $m-\m^\star\ge-(1-p_0)$, so $p_0<1$. Any $\theta$ large enough works.

\stp2{case $\PP(Y_1>\m^\star)>0$}Let $z_0>0$ with $p:=\PP(Z\ge z_0)>0$, and $\theta_s:=\frac{\m^\star-m}2\in(0,\frac12]$.
\begin{itemize}
\item \emph{$\varphi$ dips below $1$.} $\varphi(\theta_s)\le1+\theta_s(m-\m^\star)+\theta_s^2=1-\theta_s^2<1$. \ (Since $e^u\le1+u+u^2$ for $|u|\le1$.)
\item \emph{$\varphi$ comes back above $1$.} Let $\theta_{\max}:=\max(1,\log(1/p)/z_0)$. Then $\varphi(\theta_{\max})\ge pe^{\theta_{\max}z_0}\ge1$.
\item \emph{An interior minimizer.} A minimizer $\theta$ of the continuous $\varphi$ on $[0,\theta_{\max}]$ has $\varphi(\theta)\le\varphi(\theta_s)<1\le\min(\varphi(0),\varphi(\theta_{\max}))$ (as
$\theta_s\le\frac12<\theta_{\max}$), so it is interior, and $\varphi'(\theta)=\E[Ze^{\theta Z}]=0$.
\end{itemize}
Hence
\[
\Phi^\star=\varphi(\theta)<1,\qquad\Psi^\star=(1-\m^\star)\varphi(\theta)-\E[Ze^{\theta Z}]=(1-\m^\star)\varphi(\theta)<1-\m^\star .
\]

\stp3{the shift of the level}
\begin{itemize}
\item $\Phi^\star,\Psi^\star>0$ (as $\PP(Y_1<1)>0$), so $c_\theta>0$, and by definition of $c_\theta$: $\ \Phi^\star\le e^{-2c_\theta}$, $\ \Psi^\star\le(1-\m^\star)e^{-2c_\theta}$.
\item For $\varepsilon\le\varepsilon_\theta$: $3\theta\varepsilon\le c_\theta/2$.
\item Shifting the level multiplies both moments by $e^{3\theta\varepsilon}$: $\Phis_\nu(\theta,\m^\star-3\varepsilon)=e^{3\theta\varepsilon}\Phi^\star$ and
$\Psis_\nu(\theta,\m^\star-3\varepsilon)=e^{3\theta\varepsilon}\Psi^\star$.
\end{itemize}
So
\[
e^{c_\theta}\Phis_\nu(\theta,\m^\star-3\varepsilon)\ \le\ e^{-c_\theta/2}\ \le\ 1,\qquad
e^{c_\theta}\Psis_\nu(\theta,\m^\star-3\varepsilon)\ \le\ (1-\m^\star)e^{-c_\theta/2}\ \le\ 1-\m^\star-\varepsilon,
\]
the last inequality because $\varepsilon\le(1-\m^\star)(1-e^{-c_\theta/2})$.
\end{proof}

\subsection{Proof of Lemma~\ref{lem:lowerdev}}\label{sec:pf-lem:lowerdev}
\begin{proof}
\proves{lem:lowerdev}
\uses{lem:kinf, lem:hoeffding, def:kinf, def:stacks}
\stp1{a fixed test function}Let $\mathcal K:=\Ki(\nu,\mu)<\infty$, and take $\lambda\in\Lset{\mu}$ with $\E\log(1-\lambda(Y_1-\mu))\ge\mathcal K-\varepsilon'/4$. Let $\tau:=e^{-\varepsilon'/4}$ and
$g(y):=\log(1-\tau\lambda(y-\mu))$.
\begin{itemize}
\item Since $1-\tau<1-\tau\lambda(y-\mu)\le\frac1{1-\mu}$, $g$ takes its values in an interval of length $R'$.
\item By Lemma~\ref{lem:kinf}(vi), $\E g(Y_1)\ge\log\tau+\mathcal K-\varepsilon'/4=\mathcal K-\varepsilon'/2$.
\end{itemize}
\stp2{reduction to a sum}As $\tau\lambda\in\Lset{\mu}$, $N\Ki(\hnu_N,\mu)\ge\sum_{k\le N}g(Y_k)$. So the event implies
\[
\sum_{k\le N}g(Y_k)\ \le\ N(\mathcal K-\varepsilon')\ \le\ N\big(\E g(Y_1)-\varepsilon'/2\big).
\]
\stp3{Hoeffding}Lemma~\ref{lem:hoeffding}(a) applied to $(g(Y_k)-\log(1-\tau))/R'\in[0,1]$ bounds this probability by $e^{-2N(\varepsilon'/2)^2/R'^2}=e^{-\rho'N}$.
\end{proof}

\subsection{Proof of Lemma~\ref{lem:upperdev}}\label{sec:pf-lem:upperdev}
\begin{proof}
\proves{lem:upperdev}
\uses{lem:kinf, def:kinf}
\stp1{a finite net}Let $\mathcal K:=\Ki(\nu,\mu)$, $F(\lambda):=\sum_{k\le n}\log(1-\lambda(Y_k-\mu))$, $\eta':=\varepsilon'(1-\tau)/8$ and
$\lambda_j:=\min(j\eta',\frac\tau{1-\mu})$ for $0\le j\le n_{\mathrm g}-1$. Every $\lambda'\in[0,\frac\tau{1-\mu}]$ is within $\eta'$ of some $\lambda_j$.

\stp2{$F$ is Lipschitz on $[0,\frac\tau{1-\mu}]$}There the derivative of $\lambda\mapsto\log(1-\lambda(y-\mu))$ is at most $\frac1{1-\tau}$ in absolute value, so
$|F(\lambda)-F(\lambda')|\le n|\lambda-\lambda'|/(1-\tau)$.

\stp3{the event forces a large value on the net}On the event, take $\lambda$ with $F(\lambda)\ge n(\mathcal K+\frac78\varepsilon')$, $\lambda':=\tau\lambda$ and $j$ with
$|\lambda_j-\lambda'|\le\eta'$. Then
\[
F(\lambda_j)\ \ge\ F(\lambda')-\frac{n\eta'}{1-\tau}\ \ge\ F(\lambda)+n\log\tau-\frac{n\varepsilon'}8\ \ge\ n\Big(\mathcal K+\frac{\varepsilon'}2\Big)\qquad\why{Step~2; Lemma~\ref{lem:kinf}(vi)}.
\]
\stp4{Hoeffding at each point of the net}For a fixed $j$, $Z_k:=\log(1-\lambda_j(Y_k-\mu))\in[\log(1-\tau),\kap(\mu)]$, an interval of length $R'$, and $\E Z_k\le\mathcal K$. Hoeffding's
inequality gives $\PP(\sum_kZ_k\ge n(\mathcal K+\varepsilon'/2))\le e^{-\rho'n}$; a union bound over the $n_{\mathrm g}$ points concludes.
\end{proof}

\subsection{Proof of Lemma~\ref{lem:stacks}}\label{sec:pf-lem:stacks}
\begin{proof}
\proves{lem:stacks}
\uses{def:hy, def:grid, def:stacks, def:learner}
\stp{(i)}{stacks}Given the past, the reward of the $k$-th pull of $c$ has law $\nu_c$ and is independent of the past. Building the bandit from stacks drawn in advance gives
histories with the same law, by induction on the rounds \cite[Section~4.6]{LS} (to be checked).

\stp{(ii)}{\HY{} is a function of the stacks}Induction on the batches and on the rounds: the arm pulled at a round of batch $\ell+1$ is a function of $\tl_1,\dots,\tl_\ell$, of
the counts and samples frozen at $\tl_{\ell-1}$ and $\tl_\ell$ (through $\Kh_{\ell-1}$ and $\Ld_{\ell-1}$ in $\Qp_\ell$ and $\Lh_\ell$), of the current counts and of $t$. The
rewards of $c$ observed up to $\tl_\ell$ are $\Xs_{c,1},\dots,\Xs_{c,\Ng_c(\tl_\ell)}$.

\stp{(iii)}{any batched learner}Realize each policy, by the inverse distribution function and a coordinate of $U$, on the history in which the rewards of the current batch
are replaced by $0$. Since the learner is batched this does not change the policy, and the counts have the law of the original run (Ionescu-Tulcea). By induction on the
rounds, in batch $\ell'\le\ell+1$, $A_t$ is a function of $U$, of the past actions and of the rewards of the rounds $s\le\tl_{\ell'-1}\le\tl_\ell$. Such a reward is
$\Xs_{b,\Ng_b(s)}$ with $b=A_s$, and if $b=a$ then $\Ng_a(s)\le\Ng_a(\tl_\ell)$. The second claim follows by the same induction up to $\tl_\ell$, on $\{\Ng_a(\tl_\ell)=n\}$.
\end{proof}

\subsection{Proof of Lemma~\ref{lem:sums}}\label{sec:pf-lem:sums}
\begin{proof}
\proves{lem:sums}
\uses{def:grid}
\stp{(i)}{}On $[\tl_{\ell-1},\tl_\ell]$, an interval of length $\Bg_\ell$, $s^{-(1+\theta)}\ge\tl_\ell^{-(1+\theta)}$; so $\Bg_\ell\tl_\ell^{-(1+\theta)}\le\int_{\tl_{\ell-1}}^{\tl_\ell}s^{-(1+\theta)}ds$.
Summing over $\ell\ge2$ gives at most $\int_{\tl_1}^\infty s^{-(1+\theta)}ds=\tl_1^{-\theta}/\theta$.

\stp{(ii)}{}For $y\ge0$, $(1+y)^pe^{-\theta y/2}\le e^{\theta/2}(\frac{2p}{e\theta})^p$ (the maximum of $s^pe^{-\theta s/2}$ is at $s=2p/\theta$). With $y=\Lg_\ell$, each term is at
most $e^{\theta/2}(\frac{2p}{e\theta})^p\Bg_\ell\tl_\ell^{-(1+\theta/2)}$; apply (i) with $\theta/2$.

\stp{(iii)}{}$f(n)\le\int_{n-1}^nf$ for $n\ge1$. For the integral, substitute $s=(u/c)^p$.

\stp{(iv)}{}$v_i-v_{i-1}$ is the number of integers $n\in[\lambda_{i-1},\lambda_i)$, and for such $n$, $g(\lambda_i^{1/\beta})\le g(n^{1/\beta})$ ($g$ is non-increasing). The intervals are
disjoint.

\stp{(v)}{}If the left side is $v_{i^\ast}>0$ with $E_{i^\ast}$, every union $\bigcup_{i'\ge i}E_{i'}$, $i\le i^\ast$, contains $E_{i^\ast}$. So the right side is at least
$\sum_{i\le i^\ast}(v_i-v_{i-1})=v_{i^\ast}$.

\stp{(vi)}{}$y^4e^{-\theta y/2}\le(\frac8{e\theta})^4$, so each term of the first sum is at most $(\frac8{e\theta})^4\Bg_\ell\tl_\ell^{-(1+\theta/2)}$. And
$\sum_{\ell\ge2}\Bg_\ell\tl_\ell^{-(1+\theta/2)}\le\sum_{u\ge1}u^{-(1+\theta/2)}\le1+\frac2\theta$.
\end{proof}

\subsection{Proof of Lemma~\ref{lem:diag}}\label{sec:pf-lem:diag}
\begin{proof}
\proves{lem:diag}
\uses{def:grid}
Let $e_k:=\bar\varepsilon2^{-k}$, and $L^{(k)}$ the smallest $L>L^{(k-1)}$ with $\log\tl_L\ge kf(e_k)$ ($L^{(0)}:=0$); it exists since $\tl_L\to\infty$. Put $\varepsilon_L:=e_1$ for
$L<L^{(1)}$, and $\varepsilon_L:=e_k$ for $L\in[L^{(k)},L^{(k+1)})$. On this block $f(\varepsilon_L)/\log\tl_L\le1/k$, and $k\to\infty$ with $L$.
\end{proof}

\subsection{Proof of Lemma~\ref{lem:distinct}}\label{sec:pf-lem:distinct}
\begin{proof}
\proves{lem:distinct}
\uses{def:learner, def:grid}
(i) The batches $j\ge2$ in which $a$ is pulled start at distinct counts $\Ng_a(\tl_{j-1})\in[\Ng_a(\tl_1),\tl_L]\subset[1,\tl_L]$, since the count is
non-decreasing and increases strictly from one such batch to the next. (ii) Let
$j^\ast$ be the last batch with $P(j^\ast)$ and $n_{j^\ast}(a)>0$. The sum is at most $\Ng_a(\tl_{j^\ast})-\Ng_a(\tl_1)\le M_x$.
\end{proof}

\subsection{Proof of Lemma~\ref{lem:chain}}\label{sec:pf-lem:chain}
\begin{proof}
\proves{lem:chain}
Induction on $|S|$, adding a new maximum $k$.

\noindent\emph{$S=\{k\}$.} $w_{k+1}D_k\ge Y$ and $D_k\le D_{\max}$ give $w_{k+1}\ge Y/D_{\max}$, and $(w_j)$ is non-decreasing.

\noindent\emph{Induction step.} Assume that no element of $S$ has $w_jD_j\ge(1-\delta')Y$. The induction hypothesis on $S\setminus\{k\}$, at $j'=k-1$, gives
$w_k\ge(1-\delta')^{-(|S|-2)}Y/D_{\max}$. And $w_kD_k<(1-\delta')Y\le(1-\delta')w_{k+1}D_k$, so $w_k<(1-\delta')w_{k+1}$. Hence, for $j'\ge k$,
\[
w_{j'+1}\ \ge\ w_{k+1}\ >\ \frac{w_k}{1-\delta'}\ \ge\ (1-\delta')^{-(|S|-1)}\frac Y{D_{\max}} .
\]
\end{proof}

%======================================================================
\section{Proof of the lower bound (Theorem~\ref{thm:lower})}\label{sec:app-lower}

\Plan The proof has two independent parts.
\begin{itemize}
\item \emph{One batch ahead} (Subsection~\ref{sec:app-lb-one}, Lemma~\ref{lem:onebatch}): the classical change of measure, which gives the case $x=1$.
\item \emph{Two batches ahead} (Subsection~\ref{sec:app-lb-two}, Lemmas~\ref{lem:tilts}--\ref{lem:twobatches}): a change of measure that tilts only the rewards of $a$ in
batch $\ell$, which gives the factor $x^2$ when $x>1$.
\end{itemize}
Subsection~\ref{sec:app-lb-assembly} assembles the two parts; the proofs of the lemmas follow.

\par\smallskip
\Setting (fixed in the whole section)
\begin{itemize}
\item an instance $\nu\in\cP$ and a sub-optimal arm $a$;
\item a batched learner, consistent on $\cP$, realized on the stacks and on an internal randomization $U$ as in Lemma~\ref{lem:stacks}(iii).
\end{itemize}
\textbf{Notation.}
\begin{itemize}
\item For a distribution $\rho$ on $[0,1]$ with mean smaller than $1$, $\nu^\rho$ is the instance $\nu$ in which $\nu_a$ is replaced by $\rho$; $\PP_\rho$ and $\E_\rho$ are
the law and the expectation of the run on $\nu^\rho$.
\item $\mathcal H_\ell$ is the $\sigma$-field generated by $U$ and by the observations up to round $\tl_\ell$.
\item $V:=\big(U,(\Xs_{c,k})_{c\ne a,\,k\ge1}\big)$; its law $P_V$ is the same on all the instances $\nu^\rho$.
\end{itemize}
\par\smallskip

\noindent\emph{A consequence of batching, used throughout.} The actions of batch $\ell+1$ do not depend on the rewards of batch $\ell+1$. Hence
\[
\text{the actions up to round }\tl_{\ell+1},\ \text{and in particular every count }\Ng_c(\tl_{\ell+1}),\ \text{are }\mathcal H_\ell\text{-measurable.}
\]

\noindent\emph{Two classical tools}, cited (no proof here).

\begin{lemma}[{primal formula \cite[Theorem~8]{HT10}}]\label{cite:honda2010asymptotically-primal}
\uses{def:kinf}
With $\mathcal A=\{G:\mathrm{supp}(G)\subset[0,1]\}$,
$D_{\min}(F,\mu;\mathcal A)\equiv\inf_{G\in\mathcal A:\mathrm E(G)\ge\mu}D(F\|G)$,
$H(\nu;F,\mu)\equiv\mathrm E_F[\log(1-(X-\mu)\nu)]$ and $D'_{\min}(F,\mu)\equiv\max_{0\le\nu\le\frac1{1-\mu}}H(\nu;F,\mu)$:
$D_{\min}(F,\mu)=D'_{\min}(F,\mu)$ holds for all $F\in\mathcal A$ and $\mu<1$.
\end{lemma}
\noindent\emph{How it is used.} Here $D(F\|G)=\KL(F,G)$. For $\mu\in[0,1)$, $D'_{\min}(F,\mu)=\Ki(F,\mu)$: the function
$\nu\mapsto H(\nu;F,\mu)$ is concave on $[0,\frac1{1-\mu}]$, so its value at the right end is at most its limit from the left,
and the maximum over the closed interval equals the supremum over $\Lset{\mu}$ of Definition~\ref{def:kinf}. Hence, for $G$
on $[0,1]$ and $\mu\in[0,1)$,
\[
\Ki(G,\mu)=\inf\big\{\KL(G,H):\ H\text{ on }[0,1],\ \E_H[X]\ge\mu\big\}.
\]
The strict form used in Lemma~\ref{lem:onebatch} (mean $>\mu$) follows from this one at the level $\mu+\epsilon$ and from
Lemma~\ref{lem:kinf}(v); it is derived in Step~1 of that proof. (Lemma~2 of \cite{HT10} states the strict form, but as printed
only for $\mu<0$; it is not used here.) This node is used only in Lemma~\ref{lem:onebatch}.

\begin{lemma}[{data processing \cite[Lemma~1]{GMS19}}]\label{cite:garivier2019explore-dataprocessing}
Consider a measurable space $(\Gamma,\mathcal G)$ equipped with two distributions
$\mathbb P_1$ and $\mathbb P_2$, and any $\mathcal G$--measurable random variable $Z:\Omega\to[0,1]$. We denote
respectively by $\mathbb E_1$ and $\mathbb E_2$ the expectations under $\mathbb P_1$ and $\mathbb P_2$. Then,
$\mathrm{KL}(\mathbb P_1,\mathbb P_2)\ge\mathrm{kl}(\mathbb E_1[Z],\mathbb E_2[Z])$.
\end{lemma}
\noindent Here $\kl$ is the Bernoulli divergence, and ``$Z:\Omega\to[0,1]$'' is as printed (arXiv:1602.07182v3, Section~2.1;
$\Gamma$ is meant). It is used with $(\Gamma,\mathcal G)$ the space of an observation $O$, $\mathbb P_1=P^O$ and
$\mathbb P_2=Q^O$: if $Z\in[0,1]$ is a measurable function of $O$, then $\KL(P^O,Q^O)\ge\kl(\E_P[Z],\E_Q[Z])$.

%----------------------------------------------------------------------
\subsection{One batch ahead}\label{sec:app-lb-one}

\begin{lemma}[one batch ahead]\label{lem:onebatch}
\uses{def:grid, def:model, def:learner, def:kinf}
For every grid (with $\Bg_1\ge2$), every batched learner consistent on $\cP$, every $\nu\in\cP$ and every sub-optimal arm $a$,
\[
\liminf_{\ell\to\infty}\frac{\E_\nu[\Ng_a(\tl_\ell)]}{\log\tl_{\ell+1}}\ \ge\ \frac1{\Ki(\nu_a,\m^\star)} .
\]
\end{lemma}
\role{gives Theorem~\ref{thm:lower}(i), the case $x=1$.}
\noindent(see proof in Subsection~\ref{sec:pf-lem:onebatch}.)

%----------------------------------------------------------------------
\subsection{Two batches ahead}\label{sec:app-lb-two}

\Plan Four steps, one lemma each.
\begin{enumerate}
\item[(1)] Lemma~\ref{lem:tilts}: a \emph{boundary tilt} $G$ of $\nu_a$ has mean exactly $\m^\star$ and $\KL(\nu_a,G)\approx\da$; mixing it with a small atom at $0$ or at $1$
gives two instances, one where $a$ is barely sub-optimal and one where $a$ is barely optimal.
\item[(2)] Lemma~\ref{lem:com}: a measure $\QQ$ under which the rewards of $a$ after $\tl_{\ell-1}$ are drawn from $G$ compares with both instances (change of measure on the stack).
\item[(3)] Lemma~\ref{lem:typ}: under $\QQ$ the log-likelihood ratios take their typical values, except on an event of small probability that does not depend on $\ell$.
\item[(4)] Lemma~\ref{lem:twobatches}: if $\Ng_a(\tl_{\ell-1})$ were below $\log\tl_{\ell+1}/\da$, the learner could not tell the two instances apart before batch $\ell+1$,
and would lose polynomially under one of them.
\end{enumerate}

\begin{definition}[boundary tilt]\label{def:tilt}\dt{s:tiltb}
\uses{def:model}
A distribution $G$ on $[0,1]$ is a \emph{boundary tilt} of $\nu_a$ if
\[
\E_G[X]=\m^\star\qquad\text{and}\qquad G=h\,\nu_a+c_G\,\delta_1,
\]
with $c_G\ge0$ and $e^{-C}\le h\le e^C$ $\nu_a$-a.s.\ for some $C<\infty$. For $s\in(0,1)$, let
\[
\rhom:=(1-s)G+s\,\delta_0\qquad\text{and}\qquad\rhop:=(1-s)G+s\,\delta_1 .
\]
\end{definition}

\begin{lemma}[boundary tilts]\label{lem:tilts}
\uses{def:kinf, def:tilt, def:model}
\begin{enumerate}
\item[(a)] $\inf\{\KL(\nu_a,G):\ G\text{ boundary tilt of }\nu_a\}=\da$, and $\KL(\nu_a,G)\ge\da>0$ for every boundary tilt $G$.
\item[(b)] Let $G$ be a boundary tilt and $s\in(0,1)$. Then:
\begin{itemize}
\item $\E_{\rhom}[X]<\m^\star<\E_{\rhop}[X]<1$;
\item for $\rho\in\{\rhom,\rhop\}$, $\nu_a\ll\rho$ and $G\ll\rho$, and the functions $\log\frac{d\nu_a}{d\rho}$ ($\nu_a$-a.s.) and $\log\frac{dG}{d\rho}$ ($G$-a.s.) take
their values in intervals of length at most some $D_G<\infty$;
\item for $\rho\in\{\rhom,\rhop\}$,
\[
\KL(\nu_a,\rho)\ \le\ \KL(\nu_a,G)+\log\frac1{1-s}\qquad\text{and}\qquad\KL(G,\rho)\ \le\ \log\frac1{1-s}.
\]
\end{itemize}
\end{enumerate}
\end{lemma}
\role{provides the two instances of the proof of Lemma~\ref{lem:twobatches} (Step~1) and the bounds on their divergences (Step~2).}
\noindent(see proof in Subsection~\ref{sec:pf-lem:tilts}.)

\begin{definition}[the tilted measure]\label{def:Q}\dt{s:Q}
\uses{def:tilt, def:learner, def:stacks, lem:stacks}
Fix a boundary tilt $G$, $s\in(0,1)$ and $\ell\ge2$. Let $(Y_i)_{i\ge1}$ be i.i.d.\ with law $G$, independent of everything else.
\begin{itemize}
\item Let $\nzero:=\Ng_a(\tl_{\ell-1})$, the number of pulls of $a$ before batch $\ell$.
\item $\QQ=\QQ^G_\ell$ is the law of the run in which the learner is run on $\nu$, except that the stack of $a$ is
\[
(\Xs_{a,1},\dots,\Xs_{a,\nzero},\,Y_1,Y_2,\dots):
\]
every reward of $a$ received after round $\tl_{\ell-1}$ is drawn from $G$. This is well defined: by Lemma~\ref{lem:stacks}(iii), the run up to $\tl_\ell$, rewards of batch
$\ell$ excluded, only uses the first $\nzero$ rewards of $a$.
\item Let $\Mb:=\Ng_a(\tl_\ell)-\nzero$, the number of pulls of $a$ in batch $\ell$. It is $\mathcal H_{\ell-1}$-measurable.
\item Let $Z_j$ be the $j$-th reward of $a$ in the run. For $\rho\in\{\rhom,\rhop\}$, the log-likelihood ratios of the rewards of $a$ before and inside batch $\ell$ are
\[
\Sl_1^\rho:=\sum_{j\le\nzero}\log\frac{d\nu_a}{d\rho}(Z_j),\qquad\Sl_2^\rho:=\sum_{\nzero<j\le\nzero+\Mb}\log\frac{dG}{d\rho}(Z_j),\qquad\Sl^\rho:=\Sl_1^\rho+\Sl_2^\rho .
\]
\end{itemize}
\end{definition}

\begin{lemma}[change of measure on the stack]\label{lem:com}
\uses{def:tilt, def:Q}
Let $\rho\in\{\rhom,\rhop\}$.
\begin{enumerate}
\item[(i)] For every $F\in\mathcal H_\ell$, $\ \PP_\rho(F)\ \ge\ \E_\QQ\big[\1_F\,e^{-\Sl^\rho}\big]$.
\item[(ii)] For every $F\in\mathcal H_{\ell-1}$, $\ \QQ(F)=\PP_\nu(F)$.
\end{enumerate}
\end{lemma}
\role{compares the tilted measure with the two instances, in Step~3 of the proof of Lemma~\ref{lem:twobatches}.}
\noindent(see proof in Subsection~\ref{sec:pf-lem:com}.)

\begin{lemma}[typical values]\label{lem:typ}
\uses{def:tilt, lem:tilts, def:Q}
Let $D_G$ be as in Lemma~\ref{lem:tilts}(b), $\varepsilon_H>0$ and $b>0$. Let $T_b$ be the event on which, for both $\rho\in\{\rhom,\rhop\}$,
\[
\Sl_1^\rho\ \le\ \nzero\big(\KL(\nu_a,\rho)+\varepsilon_H\big)+b\qquad\text{and}\qquad\Sl_2^\rho\ \le\ \Mb\big(\KL(G,\rho)+\varepsilon_H\big)+b .
\]
Then
\[
\QQ(T_b^c)\ \le\ \eta(b):=\frac{4\,e^{-4\varepsilon_Hb/D_G^2}}{1-e^{-2\varepsilon_H^2/D_G^2}} ,
\]
which does not depend on $\ell$ and tends to $0$ as $b\to\infty$.
\end{lemma}
\role{controls the log-likelihood ratios, in Steps~2 and~3 of the proof of Lemma~\ref{lem:twobatches}.}
\noindent(see proof in Subsection~\ref{sec:pf-lem:typ}.)

\begin{lemma}[two batches ahead]\label{lem:twobatches}
\uses{def:grid, def:learner, def:kinf}
Assume $\liminf_\ell\Bg_\ell/\tl_\ell>0$.
\begin{enumerate}
\item[(i)] If the number of pulls of $a$ in batch $\ell$ is tight at the scale $\log\tl_{\ell+1}$, that is
\[
\lim_{C\to\infty}\ \limsup_{\ell\to\infty}\ \PP_\nu\big(\Ng_a(\tl_\ell)-\Ng_a(\tl_{\ell-1})>C\log\tl_{\ell+1}\big)=0,
\]
then for every $\varepsilon\in(0,1)$,
\[
\PP_\nu\Big(\Ng_a(\tl_{\ell-1})\le(1-\varepsilon)\frac{\log\tl_{\ell+1}}{\da}\Big)\ \xrightarrow[\ell\to\infty]{}\ 0 .
\]
\item[(ii)] If $\E_\nu[\Ng_a(\tl_\ell)]=O(\log\tl_{\ell+1})$, then
\[
\liminf_{\ell\to\infty}\frac{\E_\nu[\Ng_a(\tl_{\ell-1})]}{\log\tl_{\ell+1}}\ \ge\ \frac1{\da} .
\]
\end{enumerate}
\end{lemma}
\role{gives Theorem~\ref{thm:lower}(ii), the factor $x^2$ when $x>1$.}
\noindent(see proof in Subsection~\ref{sec:pf-lem:twobatches}.)

%----------------------------------------------------------------------
\subsection{Proof of Theorem~\ref{thm:lower}}\label{sec:app-lb-assembly}
\begin{proof}
\proves{thm:lower}
\uses{lem:onebatch, lem:twobatches, def:learner, def:grid, def:kinf, def:model}
\paragraph{(i), $x=1$.} For every $\ell$,
\[
\frac{\E[\Ng_a(\tl_\ell)]}{\log\tl_\ell}\ =\ \frac{\E[\Ng_a(\tl_\ell)]}{\log\tl_{\ell+1}}\cdot\xg_{\ell+1} .
\]
The first factor has $\liminf\ge1/\da$ by Lemma~\ref{lem:onebatch}, and $\xg_{\ell+1}\to1$. Hence $\liminf_\ell\E[\Ng_a(\tl_\ell)]/\log\tl_\ell\ge1/\da$.

\noindent\emph{Every grid.} Since $\xg_{\ell+1}\ge1$, the same identity gives $\E[\Ng_a(\tl_\ell)]/\log\tl_\ell\ge\E[\Ng_a(\tl_\ell)]/\log\tl_{\ell+1}$ on every grid, convergent or not,
which is the remark in parentheses.

\paragraph{(ii), $x>1$.} \emph{Lemma~\ref{lem:twobatches} applies.} Since $\tl_{\ell-1}=\tl_\ell^{1/\xg_\ell}$,
\[
\frac{\Bg_\ell}{\tl_\ell}\ =\ 1-\frac{\tl_{\ell-1}}{\tl_\ell}\ =\ 1-e^{-\Lg_\ell(1-1/\xg_\ell)}\ \xrightarrow[\ell\to\infty]{}\ 1,
\]
because $\Lg_\ell\to\infty$ and $1-1/\xg_\ell\to1-1/x>0$.

\noindent\emph{Case A: $\E[\Ng_a(\tl_\ell)]=O(\log\tl_\ell)$.} Then $\E[\Ng_a(\tl_{\ell+1})]=O(\log\tl_{\ell+1})$, hence $O(\log\tl_{\ell+2})$. Lemma~\ref{lem:twobatches}(ii), applied
at the index $\ell+1$, gives
\[
\liminf_{\ell\to\infty}\frac{\E[\Ng_a(\tl_\ell)]}{\log\tl_{\ell+2}}\ \ge\ \frac1{\da} .
\]
Since $\log\tl_{\ell+2}/\log\tl_\ell=\xg_{\ell+1}\xg_{\ell+2}\to x^2$, this gives $\liminf_\ell\E[\Ng_a(\tl_\ell)]/\log\tl_\ell\ge x^2/\da$: this is the $\liminf$ claim, and it implies
the $\limsup$ claim.

\noindent\emph{Case B: otherwise.} Then $\limsup_\ell\E[\Ng_a(\tl_\ell)]/\log\tl_\ell=\infty$, and the $\limsup$ claim holds.

\paragraph{(iii), regret.} By definition of the regret,
\[
\frac{\Reg(\tl_\ell)}{\log\tl_\ell}\ =\ \sum_{c:\Dg_c>0}\Dg_c\,\frac{\E[\Ng_c(\tl_\ell)]}{\log\tl_\ell} .
\]
The $\liminf$ of a finite sum of non-negative sequences is at least the sum of the $\liminf$s. For each sub-optimal arm $c$, the $\liminf$ of $\E[\Ng_c(\tl_\ell)]/\log\tl_\ell$
is at least $x^2/\Ki(\nu_c,\m^\star)$: by (i) if $x=1$, and by Case~A of (ii) if $x>1$ (the counts of every sub-optimal arm are assumed logarithmic). \hfill$\square$

%----------------------------------------------------------------------
\end{proof}

\subsection{Proof of Lemma~\ref{lem:onebatch}}\label{sec:pf-lem:onebatch}
\noindent\emph{Idea.} The observations available at the end of batch $\ell+1$ contain the actions of batch $\ell+1$ but only the rewards of the first $\ell$ batches. A
learner that gives batch $\ell+1$ to the right arm under both $\nu$ and an alternative where $a$ is optimal must have paid $\Ng_a(\tl_\ell)\da\gtrsim\log\tl_{\ell+1}$.

\par\smallskip
\Setting a grid with $\Bg_1\ge2$, a consistent batched learner, $\nu\in\cP$, a sub-optimal arm $a$, and $\xi>0$.
\Goal show that $\liminf_\ell\E_\nu[\Ng_a(\tl_\ell)]/\log\tl_{\ell+1}\ge1/(\da+\xi)$; letting $\xi\to0$ then gives the lemma.
\par\smallskip

\begin{proof}
\proves{lem:onebatch}
\uses{def:kinf, def:model, def:grid, def:learner, lem:kinf, cite:honda2010asymptotically-primal, cite:garivier2019explore-dataprocessing}
\stp1{an alternative instance where $a$ is optimal}Fix $\epsilon\in(0,1-\m^\star)$ with $\epsilon/(1-\m^\star-\epsilon)\le\xi/2$, and let
$\mu':=\m^\star+\epsilon\in(\m^\star,1)$. By Lemma~\ref{lem:kinf}(v) at the level $\mu'$ with the shift $\epsilon\le\mu'$,
$\Ki(\nu_a,\mu')\le\Ki(\nu_a,\m^\star)+\epsilon/(1-\mu')\le\da+\xi/2<\infty$. By the primal formula
(Lemma~\ref{cite:honda2010asymptotically-primal} and the remark after it) at the level $\mu'$, there is a distribution $H$ on $[0,1]$ with
$\E_H[X]\ge\mu'$ and $\KL(\nu_a,H)\le\Ki(\nu_a,\mu')+\xi/2$. Hence
\[
\E_H[X]>\m^\star\qquad\text{and}\qquad\KL(\nu_a,H)\le\da+\xi .
\]
$H\ne\delta_1$: otherwise $\KL(\nu_a,H)<\infty$ would force $\nu_a=\delta_1$, whose mean is $1>\m_a$. So $\m':=\E_H[X]<1$, and the instance $\nu':=\nu^H$ belongs to
$\cP$, with $a$ as unique optimal arm.

\stp2{the divergence of the observations}Let
\[
O_\ell:=\big(U,\ A_1,R_1,\dots,A_{\tl_\ell},R_{\tl_\ell},\ A_{\tl_\ell+1},\dots,A_{\tl_{\ell+1}}\big)
\]
be the observations up to $\tl_\ell$, followed by the actions of batch $\ell+1$. In $O_\ell$:
\begin{itemize}
\item each action is a function of $U$, of the previous actions and of the rewards of the batches that end before it, all of which precede it in $O_\ell$; so the actions
have the same conditional laws under $\nu$ and $\nu'$;
\item each reward $R_t$, $t\le\tl_\ell$, has conditional law $\nu_{A_t}$ under $\nu$ and $\nu'_{A_t}$ under $\nu'$, and these differ only when $A_t=a$.
\end{itemize}
The chain rule \cite[Lemma~15.1]{LS} therefore gives
\[
\KL\big(\PP_\nu^{O_\ell},\PP_{\nu'}^{O_\ell}\big)\ =\ \E_\nu[\Ng_a(\tl_\ell)]\,\KL(\nu_a,H).
\]

\stp3{data processing}Let $Z:=\Ng_a(\tl_{\ell+1})/\tl_{\ell+1}\in[0,1]$; it is a function of $O_\ell$. Let $p_\ell:=\E_\nu[Z]$ and $q_\ell:=\E_{\nu'}[Z]$. By data processing and
the elementary inequality $\kl(p,q)\ge(1-p)\log\frac1{1-q}-\log2$,
\begin{align*}
\E_\nu[\Ng_a(\tl_\ell)]\,\KL(\nu_a,H)\ &\ge\ \kl(p_\ell,q_\ell)&&\why{data processing, Step~2}\\
&\ge\ (1-p_\ell)\log\frac1{1-q_\ell}-\log2 .
\end{align*}
By consistency under $\nu$, $p_\ell\to0$. The left side is finite (it is at most $\tl_\ell\KL(\nu_a,H)$), so $q_\ell<1$ as soon as $p_\ell<1$, and the right side is finite.

\stp4{consistency under $\nu'$}Fix $\alpha\in(0,1)$. Under $\nu'$ every arm $c\ne a$ is sub-optimal, so by consistency there is $C_\alpha<\infty$ with, for every $\ell$,
\[
1-q_\ell\ =\ \sum_{c\ne a}\frac{\E_{\nu'}[\Ng_c(\tl_{\ell+1})]}{\tl_{\ell+1}}\ \le\ C_\alpha\,\tl_{\ell+1}^{\alpha-1},
\qquad\text{that is}\qquad
\log\frac1{1-q_\ell}\ \ge\ (1-\alpha)\log\tl_{\ell+1}-\log C_\alpha .
\]

\stp5{conclusion}Steps~3 and~4, with $\KL(\nu_a,H)\le\da+\xi$, give
\[
\E_\nu[\Ng_a(\tl_\ell)]\,(\da+\xi)\ \ge\ (1-p_\ell)\big[(1-\alpha)\log\tl_{\ell+1}-\log C_\alpha\big]-\log2 .
\]
Divide by $\log\tl_{\ell+1}\to\infty$ and let $\ell\to\infty$ ($p_\ell\to0$): the $\liminf$ is at least $(1-\alpha)/(\da+\xi)$. Let $\alpha\to0$.
\end{proof}

%----------------------------------------------------------------------
\subsection{Proof of Lemma~\ref{lem:tilts}}\label{sec:pf-lem:tilts}
\noindent\emph{Idea.} (b) is a computation of densities. For (a), ``$\ge$'' is the easy direction of the duality (Gibbs' inequality against the tilt by
$\log(1-\lambda(x-\m^\star))$). For ``$\le$'', the optimality condition of the dual formula at $\lambda$ defines a reweighting $h\nu_a$ of $\nu_a$, completed by an atom at $1$,
of mean exactly $\m^\star$, whose divergence from $\nu_a$ is the value of the dual objective.

\begin{proof}
\proves{lem:tilts}
\uses{lem:kinf, def:tilt, def:model, def:kinf}
\stp{(b)}{densities}\emph{Means.} $\E_{\rhom}[X]=(1-s)\m^\star<\m^\star$, and $\E_{\rhop}[X]=(1-s)\m^\star+s$, which lies in $(\m^\star,1)$ because $\m^\star<1$.

\noindent\emph{Absolute continuity and the two inequalities.} Since $\rho\ge(1-s)G$ and $G\ge e^{-C}\nu_a$, we have $\nu_a\ll G\ll\rho$ and
\[
\frac{dG}{d\rho}\ \le\ \frac1{1-s},\qquad\frac{d\nu_a}{d\rho}\ =\ \frac{d\nu_a}{dG}\cdot\frac{dG}{d\rho}\ \le\ \frac1{1-s}\,\frac{d\nu_a}{dG}.
\]
Integrating the logarithms, against $G$ for the first and against $\nu_a$ for the second, gives $\KL(G,\rho)\le\log\frac1{1-s}$ and
$\KL(\nu_a,\rho)\le\KL(\nu_a,G)+\log\frac1{1-s}$.

\noindent\emph{The constant $D_G$.} All the measures involved ($\nu_a$, $G$, $\rhom$, $\rhop$) are of the form $g\nu_a+c_0\delta_0+c_1\delta_1$ with $c_0,c_1\ge0$ and $g$
bounded above and below by positive constants. So their mutual densities are bounded above and below by positive constants on the supports, and their logarithms take their
values in a bounded interval; $D_G$ is its length.

\stp{(a), $\ge$}{Gibbs' inequality}
\par\smallskip
\Setting a distribution $G$ on $[0,1]$ with $\E_G[X]\ge\m^\star$ and $\KL(\nu_a,G)<\infty$ (every boundary tilt is such a $G$).
\Goal $\KL(\nu_a,G)\ge\da$.
\par\smallskip
Let $\lambda\in\Lset{\m^\star}$ and $f(y):=\log(1-\lambda(y-\m^\star))$, a bounded function on $[0,1]$. Let $G_f:=e^fG/\E_G[e^f]$. Since
$\E_G[e^f]=1-\lambda(\E_G[X]-\m^\star)\le1$,
\begin{align*}
0\ \le\ \KL(\nu_a,G_f)\ &=\ \KL(\nu_a,G)-\E_{\nu_a}[f]+\log\E_G[e^f]&&\why{Gibbs' inequality; definition of $G_f$}\\
&\le\ \KL(\nu_a,G)-\E_{\nu_a}[f].
\end{align*}
So $\E_{\nu_a}[\log(1-\lambda(X-\m^\star))]\le\KL(\nu_a,G)$ for every $\lambda\in\Lset{\m^\star}$; the supremum over $\lambda$ is $\da$. Finally $\da>0$ by
Lemma~\ref{lem:kinf}(iv).

\stp{(a), $\le$}{a tilt from the dual formula}
\par\smallskip
\Goal for every $\gamma'>0$, a boundary tilt $G$ with $\KL(\nu_a,G)\le\da+\gamma'$.
\par\smallskip
\noindent\emph{Notation.} Let $\Lambda^\star:=\frac1{1-\m^\star}$. For $\lambda\in[0,\Lambda^\star)$ and $y\in[0,1]$, let
\[
\varphi(\lambda):=\E_{\nu_a}\big[\log(1-\lambda(X-\m^\star))\big],\qquad h_\lambda(y):=\frac1{1-\lambda(y-\m^\star)},\qquad\chi(\lambda):=\E_{\nu_a}\big[(X-\m^\star)h_\lambda(X)\big],
\]
so that $\da=\sup_\lambda\varphi(\lambda)$. Two identities follow from $h_\lambda=1+\lambda(y-\m^\star)h_\lambda$ and $y\,h_\lambda=\m^\star h_\lambda+(y-\m^\star)h_\lambda$:
\[
\E_{\nu_a}[h_\lambda]=1+\lambda\chi(\lambda),\qquad\E_{\nu_a}[Xh_\lambda]=\m^\star\big(1+\lambda\chi(\lambda)\big)+\chi(\lambda).
\]

\noindent\emph{Sub-step 1: $\chi$ is non-decreasing and continuous, with $\chi(0)<0$.} With $\log u\le u-1$, for $\lambda,\lambda'\in[0,\Lambda^\star)$,
\[
\varphi(\lambda')-\varphi(\lambda)\ =\ \E_{\nu_a}\Big[\log\frac{1-\lambda'(X-\m^\star)}{1-\lambda(X-\m^\star)}\Big]\ \le\ -(\lambda'-\lambda)\,\chi(\lambda).
\]
Adding this inequality to the one with $\lambda$ and $\lambda'$ exchanged gives $(\lambda'-\lambda)(\chi(\lambda')-\chi(\lambda))\ge0$: $\chi$ is non-decreasing. Moreover
$\chi(0)=\m_a-\m^\star<0$, and $\chi$ is Lipschitz on $[0,\lambda_1]$ for every $\lambda_1<\Lambda^\star$ (its integrand has a bounded derivative there).

\noindent\emph{Sub-step 2: a boundary tilt for every $\lambda$ with $\chi(\lambda)\le0$.} Fix such a $\lambda$ and let
\[
D_\lambda:=(1-\m^\star)+\big(-\chi(\lambda)\big)\big(1-\lambda(1-\m^\star)\big)\ \ge\ 1-\m^\star,\qquad r:=\frac{1-\m^\star}{D_\lambda}\in(0,1],
\]
\[
h:=r\,h_\lambda,\qquad c:=1-r\big(1+\lambda\chi(\lambda)\big),\qquad G:=h\,\nu_a+c\,\delta_1 .
\]
\begin{itemize}
\item $c\ge0$: this is $(1-\m^\star)(1+\lambda\chi)\le D_\lambda$, i.e.\ $(1-\m^\star)\lambda\chi\le(-\chi)(1-\lambda(1-\m^\star))$, true since the left side is $\le0$ and the
right side $\ge0$.
\item $G$ has mass $r(1+\lambda\chi)+c=1$ (first identity).
\item $G$ has mean $r\big[\m^\star(1+\lambda\chi)+\chi\big]+c=1-rD_\lambda=\m^\star$ (second identity).
\item $h$ is bounded above and below by positive constants, since $1-\lambda(1-\m^\star)\le1-\lambda(y-\m^\star)\le1+\lambda\m^\star$.
\end{itemize}
So $G$ is a boundary tilt. Since $G\ge h\nu_a$, $\frac{d\nu_a}{dG}\le\frac1h$ $\nu_a$-a.s., and
\begin{align*}
\KL(\nu_a,G)\ &\le\ \E_{\nu_a}[-\log h]&&\why{$d\nu_a/dG\le1/h$}\\
&=\ \varphi(\lambda)-\log r&&\why{$-\log h_\lambda=\log(1-\lambda(X-\m^\star))$}\\
&\le\ \varphi(\lambda)+\big(-\chi(\lambda)\big)(\Lambda^\star-\lambda)&&\why{$-\log r=\log\big(1+\tfrac{(-\chi)(1-\lambda(1-\m^\star))}{1-\m^\star}\big)$, $\log(1+u)\le u$}.
\end{align*}

\noindent\emph{Sub-step 3: choice of $\lambda$.} Let $\gamma'>0$.
\begin{itemize}
\item If $\chi(\lambda_1)>0$ for some $\lambda_1<\Lambda^\star$: by continuity there is $\lambda^\star\in(0,\lambda_1)$ with $\chi(\lambda^\star)=0$, and Sub-step~2 gives
$\KL(\nu_a,G)\le\varphi(\lambda^\star)\le\da$.
\item Otherwise $\chi\le0$ on $[0,\Lambda^\star)$, and $-\chi\le-\chi(0)=\m^\star-\m_a$ since $\chi$ is non-decreasing. Take $\lambda:=\max\big(0,\Lambda^\star-\gamma'/(\m^\star-\m_a+1)\big)$:
then $(-\chi(\lambda))(\Lambda^\star-\lambda)\le\gamma'$, and Sub-step~2 gives $\KL(\nu_a,G)\le\da+\gamma'$.
\end{itemize}
\end{proof}

%----------------------------------------------------------------------
\subsection{Proof of Lemma~\ref{lem:com}}\label{sec:pf-lem:com}
\noindent\emph{Idea.} Under $\QQ$ the rewards of $a$ after $\tl_{\ell-1}$ come from $G$, under $\PP_\rho$ from $\rho$. The likelihood ratio of the first $\Ng_a(\tl_\ell)$
rewards of $a$ is $e^{-\Sl^\rho}$, and the events of $\mathcal H_\ell$ only depend on these rewards and on $V$.

\begin{proof}
\proves{lem:com}
\uses{lem:stacks, lem:tilts, def:learner, def:grid, def:stacks, def:Q}
\stp1{reduction to fixed counts}Let $F\in\mathcal H_\ell$. For integers $n\le N$ let $F_{n,N}:=F\cap\{\nzero=n,\ \Ng_a(\tl_\ell)=N\}$. By Lemma~\ref{lem:stacks}(iii) there is a
Borel set $B_{n,N}$ such that, under $\QQ$ as well as under $\PP_\rho$,
\[
F_{n,N}=\big\{(V,Z_1,\dots,Z_N)\in B_{n,N}\big\}.
\]
It suffices to prove (i) for each $F_{n,N}$ and to sum over $n\le N$.

\stp2{the two laws of $(V,Z_1,\dots,Z_N)$}
\begin{itemize}
\item Under $\QQ$, on $\{\nzero=n\}$, $(V,Z_1,\dots,Z_N)=(V,\Xs_{a,1},\dots,\Xs_{a,n},Y_1,\dots,Y_{N-n})$, of law $P_V\otimes\nu_a^{\otimes n}\otimes G^{\otimes(N-n)}$.
\item Under $\PP_\rho$ it has the law $P_V\otimes\rho^{\otimes N}$.
\end{itemize}

\stp3{comparison}For every measurable $\psi\ge0$,
\[
\int\psi\,d\rho\ \ge\ \int\psi\Big(\frac{d\nu_a}{d\rho}\Big)^{-1}d\nu_a\qquad\text{and}\qquad\int\psi\,d\rho\ \ge\ \int\psi\Big(\frac{dG}{d\rho}\Big)^{-1}dG
\]
(Lemma~\ref{lem:tilts}(b): $\nu_a\ll\rho$, $G\ll\rho$). Applying these inequalities coordinate by coordinate (Tonelli) to $\psi=\1_{B_{n,N}}$ gives
$\PP_\rho(F_{n,N})\ge\E_\QQ\big[\1_{F_{n,N}}e^{-\Sl^\rho}\big]$, since on $F_{n,N}$ the product of the inverse densities is $e^{-\Sl^\rho}$. Summing over $n\le N$ gives (i).

\stp4{(ii)}The run up to $\tl_{\ell-1}$, actions of batch $\ell$ included, only uses the rewards of $a$ received before $\tl_{\ell-1}$, which are not tilted. So $\QQ$ and
$\PP_\nu$ coincide on $\mathcal H_{\ell-1}$.
\end{proof}

%----------------------------------------------------------------------
\subsection{Proof of Lemma~\ref{lem:typ}}\label{sec:pf-lem:typ}
\par\smallskip
\Goal bound by $\eta(b)/4$ the $\QQ$-probability of each of the four events (two sums, two values of $\rho$) whose union is $T_b^c$.
\par\smallskip
\begin{proof}
\proves{lem:typ}
\uses{lem:hoeffding, lem:tilts, def:tilt, def:stacks, def:Q}
\stp1{the sum $\Sl_1^\rho$}Under $\QQ$, $\Sl_1^\rho=\sum_{j\le\nzero}\xi_j$ with $\xi_j:=\log\frac{d\nu_a}{d\rho}(\Xs_{a,j})$. The $\xi_j$ are i.i.d., of mean $\KL(\nu_a,\rho)$ and of
range at most $D_G$. A union bound over the value $n$ of $\nzero$ (the value $n=0$ gives $\Sl_1^\rho=0\le b$) and Hoeffding's inequality give
\begin{align*}
\QQ\big(\Sl_1^\rho>\nzero(\KL(\nu_a,\rho)+\varepsilon_H)+b\big)\ &\le\ \sum_{n\ge1}e^{-2(n\varepsilon_H+b)^2/(nD_G^2)}&&\why{union bound, Lemma~\ref{lem:hoeffding}(a)}\\
&\le\ \sum_{n\ge1}e^{-2n\varepsilon_H^2/D_G^2}\,e^{-4\varepsilon_Hb/D_G^2}&&\why{$(n\varepsilon_H+b)^2\ge n^2\varepsilon_H^2+2n\varepsilon_Hb$}\\
&\le\ \frac{e^{-4\varepsilon_Hb/D_G^2}}{1-e^{-2\varepsilon_H^2/D_G^2}}&&\why{geometric series}.
\end{align*}

\stp2{the sum $\Sl_2^\rho$}Under $\QQ$, $\Sl_2^\rho=\sum_{i\le\Mb}\log\frac{dG}{d\rho}(Y_i)$, whose terms are i.i.d.\ with mean $\KL(G,\rho)$ and range at most $D_G$. The same
computation, with a union bound over the value of $\Mb$, gives the same bound.

\stp3{conclusion}Sum the four bounds.
\end{proof}

%----------------------------------------------------------------------
\subsection{Proof of Lemma~\ref{lem:twobatches}}\label{sec:pf-lem:twobatches}
\noindent\emph{Idea.} On the histories in which the rewards of $a$ in batch $\ell$ look like $G$, the learner cannot distinguish $\nu^{\rhom}$ ($a$ barely sub-optimal) from
$\nu^{\rhop}$ ($a$ barely optimal) unless $\Ng_a(\tl_{\ell-1})\da\gtrsim\log\tl_{\ell+1}$. Whatever it does with batch $\ell+1$, of size $\gtrsim\tl_{\ell+1}$, it then loses
polynomially under one of them.

\par\smallskip
\Setting\dt{s:lbproof} a constant $\cB>0$ and a rank $\ell_B$ with $\Bg_\ell\ge\cB\tl_\ell$ for every $\ell\ge\ell_B$ (they exist since $\liminf_\ell\Bg_\ell/\tl_\ell>0$).
For (i): $\varepsilon\in(0,1)$ and $C<\infty$ are fixed; we write
\[
\alpha:=\varepsilon/8\qquad\text{and}\qquad z:=\log\tl_{\ell+1}\quad(\text{which depends on }\ell).
\]
\Goal (i) $\displaystyle\limsup_{\ell}\PP_\nu\Big(\nzero\le(1-\varepsilon)^2\frac z{\da}\Big)\le\limsup_\ell\PP_\nu(\Mb>Cz)$, where $\nzero=\Ng_a(\tl_{\ell-1})$ and
$\Mb=\Ng_a(\tl_\ell)-\nzero$. Letting $C\to\infty$ gives (i), since $(1-\varepsilon)^2$ ranges over $(0,1)$. Then (ii).
\par\smallskip

\begin{proof}
\proves{lem:twobatches}
\uses{lem:tilts, lem:typ, lem:com, def:kinf, def:tilt, def:model, def:grid, def:learner, def:Q}
\stp1{the two instances}By Lemma~\ref{lem:tilts}(a) (infimum equal to $\da>0$), take a boundary tilt $G$ with
\[
\KL(\nu_a,G)\ \le\ \frac{\da}{1-\varepsilon}.
\]
Since $\frac{1-4\alpha}{1-\varepsilon}=\frac{1-\varepsilon/2}{1-\varepsilon}>1$ and $\KL(\nu_a,G)>0$, we can choose $s\in(0,1)$ and $\varepsilon_H>0$ so small that
\[
k_1:=\KL(\nu_a,G)+\log\frac1{1-s}+\varepsilon_H\ \le\ \frac{1-4\alpha}{1-\varepsilon}\,\KL(\nu_a,G)
\qquad\text{and}\qquad
C\,k_2\ \le\ \alpha,\quad k_2:=\log\frac1{1-s}+\varepsilon_H .
\]
By Lemma~\ref{lem:tilts}(b), $\nu^{\rhom}\in\cP$ with $a$ sub-optimal, and $\nu^{\rhop}\in\cP$ with $a$ as unique optimal arm. Both are consistent instances for the learner.

\stp2{a bound on the likelihood ratios}For $b>0$ let
\[
A_b:=\big\{\nzero k_1+\Mb k_2+2b\le(1-2\alpha)z\big\},
\]
an event of $\mathcal H_{\ell-1}$ ($\nzero$ and $\Mb$ are $\mathcal H_{\ell-1}$-measurable). Let $\rho\in\{\rhom,\rhop\}$. On $A_b\cap T_b$ (Lemma~\ref{lem:typ}),
\begin{align*}
\Sl^\rho\ &\le\ \nzero\big(\KL(\nu_a,\rho)+\varepsilon_H\big)+\Mb\big(\KL(G,\rho)+\varepsilon_H\big)+2b&&\why{definition of $T_b$}\\
&\le\ \nzero k_1+\Mb k_2+2b&&\why{Lemma~\ref{lem:tilts}(b)}\\
&\le\ (1-2\alpha)z&&\why{definition of $A_b$}.
\end{align*}

\stp3{whatever batch $\ell+1$ does, $A_b$ is unlikely}Let
\[
E:=\big\{a\text{ receives at most }\Bg_{\ell+1}/2\text{ pulls in batch }\ell+1\big\}\ \in\ \mathcal H_\ell .
\]
The events $A_b$, $T_b$ (which only involves the rewards of $a$ observed by $\tl_\ell$) and $E$ belong to $\mathcal H_\ell$. Assume $\ell+1\ge\ell_B$.

\noindent\emph{On $E^c$, $a$ gets the bulk of batch $\ell+1$}, so $\Ng_a(\tl_{\ell+1})>\Bg_{\ell+1}/2\ge\cB\tl_{\ell+1}/2$. Hence
\begin{align*}
\QQ(A_b\cap T_b\cap E^c)\ &\le\ e^{(1-2\alpha)z}\,\E_\QQ\big[\1_{A_b\cap T_b\cap E^c}\,e^{-\Sl^{\rhom}}\big]&&\why{Step~2: $\Sl^{\rhom}\le(1-2\alpha)z$ there}\\
&\le\ \tl_{\ell+1}^{1-2\alpha}\,\PP_{\rhom}\big(A_b\cap T_b\cap E^c\big)&&\why{Lemma~\ref{lem:com}(i)}\\
&\le\ \tl_{\ell+1}^{1-2\alpha}\,\PP_{\rhom}\big(\Ng_a(\tl_{\ell+1})\ge\cB\tl_{\ell+1}/2\big)&&\why{inclusion}\\
&\le\ \frac2{\cB}\,\tl_{\ell+1}^{-\alpha}\,\frac{\E_{\rhom}[\Ng_a(\tl_{\ell+1})]}{\tl_{\ell+1}^\alpha}&&\why{Markov}\\
&\xrightarrow[\ell\to\infty]{}\ 0&&\why{consistency at $\nu^{\rhom}$, where $a$ is sub-optimal}.
\end{align*}

\noindent\emph{On $E$, the other arms get the bulk of batch $\ell+1$}, so $\sum_{c\ne a}\Ng_c(\tl_{\ell+1})\ge\cB\tl_{\ell+1}/2$. The same chain with $\rhop$ gives
\[
\QQ(A_b\cap T_b\cap E)\ \le\ \frac2{\cB}\,\tl_{\ell+1}^{-\alpha}\sum_{c\ne a}\frac{\E_{\rhop}[\Ng_c(\tl_{\ell+1})]}{\tl_{\ell+1}^\alpha}\ \xrightarrow[\ell\to\infty]{}\ 0
\qquad\why{consistency at $\nu^{\rhop}$, where every $c\ne a$ is sub-optimal}.
\]

\noindent\emph{Back to $\PP_\nu$.} Since $A_b\in\mathcal H_{\ell-1}$,
\begin{align*}
\PP_\nu(A_b)\ &=\ \QQ(A_b)&&\why{Lemma~\ref{lem:com}(ii)}\\
&\le\ \QQ(T_b^c)+\QQ(A_b\cap T_b\cap E^c)+\QQ(A_b\cap T_b\cap E)\\
&\le\ \eta(b)+o(1)&&\why{Lemma~\ref{lem:typ}, and the two limits above}.
\end{align*}

\stp4{few pulls of $a$ before batch $\ell$ imply $A_b$}We check that, for $\ell$ large enough that $2b\le\alpha z$,
\[
\Big\{\nzero\le(1-\varepsilon)\frac z{\KL(\nu_a,G)}\Big\}\cap\big\{\Mb\le Cz\big\}\ \subset\ A_b .
\]
On the left event, the three terms of $A_b$ satisfy
\[
\nzero k_1\ \le\ (1-\varepsilon)\frac z{\KL(\nu_a,G)}\cdot\frac{1-4\alpha}{1-\varepsilon}\KL(\nu_a,G)=(1-4\alpha)z,\qquad\Mb k_2\le Czk_2\le\alpha z,\qquad2b\le\alpha z,
\]
whose sum is at most $(1-2\alpha)z$. Moreover $(1-\varepsilon)z/\KL(\nu_a,G)\ge(1-\varepsilon)^2z/\da$ by the choice of $G$ in Step~1. Hence
\begin{align*}
\PP_\nu\Big(\nzero\le(1-\varepsilon)^2\frac z{\da}\Big)\ &\le\ \PP_\nu\Big(\nzero\le(1-\varepsilon)\frac z{\KL(\nu_a,G)},\ \Mb\le Cz\Big)+\PP_\nu(\Mb>Cz)\\
&\le\ \PP_\nu(A_b)+\PP_\nu(\Mb>Cz)&&\why{inclusion above}\\
&\le\ \eta(b)+o(1)+\PP_\nu(\Mb>Cz)&&\why{Step~3}.
\end{align*}

\stp5{limits, proof of (i)}Take the $\limsup$ in $\ell$, then let $b\to\infty$ ($\eta(b)\to0$; $G$, $s$, $\varepsilon_H$ depend on $C$ but not on $b$): this is the Goal. The left
side depends neither on $b$ nor on $C$; letting $C\to\infty$ and using the assumption of (i) gives (i).

\stp6{proof of (ii)}Assume $\E_\nu[\Ng_a(\tl_\ell)]\le C_0\log\tl_{\ell+1}$ for every $\ell$ large. By Markov's inequality,
$\PP_\nu(\Mb>C\log\tl_{\ell+1})\le\PP_\nu(\Ng_a(\tl_\ell)>C\log\tl_{\ell+1})\le C_0/C$, so the assumption of (i) holds. Then, with $z=\log\tl_{\ell+1}$,
\begin{align*}
\E_\nu[\Ng_a(\tl_{\ell-1})]\ &\ge\ (1-\varepsilon)\frac z{\da}\,\PP_\nu\Big(\Ng_a(\tl_{\ell-1})>(1-\varepsilon)\frac z{\da}\Big)&&\why{Markov}\\
&=\ (1-\varepsilon)\big(1-o(1)\big)\frac z{\da}&&\why{(i)}.
\end{align*}
Divide by $z$, let $\ell\to\infty$, then $\varepsilon\to0$.
\end{proof}

%======================================================================
\section{Deterministic properties of HY}\label{sec:app-det}

\par\smallskip
\Setting every lemma of this section holds for \emph{every} run of \HY{} on \emph{every} sequence of rewards in $[0,1]$: no probability is involved. Recall that
$\Ng_c(\tl_\ell)\ge1$ for every arm $c$ and every $\ell\ge1$ (batch $1$ pulls every arm).
\Plan four properties, used by both upper bounds.
\begin{enumerate}
\item[(1)] Lemma~\ref{lem:imedround} (\emph{IMED rounds}): an arm pulled at an \term{s:act}{IMED round} has a small index, hence a bounded count.
\item[(2)] Lemma~\ref{lem:service} (\emph{queue lemma}): in a \term{s:act}{paced phase}, every paced demand is served, up to $K-1$ pulls, after an initial burst.
\item[(3)] Lemma~\ref{lem:W}: the \term{s:act}{IMED phase} of a long batch lifts the \term{s:water}{water level} above $\log\tl_{j-1}$.
\item[(4)] Lemma~\ref{lem:sparse} (\emph{sparse activations}): \term{s:act}{activated} batches are sparse on the logarithmic scale.
\end{enumerate}
\par\smallskip


\begin{lemma}[IMED rounds]\label{lem:imedround}
\uses{def:hy, def:learner, def:grid, def:hy-vocab}
Let $\ell\ge1$ and let $c\ne\Ld_\ell$ be pulled at an \term{s:act}{IMED round} $t$ of batch $\ell+1$. Then
\[
\Kt_\ell(c)<\infty\qquad\text{and}\qquad\Ng_c(t-1)\Kt_\ell(c)+\yg_\ell\log\Ng_c(t-1)\ <\ \yg_\ell\log t .
\]
Consequently, if $\Kt_\ell(c)\ge D>0$, then every round $t'$ of batch $\ell+1$ with $t'\le\sA_\ell$ satisfies
\[
\Ng_c(t')\ \le\ \max\Big(\Ng_c(\tl_\ell),\ 1+\frac{\yg_\ell\log t'}D\Big).
\]
\end{lemma}
\role{bounds the pulls of an arm in the IMED phase: Lemmas~\ref{lem:goodA}, \ref{lem:decompA}, \ref{lem:pullsC}.}
\noindent(see proof in Subsection~\ref{sec:pf-lem:imedround}.)

%----------------------------------------------------------------------
\subsection{The paced phase: a queue lemma}

\begin{definition}[non-leaders and backlog]\label{def:backlog}
\uses{def:hy, def:hy-vocab, def:learner}
For $\ell\in\Act$ (batch $\ell+1$ is activated), let $\mathcal C_\ell:=\cA\setminus\{\Ld_\ell\}$ be the non-leaders. For a round
$t\in[\sA_\ell,\tl_{\ell+1}]$, the \emph{backlog} of the paced demands after round $t$ is
\[
Q_\ell(t):=\sum_{c\in\mathcal C_\ell}\big(\nt_{\ell,c}(t)-\Ng_c(t)\big)^+ .
\]
\end{definition}

\begin{lemma}[service: the paced demands are served at every round]\label{lem:service}
\uses{def:hy-vocab, def:hy, def:grid, def:learner, def:backlog}
Let $\ell\in\Act$ and let $t\in(\sA_\ell,\tl_{\ell+1}]$ be a paced round.
\begin{enumerate}
\item[(a)] \emph{(No overshoot.)} Every $c\in\mathcal C_\ell$ satisfies $\Ng_c(t)\le\max\big(\Ng_c(\sA_\ell),\nt_{\ell,c}(t)\big)$, and is pulled at most $\Ug_{\ell+1}+K-1$ times in the
\term{s:act}{paced phase}.
\item[(b)] \emph{(Service.)} If $\sA_\ell\log\sA_\ell\ge2(K-1)\Ug_{\ell+1}$, then
\[
Q_\ell(t)\ \le\ \max\big((K-1)(\Ug_{\ell+1}+K-1)-(t-\sA_\ell),\ K-1\big).
\]
In particular, if $t-\sA_\ell\ge(K-1)(\Ug_{\ell+1}+K-1)$, then $\Ng_c(t)\ge\nt_{\ell,c}(t)-(K-1)$ for every $c\in\mathcal C_\ell$.
\item[(c)] \emph{(The \term{t:leader}{leader}.)} $\Ld_\ell$ is pulled at least $(t-\sA_\ell)-(K-1)(\Ug_{\ell+1}+K-1)$ times in $(\sA_\ell,t]$.
\end{enumerate}
\end{lemma}
\role{every paced demand is served: Lemmas~\ref{lem:structA} and~\ref{lem:pullsC}.}
\noindent(see proof in Subsection~\ref{sec:pf-lem:service}.)

%----------------------------------------------------------------------
\subsection{The IMED prelude lifts the water level}\label{sec:app-water-det}

\begin{definition}[the water level, (W1)--(W2) and the rank $\lW$]\label{def:water}\dt{s:water}
\uses{def:hy, def:hy-vocab, def:grid, def:learner}
\emph{The water level.} Fix a batch $j\ge2$, made of the rounds $(\tl_{j-1},\tl_j]$. For an arm $c$ and a round $t\in(\tl_{j-1},\tl_j+1]$, let
\[
I_t(c):=\frac{\Ng_c(t-1)\,\Kt_{j-1}(c)}{\yg_{j-1}}+\log\Ng_c(t-1),
\]
the IMED index of $c$ divided by $\yg_{j-1}$. The \emph{water level} of batch $j$ is the smallest index at the end of the batch,
\[
\Lam_j:=\min_cI_{\tl_j+1}(c)=\min_c\Big(\frac{\Ng_c(\tl_j)\,\Kt_{j-1}(c)}{\yg_{j-1}}+\log\Ng_c(\tl_j)\Big).
\]
It is computed with the IMED divergences of the batch, whatever rule pulled the arms. Two facts are used below.
\begin{itemize}
\item $\Kt_{j-1}\ge0$ and $\yg_{j-1}>0$ are fixed during the batch, so $t\mapsto I_t(c)$ is non-decreasing.
\item At an IMED round $t$, the pulled arm minimizes $\Jg_{j-1}(\cdot,t)=\yg_{j-1}I_t(\cdot)$, hence minimizes $I_t$.
\end{itemize}

\emph{Two properties.}
\begin{itemize}
\item[(W1)] Every arm satisfies $\dfrac{\Ng_c(\tl_j)\Kt_{j-1}(c)}{\yg_{j-1}}+\log\Ng_c(\tl_j)\ge\Lam_j$. In particular, since $\Kt_{j-1}(\Ld_{j-1})=0$,
\[
\log\Ng_{\Ld_{j-1}}(\tl_j)\ \ge\ \Lam_j .
\]
\item[(W2)] Every arm $c$ pulled in batch $j$ satisfies $\log\Ng_c(\tl_{j-1})\le\Lam_j$.
\end{itemize}
(W1) holds for every batch, since $\Lam_j$ is a minimum.

\emph{The rank $\lW$.} Let $\lW$ be the smallest $j\ge2$ such that every $j'\ge j$ satisfies the three conditions
\[
\tl_{j'-1}\ge e^e,\qquad\Ug_{j'}+K\le\tl_{j'-1},\qquad\tl_{j'-1}^{1+\thv_{j'-1}}-\tl_{j'-1}-1\ge K\tl_{j'-1} .
\]
It exists for the default $\thv$: $\tl_j\to\infty$, $\lceil\Lg^\beta\rceil+K\le e^{\Lg}$ for $\Lg$ large, and $\tl_j^{\thv_j}=e^{\Lg_j/\log\Lg_j}\to\infty$. It depends only on
$K$, $\beta$ and the grid.
\end{definition}

\begin{lemma}[the IMED phase lifts the water level]\label{lem:W}
\uses{def:hy, def:hy-vocab, def:grid, def:water}
With the default $\thv$, for every run of \HY{} and every batch $j\ge\lW$:
\begin{enumerate}
\item[(i)] (W2) holds in batch $j$;
\item[(ii)] if batch $j$ is \term{s:act}{activated} ($j-1\in\Act$), then $\ \Lam_j\ \ge\ \log\big((\sA_{j-1}-\tl_{j-1})/K\big)\ \ge\ \Lg_{j-1}$.
\end{enumerate}
\end{lemma}
\role{gives (W2) for \HY, on which the whole water-level argument of Section~\ref{sec:app-xone} rests (Lemma~\ref{lem:lowmass}).}
\noindent(see proof in Subsection~\ref{sec:pf-lem:W}.)

\begin{lemma}[activated batches are sparse on the logarithmic scale]\label{lem:sparse}
\uses{def:hy, def:hy-vocab, def:grid}
With the default $\thv$, for every integer $k\ge1$, the number of $\ell\in\Act$ with $\tl_\ell\ge e^e$ and $\Lg_\ell\in[2^k,2^{k+1})$ is at most $k+1$. Hence, for every
$f:[1,\infty)\to[0,\infty)$,
\[
\sum_{\ell\in\Act,\ \tl_\ell\ge e^e}f(\Lg_\ell)\ \le\ \sum_{k\ge1}(k+1)\sup_{y\in[2^k,2^{k+1})}f(y).
\]
\end{lemma}
\role{makes the nets of the consistency threshold summable, in Lemma~\ref{lem:amb}.}
\noindent(see proof in Subsection~\ref{sec:pf-lem:sparse}.)

%----------------------------------------------------------------------
\subsection{Proof of Lemma~\ref{lem:imedround}}\label{sec:pf-lem:imedround}
\begin{proof}
\proves{lem:imedround}
\uses{lem:kinf, def:hy, def:grid, def:learner}
\stp1{the leader has index below $\yg_\ell\log t$}$\Kt_\ell(\Ld_\ell)=0$: if $\tms_\ell<1$ this is Lemma~\ref{lem:kinf}(ii), since the mean of $\tnu_{\Ld_\ell}(\tl_\ell)$ is
$\tms_\ell$; if $\tms_\ell=1$ it is the convention. Since $\Ng_{\Ld_\ell}(t-1)\le t-1$,
\[
\Jg_\ell(\Ld_\ell,t)\ =\ \yg_\ell\log\Ng_{\Ld_\ell}(t-1)\ \le\ \yg_\ell\log(t-1)\ <\ \yg_\ell\log t .
\]

\stp2{the pulled arm}The pulled arm $c$ is a minimizer of $\Jg_\ell(\cdot,t)$, and the leader is not one (it would have been pulled). Hence
$\Jg_\ell(c,t)<\Jg_\ell(\Ld_\ell,t)<\yg_\ell\log t$; in particular $\Kt_\ell(c)<\infty$ (with $\infty\cdot n=\infty$ for $n\ge1$).

\stp3{the consequence}Let $\Kt_\ell(c)\ge D>0$ (so $c\ne\Ld_\ell$) and $t'\le\sA_\ell$ a round of batch $\ell+1$. If $c$ is not pulled in $(\tl_\ell,t']$, then
$\Ng_c(t')=\Ng_c(\tl_\ell)$. Otherwise let $t''$ be its last pull there; it is an IMED round, and $\Ng_c(t')=\Ng_c(t''-1)+1$. By Step~2 at $t''$, since $\log\Ng_c(t''-1)\ge0$,
\[
\Ng_c(t''-1)\,D\ \le\ \Jg_\ell(c,t'')\ <\ \yg_\ell\log t''\ \le\ \yg_\ell\log t',
\]
so $\Ng_c(t')<1+\yg_\ell\log t'/D$.
\end{proof}

%----------------------------------------------------------------------
\subsection{Proof of Lemma~\ref{lem:service}}\label{sec:pf-lem:service}
\noindent\emph{Idea.} The backlog is a queue served at rate one. The targets of the rejected arms are concave in $t$ and capped, so after an initial burst they bring less
than half a unit per round, plus one unit per arm for the ceilings.

\par\smallskip
\Setting $\ell\in\Act$. Write $n_c:=\nt_{\ell,c}$, $U:=\Ug_{\ell+1}$, $Q:=Q_\ell$ and $\sigma:=\sA_\ell$; note $\sigma\ge\tl_\ell\ge2$.
\Goal (a), (b), (c) for every paced round $t\in(\sigma,\tl_{\ell+1}]$.
\par\smallskip

\noindent\emph{Two facts.}
\begin{itemize}
\item[(F1)] For $c\in\mathcal C_\ell$, $\Kh_\ell(c)$, the verdict, $\om_\ell$ and $\yg_\ell$ are fixed during the batch. So $t\mapsto n_c(t)$ is non-decreasing on the integers,
and $0\le n_c\le U+K-1$.
\item[(F2)] At a paced round $t$, \HY{} pulls the first $c\in\mathcal C_\ell$ with $\Ng_c(t-1)<n_c(t)$, and $\Ld_\ell$ if there is none.
\end{itemize}

\begin{proof}
\proves{lem:service}
\uses{def:learner, def:grid, def:hy, def:backlog}
\stp{(a)}{no overshoot}If $c$ is not pulled in $(\sigma,t]$, then $\Ng_c(t)=\Ng_c(\sigma)$. Otherwise let $u$ be its last pull there. By (F2), $\Ng_c(u-1)<n_c(u)$, so
\[
\Ng_c(t)\ =\ \Ng_c(u-1)+1\ \le\ n_c(u)\ \le\ n_c(t)\qquad\why{(F1)}.
\]
Each pull of $c$ in the paced phase brings its count to a distinct value in $(\Ng_c(\sigma),U+K-1]$; so there are at most $U+K-1$ of them.

\stp{(b1)}{one round of the queue}For $t\in[\sigma,\tl_{\ell+1}]$ let $e_c(t):=(n_c(t)-\Ng_c(t))^+$, so that $Q(t)=\sum_ce_c(t)$. For $t>\sigma$, the \emph{arrivals} and
the \emph{net input} of round $t$ are
\[
\delta(t):=\sum_{c\in\mathcal C_\ell}\big(n_c(t)-n_c(t-1)\big)\ \ge\ 0,\qquad X_t:=\delta(t)-1 .
\]
\emph{Claim:} $Q(t)\le(Q(t-1)+X_t)^+$. Let $d_c:=(n_c(t)-\Ng_c(t-1))^+$, the demand of $c$ before the pull of round $t$. Then
$d_c\le e_c(t-1)+n_c(t)-n_c(t-1)$, so $\sum_cd_c\le Q(t-1)+\delta(t)$.
\begin{itemize}
\item If $\sum_cd_c=0$, the leader is pulled and $Q(t)=0$.
\item Otherwise the pulled arm $c^\ast$ has $e_{c^\ast}(t)=d_{c^\ast}-1$, and $e_c(t)=d_c$ for $c\ne c^\ast$; so $Q(t)=\sum_cd_c-1\le Q(t-1)+X_t$.
\end{itemize}

\stp{(b2)}{arrivals are slow}\emph{Claim:} for integers $\sigma\le u<t$,
\[
\sum_{c\in\mathcal C_\ell}\big(n_c(t)-n_c(u)\big)\ \le\ (K-1)+\frac{t-u}2 .
\]
If $c$ is ambiguous, or rejected with $\Kh_\ell(c)=\infty$, then $n_c$ is constant. Let $c$ be rejected with $\Kh_\ell(c)\in(0,\infty)$ (rejection requires
$\Kh_\ell(c)>0$), let $f:=(1+\om_\ell)\yg_\ell/\Kh_\ell(c)$, and assume $n_c(t)>n_c(u)$. Then the cap $U$ is not reached at $u$, i.e.\ $f\log u<U$, and
\begin{align*}
n_c(t)-n_c(u)\ &\le\ \lceil f\log t\rceil-\lceil f\log u\rceil&&\why{cap not reached at $u$}\\
&<\ f\log(t/u)+1\\
&\le\ \frac U{\log u}\cdot\frac{t-u}u+1&&\why{$f<U/\log u$, $\log(t/u)\le(t-u)/u$}\\
&\le\ \frac{U(t-u)}{\sigma\log\sigma}+1&&\why{$u\log u\ge\sigma\log\sigma$}\\
&\le\ \frac{t-u}{2(K-1)}+1&&\why{assumption of (b)}.
\end{align*}
Summing over the $K-1$ arms of $\mathcal C_\ell$ gives the claim.

\stp{(b3)}{unrolling the queue}By induction on $t$ with Step~(b1) (Lindley's recursion), for $t\in(\sigma,\tl_{\ell+1}]$,
\[
Q(t)\ \le\ \max\Big(Q(\sigma)+\sum_{r=\sigma+1}^tX_r,\ \max_{\sigma<u\le t}\sum_{r=u+1}^tX_r\Big)\qquad(\text{the sum for }u=t\text{ is }0).
\]

\stp{(b4)}{bounding the two terms}\emph{The second term.} For $\sigma<u\le t$, Step~(b2) gives
\[
\sum_{r=u+1}^tX_r\ =\ \sum_{c\in\mathcal C_\ell}\big(n_c(t)-n_c(u)\big)-(t-u)\ \le\ K-1-\frac{t-u}2\ \le\ K-1 .
\]
\emph{The first term.} It equals
\[
Q(\sigma)+\sum_{r=\sigma+1}^tX_r\ =\ \sum_{c\in\mathcal C_\ell}\Big[\big(n_c(\sigma)-\Ng_c(\sigma)\big)^++n_c(t)-n_c(\sigma)\Big]-(t-\sigma)\ \le\ (K-1)(U+K-1)-(t-\sigma),
\]
because each bracket equals $\max\big(n_c(t)-\Ng_c(\sigma),\,n_c(t)-n_c(\sigma)\big)\le n_c(t)\le U+K-1$. (The counts left by the IMED phase do not matter.)

\noindent\emph{Conclusion of (b).} Steps~(b3) and~(b4) give the bound on $Q(t)$. If $t-\sigma\ge(K-1)(U+K-1)$, then $Q(t)\le K-1$, so each $e_c(t)\le K-1$, that is
$\Ng_c(t)\ge n_c(t)-(K-1)$.

\stp{(c)}{the leader}By (a), the arms of $\mathcal C_\ell$ are pulled at most $(K-1)(U+K-1)$ times in the paced phase. Every other paced round of $(\sigma,t]$ is a pull of
$\Ld_\ell$.
\end{proof}

%----------------------------------------------------------------------
\subsection{Proof of Lemma~\ref{lem:W}}\label{sec:pf-lem:W}
\noindent\emph{Idea.} During an IMED round the pulled arm minimizes the indices, and the indices only grow. So the arm with most IMED pulls certifies, at its last IMED
pull, that all the indices are at least the logarithm of its count. After $\approx\tl^{1+\thv}$ IMED rounds this is above $\log\tl_{j-1}$, hence above the logarithm of every
count at the start of the batch, which gives (W2) at the paced rounds too.

\par\smallskip
\Setting a batch $j\ge\lW$; write $\ell:=j-1$, so that batch $j$ is batch $\ell+1$.
\Goal (i) (W2) in batch $j$; (ii) if $\ell\in\Act$, $\Lam_j\ge\log((\sA_\ell-\tl_\ell)/K)\ge\Lg_\ell$.
\par\smallskip

\begin{proof}
\proves{lem:W}
\uses{def:learner, def:grid, def:hy, def:hy-vocab, def:water}
\stp1{(W2) at the IMED rounds}Let $c$ be pulled at an IMED round $t$ of batch $j$. For every arm $c'$,
\begin{align*}
\log\Ng_c(\tl_\ell)\ &\le\ \log\Ng_c(t-1)&&\why{counts are non-decreasing}\\
&\le\ I_t(c)&&\why{$\Kt_\ell\ge0$}\\
&\le\ I_t(c')&&\why{$c$ minimizes $I_t$}\\
&\le\ I_{\tl_j+1}(c')&&\why{$I_\cdot(c')$ is non-decreasing}.
\end{align*}
Taking the minimum over $c'$ gives $\log\Ng_c(\tl_\ell)\le\Lam_j$. If batch $j$ is not activated, all its rounds are IMED rounds, and (i) is proved.

\stp2{(ii): the water level of an activated batch}Let $\ell\in\Act$. The $R:=\sA_\ell-\tl_\ell$ rounds $\tl_\ell+1,\dots,\sA_\ell$ are IMED rounds of batch $j$. Some arm
$c^\ast$ is pulled at least $R/K$ times among them; let $t^\ast$ be its last IMED pull. Then
\[
\Ng_{c^\ast}(t^\ast-1)\ \ge\ \Ng_{c^\ast}(\tl_\ell)+\lceil R/K\rceil-1\ \ge\ R/K ,
\]
and, since $c^\ast$ minimizes $I_{t^\ast}$, for every arm $c'$,
\[
I_{\tl_j+1}(c')\ \ge\ I_{t^\ast}(c')\ \ge\ I_{t^\ast}(c^\ast)\ \ge\ \log\Ng_{c^\ast}(t^\ast-1)\ \ge\ \log(R/K).
\]
So $\Lam_j\ge\log(R/K)$. By the third condition defining $\lW$, $R\ge\tl_\ell^{1+\thv_\ell}-1-\tl_\ell\ge K\tl_\ell$, so $\log(R/K)\ge\Lg_\ell$.

\stp3{(W2) at the paced rounds}Let $\ell\in\Act$ and $t$ a paced round of batch $j$.
\begin{itemize}
\item If $c\ne\Ld_\ell$ is pulled at $t$, then $\Ng_c(\tl_\ell)\le\Ng_c(t-1)<\nt_{\ell,c}(t)\le\Ug_j+K-1$, so
$\log\Ng_c(\tl_\ell)\le\log(\Ug_j+K)\le\Lg_\ell\le\Lam_j$ (second condition defining $\lW$, then Step~2).
\item If $\Ld_\ell$ is pulled at $t$, then $\log\Ng_{\Ld_\ell}(\tl_\ell)\le\log\tl_\ell=\Lg_\ell\le\Lam_j$ (Step~2).
\end{itemize}
With Step~1, (W2) holds at every round of batch $j$.
\end{proof}

%----------------------------------------------------------------------
\subsection{Proof of Lemma~\ref{lem:sparse}}\label{sec:pf-lem:sparse}
\begin{proof}
\proves{lem:sparse}
\uses{def:grid, def:hy-vocab, def:hy}
\stp1{an activated batch makes $\Lg$ jump}If $\ell\in\Act$, then $\tl_{\ell+1}\ge\sA_\ell+1>\tauA_\ell=\tl_\ell^{1+\thv_\ell}$, so $\Lg_{\ell+1}>(1+\thv_\ell)\Lg_\ell$.

\stp2{counting the jumps in a dyadic block}Fix $k\ge1$, let $\ell_1<\dots<\ell_m$ be the indices of the statement, and $u:=1/((k+1)\log2)\in(0,1)$.
\begin{itemize}
\item \emph{Each jump is by a factor at least $1+u$.} For $i<m$, $\Lg_{\ell_i}<2^{k+1}$, so $\thv_{\ell_i}=1/\log\Lg_{\ell_i}>u$, and
\[
\Lg_{\ell_{i+1}}\ \ge\ \Lg_{\ell_i+1}\ >\ (1+u)\,\Lg_{\ell_i}\qquad\why{Step~1}.
\]
\item \emph{The block has ratio $2$.} So $2^{k+1}>\Lg_{\ell_m}>(1+u)^{m-1}2^k$, that is $(m-1)\log(1+u)<\log2$.
\end{itemize}
As $u<1$, $\log(1+u)\ge u/2$, so
\[
m-1\ <\ 2(k+1)(\log2)^2\ <\ k+1 .
\]

\stp3{the sum}Group the indices according to the $k$ with $\Lg_\ell\in[2^k,2^{k+1})$; since $\tl_\ell\ge e^e$, $\Lg_\ell\ge e>2$, so $k\ge1$.
\end{proof}

%======================================================================
\section{Upper bound when \texorpdfstring{$x>1$}{x>1}}\label{sec:app-xsup}

\par\smallskip
\Setting (in the whole section) $\xg_\ell\to x\in(1,\infty)$; $\thv_\ell\to0$ and $\gam_\ell\to0$ are deterministic; $\wz\ge0$, $\beta\ge2$; the margins $(\om_\ell)$
are a non-increasing sequence in $(0,1]$, fixed in each statement. The default parameters satisfy these conditions.
\Goal Proposition~\ref{prop:dyA}: with the dyadic margins $\ody$, $\ \limsup_{t\to\infty}\E[\Ng_a(t)]/\log t\le x^2/\da$.
\Plan the section goes from the top down.
\begin{enumerate}
\item[(1)] Constants, ranks and levels (Subsection~\ref{sec:app-constA}).
\item[(2)] \emph{The skeleton} (Subsection~\ref{sec:app-skelA}): the proposition; the events; the decomposition of $\Ng_a(t)$ over the batches into a plan, \term{t:net}{nets} and
lost batches; the probabilities and sums that make the \term{t:net}{nets} and the lost batches summable.
\item[(3)] \emph{Why is the threshold deterministic?} (Subsection~\ref{sec:app-structA}): from some rank on, every batch is \term{s:act}{activated} and its \term{s:act}{paced phase} is long.
\item[(4)] \emph{Why is the \term{t:leader}{leader} optimal?} (Subsection~\ref{sec:app-leaderA}).
\item[(5)] \emph{What happens in a \term{s:evA}{good batch}?} (Subsection~\ref{sec:app-goodA}): $a$ is \term{t:verdict}{rejected} and follows its optimal plan.
\end{enumerate}
\par\smallskip

\noindent\emph{Map.} A line ``$B\leftarrow A_1;\ A_2$'' means that $A_1$ and $A_2$ are used in the proof of $B$; the tools of Section~\ref{sec:app-tools} are not drawn, and Lemmas~\ref{lem:constA}, \ref{lem:ranksA} (used by almost every lemma) are only
linked to the proposition.
\begin{itemize}
\item Prop.~\ref{prop:dyA}, dyadic margins $\leftarrow$ Lemma~\ref{lem:decompA}, decomposition; Lemmas~\ref{lem:probasA}, \ref{lem:sumsA}, bad events, sums; Lemmas~\ref{lem:constA}, \ref{lem:ranksA}, \ref{lem:H}, constants, ranks, levels.
\item Lemma~\ref{lem:decompA}, decomposition $\leftarrow$ Lemma~\ref{lem:goodA}, good batch; Lemma~\ref{lem:leaderA}, the leader; Lemmas~\ref{lem:service}, \ref{lem:imedround}, queue, IMED rounds; Lemma~\ref{lem:structA}, structure of a batch.
\item Lemma~\ref{lem:goodA}, good batch $\leftarrow$ Lemma~\ref{lem:countsA}, counts of $a$; Lemma~\ref{lem:leaderA}, the leader; Lemma~\ref{lem:thrA}, threshold; Lemmas~\ref{lem:service}, \ref{lem:imedround}, queue, IMED rounds; Lemma~\ref{lem:structA}, structure of a batch.
\item Lemma~\ref{lem:countsA}, counts of $a$ $\leftarrow$ Lemma~\ref{lem:leaderA}, the leader; Lemma~\ref{lem:structA}, structure of a batch.
\item Lemma~\ref{lem:leaderA}, the leader $\leftarrow$ Lemma~\ref{lem:thrA}, threshold; Lemma~\ref{lem:structA}, structure of a batch.
\item Lemma~\ref{lem:thrA}, threshold $\leftarrow$ Lemma~\ref{lem:structA}, structure of a batch.
\item Lemma~\ref{lem:structA}, structure of a batch $\leftarrow$ Lemmas~\ref{lem:service}, \ref{lem:imedround}, queue, IMED rounds.
\end{itemize}


%----------------------------------------------------------------------
\subsection{Constants, ranks and levels}\label{sec:app-constA}

\begin{definition}[constants, $x>1$]\label{def:constA}\dt{s:constA}
\uses{def:model, def:kinf, def:hy, def:grid}
\emph{Fixed constants.} They depend on the instance, on $a$ and on $x$, not on the accuracy $\varepsilon$. For every sub-optimal arm $c$, Lemma~\ref{lem:theta} applied to
$\nu_c$ gives $\theta_c>0$, $c_{\theta,c}>0$ and $\varepsilon_{\theta,c}>0$. The minima below are over the sub-optimal arms.
\begin{align*}
\eb&:=\min\Big(\frac\Dg5,\ \frac{1-\m^\star}2,\ \min_c\varepsilon_{\theta,c}\Big)&&\text{(largest accuracy)},\\
\eg&:=\min\Big(1,\ \frac{\Dg-\eb}{2(\mtwo+\eb)},\ \min_c\frac{c_{\theta,c}}{2\kap(\m^\star+\eb)}\Big)&&\text{(gain in the deviation inequalities)},\\
\Ceta&:=\big(4+2\kap(\m^\star+\eb)\big)e^{2(1+2\eg)}&&\text{(constant of these inequalities)},\\
\cst&:=\frac{2-\m^\star}{1-\m^\star},\qquad\xp:=2x,\qquad\hb_\ell:=(1-\gam_\ell)\,\yg_{\ell-1}\Lg_\ell&&\text{(deterministic threshold)},
\end{align*}
and, for a margin $\om\in(0,1]$, $\ew:=\om\da/(16\cst)$ (accuracy allowed by the margin $\om$).

\emph{Constants depending on $\varepsilon\in(0,\eb]$.}
\begin{align*}
\kape_\varepsilon&:=\kap(\m^\star+\varepsilon),\qquad\De_\varepsilon:=\Ki(\nu_a,\m^\star-2\varepsilon)-\varepsilon,\qquad\Kp_\varepsilon:=\Ki(\nu_a,\m^\star+\varepsilon)+2\varepsilon,\\
\Rt_\varepsilon&:=\kap(\m^\star+\varepsilon)-\log\big(1-e^{-\varepsilon/4}\big),\qquad\rt_\varepsilon:=\frac{\varepsilon^2}{2\Rt_\varepsilon^2},\qquad
\ngr_\varepsilon:=\Big\lceil\frac8{\varepsilon(1-e^{-\varepsilon/4})(1-\m^\star-\varepsilon)}\Big\rceil+1,\\
\cAc(\varepsilon)&:=\frac{2K}{1-e^{-2\varepsilon^2}},\qquad\CLU(\varepsilon):=\frac{1+\ngr_\varepsilon}{1-e^{-\rt_\varepsilon}} .
\end{align*}
$\De_\varepsilon$ and $\Kp_\varepsilon$ are a lower and an upper estimate of $\da$; the others are constants of deviation inequalities.
\end{definition}

\begin{lemma}[constants]\label{lem:constA}
\uses{def:constA, def:model, def:kinf, def:stacks, def:hy}
Let $\varepsilon\in(0,\eb]$.
\begin{enumerate}
\item[(a)] $\m^\star+\varepsilon<1$; every sub-optimal $c$ has $\m_c+\varepsilon<\m^\star-3\varepsilon$; and $\mtwo+\varepsilon<\m^\star-2\varepsilon$.
\item[(b)] (Gain, optimal arm.) For $n\ge1$ and $z\in\mathbb R$,
\[
\PP\big(n\Ki(\hnu_{a^\star,n},\mtwo+\varepsilon)\ge z\big)\ \le\ \Ceta\,n^2e^{-(1+2\eg)z}.
\]
\item[(c)] (Gain, pairs.) For a sub-optimal $c$, integers $1\le n\le n'$ and $z\in\mathbb R$,
\[
\PP\big(n\Ki(\hnu_{c,n},\m^\star+\varepsilon)\ge z,\ \hm_{c,n'}\ge\m^\star-3\varepsilon\big)\ \le\ \Ceta\,n^2e^{-(1+2\eg)z}.
\]
\item[(d)] If moreover $\varepsilon\le\ew$ for a margin $\om\in(0,1]$:
\[
\De_\varepsilon\ge\da-2\cst\varepsilon\ge\tfrac78\da,\qquad\frac1{\De_\varepsilon}\le\frac1{\da}+\frac{3\cst}{\da^2}\varepsilon,\qquad\Kp_\varepsilon\le\da+2\cst\varepsilon,\qquad(1+\om)\De_\varepsilon\ge\Kp_\varepsilon .
\]
\end{enumerate}
\end{lemma}
\noindent(see proof in Subsection~\ref{sec:pf-lem:constA}.)

\begin{definition}[ranks when $x>1$]\label{def:ranksA}\dt{s:ranksA}
\uses{def:constA, def:hy, def:hy-vocab, def:grid}
The rank $\lA$ is the smallest $\ell\ge5$ such that every $\ell'\ge\ell-3$ satisfies:
\begin{itemize}
\item[(R1)] $\Lg_{\ell'-1}\ge1$, $\ \tl_{\ell'-1}\ge e^e$ and $\xg_{\ell'}\le\xp$;
\item[(R2)] (long paced phase) $\ell'-1\in\Act$, $\ \tl_{\ell'}-\sA_{\ell'-1}\ge K(\Ug_{\ell'}+K)$, $\ \sA_{\ell'-1}\log\sA_{\ell'-1}\ge2(K-1)\Ug_{\ell'}$ and $\tl_{\ell'}\ge K\Ug_{\ell'}$;
\item[(R3)] (the threshold beats two batches ahead) $(1+2\eg)\big(\hb_{\ell'}-\wz\kap(\m^\star+\eb)\big)\ge(1+\eg)\Lg_{\ell'+2}$.
\end{itemize}
For $\varepsilon\in(0,\eb]$, the \emph{count threshold} is $\ug_\ell(\varepsilon):=\Lg_{\ell-2}/\kape_\varepsilon$. For $\varepsilon,\varepsilon'\in(0,\eb]$ with $\De_\varepsilon>0$, the rank
$\lone(\varepsilon,\varepsilon')$ is the smallest $\ell\ge\lA+3$ such that every $\ell'\ge\ell$ satisfies:
\begin{itemize}
\item[(T1)] $\kape_\varepsilon\Lg_{\ell'-2}^{\beta-1}\ge1$;
\item[(T2)] $\De_\varepsilon\Lg_{\ell'-1}^{\beta-1}\ge2\xp^3$;
\item[(T3)] $2\wz\le\min(\varepsilon,\varepsilon')\,\Ug_{\ell'-2}\ $ and $\ \wz\kap(\m^\star+\eb)\le\varepsilon\Lg_{\ell'-2}/\kap(\m^\star+\eb)$;
\item[(T4)] $K+1+8\xp^2\Lg_{\ell'+1}/\varepsilon'^2\le\Ug_{\ell'+1}$.
\end{itemize}
Two remarks on (T3), used below.
\begin{itemize}
\item Since $\kape_\varepsilon\le\kap(\m^\star+\eb)$, (T3) gives $\wz\kap(\m^\star+\varepsilon)\le\varepsilon\,\ug_{\ell'}(\varepsilon)$.
\item Both parts of (T3) are non-decreasing in the accuracies: (T3) at some accuracies implies (T3) at every larger accuracies.
\end{itemize}
\end{definition}

\begin{lemma}[ranks]\label{lem:ranksA}
\uses{def:grid, def:constA, def:ranksA}
$\lA$ exists, and $\lone(\varepsilon,\varepsilon')$ exists. For $k\ge\lA-4$ and $i\ge0$, $\Lg_{k+i}\le\xp^i\Lg_k$.
\end{lemma}
\noindent(see proof in Subsection~\ref{sec:pf-lem:ranksA}.)

\begin{definition}[levels]\label{def:H}\dt{s:H}
\uses{def:hy, def:grid}
The paced plan of batch $k$ aims at the level $\yg_{k-1}\Lg_k$ at its end. The largest level reached so far is, for $j\ge1$ and for a round
$t\in(\tl_j,\tl_{j+1}]$,
\[
\HH_j:=\max_{1\le k\le j}\yg_{k-1}\Lg_k,\qquad\HH(t):=\max\big(\HH_j,\ \yg_j\log t\big),
\]
and $\HH(t):=0$ for $t\le\tl_1$. So $\HH(\tl_j)=\HH_j$ for $j\ge2$, and $\HH(\tl_1)=0$.
\end{definition}

\begin{lemma}[the levels]\label{lem:H}
\uses{def:ass-x, def:grid, def:H}
Under Assumption~\ref{def:ass-x} (any $x\ge1$), $\HH_j/\Lg_j\to x^2$ as $j\to\infty$, and $\HH(t)/\log t\to x^2$ as $t\to\infty$.
\end{lemma}
\noindent(see proof in Subsection~\ref{sec:pf-lem:H}.)

%----------------------------------------------------------------------
\subsection{The skeleton: proposition, events, decomposition and sums}\label{sec:app-skelA}

\begin{proposition}[upper bound with the dyadic margins, $x>1$]\label{prop:dyA}
\uses{def:hy, def:learner, def:kinf}
\HY$(\wz,\beta,\ody,\thv,\gam)$ with $\thv_\ell\to0$ and $\gam_\ell\to0$ satisfies
\[
\limsup_{t\to\infty}\frac{\E[\Ng_a(t)]}{\log t}\ \le\ \frac{x^2}{\da};\qquad\text{in particular }\E[\Ng_a(t)]=O(\log t).
\]
\end{proposition}
\noindent(see proof in Subsection~\ref{sec:pf-prop:dyA}.)

\noindent\emph{Outline of the proof} (written in full in Subsection~\ref{sec:pf-prop:dyA}).
\begin{itemize}
\item Lemma~\ref{lem:decompA} (\emph{decomposition}) bounds $\Ng_a(t)$ by the largest plan $\approx(1+\ody)\HH(t)/\da$, plus the largest net lost on a bad batch, plus the
batches lost on the rare events $\Fg$.
\item The lost batches are summable (Lemma~\ref{lem:sumsA}(b)).
\item The nets are bounded by a layer-cake argument over blocks of constant accuracy, with Lemma~\ref{lem:sumsA}(a) and the explicit rates of
Lemma~\ref{lem:probasA}.
\item Since $\ody\to0$ and $\HH(t)/\log t\to x^2$ (Lemma~\ref{lem:H}), the $\limsup$ is at most $x^2/\da$.
\end{itemize}

\begin{definition}[events when $x>1$]\label{def:evA}\dt{s:evA}
\uses{def:constA, def:hy, def:kinf, def:model, def:stacks, def:ranksA}
Let $\varepsilon\in(0,\eb]$ and $\ell\ge3$; $n,n',N$ denote integers, and $z_\ell:=\hb_{\ell-1}-\wz\kap(\m^\star+\eb)$.
\begin{itemize}
\item \emph{Accurate means of the \term{t:ref}{well-sampled} arms:} $\ \Ag_\ell:=\big\{\forall c,\ \forall n\ge\tfrac12\Ug_\ell:\ |\hm_{c,n}-\m_c|\le\varepsilon\big\}$.
\item \emph{An optimal arm wrongly \term{t:verdict}{rejected} against a sub-optimal level:} $\ \Vg_\ell:=\big\{\exists n\in[1,\tfrac12\Ug_\ell):\ n\Ki(\hnu_{a^\star,n},\mtwo+\varepsilon)\ge z_\ell\big\}$.
\item \emph{A rejected sub-optimal arm whose mean climbs:}
\[
\Pig_\ell:=\bigcup_{c:\m_c<\m^\star}\big\{\exists\,1\le n\le n'<\tfrac12\Ug_\ell:\ n\Ki(\hnu_{c,n},\m^\star+\varepsilon)\ge z_\ell,\ \hm_{c,n'}\ge\m^\star-3\varepsilon\big\}.
\]
\item \emph{Low and high empirical divergence of $a$:}
\begin{gather*}
\Lo_\ell:=\big\{\exists N\ge\ug_\ell(\varepsilon):\ \Ki(\hnu_{a,N},\m^\star-2\varepsilon)<\De_\varepsilon\big\},\\
\Up_\ell:=\big\{\exists n\ge\ug_\ell(\varepsilon):\ \Ki(\hnu_{a,n},\m^\star+\varepsilon)>\Ki(\nu_a,\m^\star+\varepsilon)+\varepsilon\big\}.
\end{gather*}
\item \emph{Good batch} ($\ell\ge5$) and \emph{failure}:
\[
\Gg_\ell:=\Ag_{\ell-2}\cap\Ag_{\ell-1}\cap\Ag_\ell\cap\Vg_{\ell-1}^c\cap\Vg_\ell^c\cap\Pig_{\ell-1}^c\cap\Pig_\ell^c\cap\Lo_\ell^c\cap\Up_\ell^c,\qquad
\Fg_\ell:=\Ag_{\ell-1}^c\cup\Ag_\ell^c\cup\Vg_\ell\cup\Pig_\ell .
\]
\end{itemize}
We write $\Ag_\ell(\varepsilon)$, \dots, when the accuracy must be specified. These are events on the stacks (the thresholds are deterministic). The events $\Ag_\ell^c$, $\Lo_\ell$,
$\Up_\ell$ are non-increasing in $\ell$.
\end{definition}

\begin{lemma}[decomposition, $x>1$]\label{lem:decompA}\dt{s:decA}
\uses{def:constA, def:hy, def:grid, def:learner, def:evA, def:ranksA}
Let $\varepsilon,\varepsilon'\in(0,\eb]$. For a pair $(\varepsilon,\ell)$ and a round $t$, the \emph{plan} of $a$ is $\ \mt_\ell(\varepsilon,t):=K+(1+\om_\ell)\yg_\ell\log t/\De_\varepsilon$.
\begin{enumerate}
\item[(a)] (One batch.) Let $(\varepsilon,\ell)$ satisfy the assumptions of Lemma~\ref{lem:goodA}, with $\ell\ge\lone(\varepsilon,\varepsilon')$. For every outcome and every
$t\in(\tl_\ell,\tl_{\ell+1}]$,
\begin{equation}\label{eq:DA}
\Ng_a(t)\ \le\ \max\Big(\underbrace{\Ng_a(\tl_\ell)}_{\text{past}},\ \underbrace{\mt_\ell(\varepsilon,t)}_{\text{plan}},\ \underbrace{2\Ug_{\ell+1}\1\big(\Gg_\ell(\varepsilon)^c\big)}_{\text{net}}\Big)
+\underbrace{(t-\tl_\ell)\,\1\big(\Fg_\ell(\varepsilon')\big)}_{\text{whole batch}}.
\end{equation}
\item[(b)] (Iteration.) Let $i_0\le j$ and accuracies $\varepsilon_i\in(0,\eb]$, $i_0\le i\le j$, such that each $(\varepsilon_i,i)$ satisfies the assumptions of (a). Let
$\bar M\ge\mt_i(\varepsilon_i,\tl_{i+1})$ for $i_0\le i<j$, and $\Yg_j:=\max_{i_0\le i\le j}2\Ug_{i+1}\1(\Gg_i(\varepsilon_i)^c)$. Then for every $t\in(\tl_j,\tl_{j+1}]$,
\[
\Ng_a(t)\ \le\ \tl_{i_0}+\max\big(\bar M,\mt_j(\varepsilon_j,t)\big)+\Yg_j+\sum_{i=i_0}^j\Bg_{i+1}\1\big(\Fg_i(\varepsilon')\big).
\]
\end{enumerate}
\end{lemma}
\role{the skeleton of the proof of Proposition~\ref{prop:dyA}. In a \term{s:evA}{good batch}, $a$ follows its optimal plan (Lemma~\ref{lem:goodA}); otherwise,
if the leader is still optimal, $a$ receives at most a \term{t:net}{net}; otherwise the whole batch is charged to the rare event $\Fg_\ell$.}
\noindent(see proof in Subsection~\ref{sec:pf-lem:decompA}.)

\begin{lemma}[bad events, $x>1$]\label{lem:probasA}
\uses{def:constA, def:evA, def:hy, def:grid, def:ranksA}
Let $\varepsilon\in(0,\eb]$.
\begin{enumerate}
\item[(a)] For $\ell\ge3$, $\ \PP(\Ag_\ell^c)\le\cAc(\varepsilon)\,e^{-\varepsilon^2\Ug_\ell}$.
\item[(b)] For $\ell\ge\lA$, $\ \PP(\Vg_\ell)\le\Ceta\,\Ug_\ell^3\,\tl_{\ell+1}^{-(1+\eg)}\ $ and $\ \PP(\Pig_\ell)\le K\Ceta\,\Ug_\ell^4\,\tl_{\ell+1}^{-(1+\eg)}$.
\item[(c)] For $\ell\ge3$, $\ \PP(\Lo_\ell)+\PP(\Up_\ell)\le\CLU(\varepsilon)\,e^{-\rt_\varepsilon\ug_\ell(\varepsilon)}$.
\end{enumerate}
\end{lemma}
\role{bounds the probabilities of the events that make a batch bad (the nets) or lost (the whole batch).}
\noindent(see proof in Subsection~\ref{sec:pf-lem:probasA}.)

\begin{lemma}[the two summable costs]\label{lem:sumsA}\dt{s:SE}
\uses{def:constA, def:evA, def:hy, def:grid, def:ranksA}
Let $\varepsilon\in(0,\eb]$ and $\SE_\varepsilon:=\sup_{y\ge0}\big(2\xp^3y-\varepsilon^2y^\beta\big)<\infty$. Then (in (a) the events may even be taken at an accuracy
$\varepsilon_k\in(0,\eb]$ depending on $k$, since the bounds of Lemma~\ref{lem:probasA}(b) do not depend on the accuracy):
\begin{enumerate}
\item[(a)] (Nets lost on $\Vg$, $\Pig$.) $\displaystyle\sum_{k\ge\lA}\Ug_{k+2}\big(\PP(\Vg_k)+\PP(\Pig_k)\big)\ \le\ (1+K)\Ceta\,\sig_{5\beta}(\eg)$;
\item[(b)] (Whole batches.) $\displaystyle\sum_{i\ge\lA}\Bg_{i+1}\PP\big(\Fg_i(\varepsilon)\big)\ \le\ \cAc(\varepsilon)\,e^{\SE_\varepsilon}+(1+K)\Ceta\,\sig_{4\beta}(\eg)$.
\end{enumerate}
\end{lemma}
\role{shows that the nets lost on $\Vg,\Pig$ and the whole batches lost on $\Fg$ cost a finite amount in total (Steps~3 and~5 of the proof of
Proposition~\ref{prop:dyA}).}
\noindent(see proof in Subsection~\ref{sec:pf-lem:sumsA}.)

\noindent The rest of the section proves the lemmas used by the decomposition, grouped by the question that they answer.

%----------------------------------------------------------------------
\subsection{Why is the threshold deterministic?}\label{sec:app-structA}

\begin{lemma}[structure of a batch when $x>1$]\label{lem:structA}
\uses{def:hy, def:grid, def:learner, def:kinf, def:hy-vocab, def:ranksA}
Let $\ell\ge\lA-4$, so that (R2) holds at $\ell+1$. For every run:
\begin{enumerate}
\item[(S1)] (Service.) For $c\ne\Ld_\ell$ and $t\in(\sA_\ell,\tl_{\ell+1}]$, $\ \Ng_c(t)\le\max\big(\Ng_c(\sA_\ell),\nt_{\ell,c}(t)\big)$; and $\Ng_c(\tl_{\ell+1})\ge\nt_{\ell,c}(\tl_{\ell+1})-(K-1)$.
\item[(S2)] (The \term{t:leader}{leader} is \term{t:ref}{well sampled}.) $\Ld_\ell$ is pulled at least $\Ug_{\ell+1}$ times in the \term{s:act}{paced phase}, so $\Ld_\ell\in\Wg_{\ell+1}$.
\item[(S3)] (A poorly sampled arm was \term{t:verdict}{rejected}.) If $c\notin\Wg_{\ell+1}$, then $c\ne\Ld_\ell$, $c$ is rejected at $\tl_\ell$, and
\[
\Ng_c(\tl_{\ell+1})\,\Kh_\ell(c)\ \ge\ (1+\om_\ell)\,\yg_\ell\Lg_{\ell+1}.
\]
In particular $\rr_\ell\in\Wg_{\ell+1}$.
\item[(S4)] (Every arm is explored.) If $\mh_\ell<1$, every arm satisfies $\Ng_c(\tl_{\ell+1})\ge\min\big(\Ug_{\ell+1},\ \yg_\ell\Lg_{\ell+1}/\kap(\mh_\ell)\big)$.
\end{enumerate}
\end{lemma}
\role{used in Lemmas~\ref{lem:thrA}, \ref{lem:leaderA}, \ref{lem:countsA}, \ref{lem:goodA}, \ref{lem:decompA}: from rank $\lA$ on, every batch is
\term{s:act}{activated} with a long paced phase, so every paced demand is served.}
\noindent(see proof in Subsection~\ref{sec:pf-lem:structA}.)

\begin{lemma}[the threshold when $x>1$]\label{lem:thrA}
\uses{def:hy, def:grid, def:constA, def:ranksA}
For every $\ell\ge\lA-3$ and every run, $\ \Lh_\ell=\Lg_\ell\ $ and $\ \hg_\ell=\hb_\ell=(1-\gam_\ell)\,\yg_{\ell-1}\Lg_\ell$: the threshold is deterministic.
\end{lemma}
\role{used in Lemmas~\ref{lem:leaderA}, \ref{lem:goodA}: the test at $\tl_\ell$ compares the evidence with a deterministic level.}
\noindent(see proof in Subsection~\ref{sec:pf-lem:thrA}.)

%----------------------------------------------------------------------
\subsection{Why is the leader optimal?}\label{sec:app-leaderA}

\begin{lemma}[the leader when $x>1$]\label{lem:leaderA}
\uses{def:constA, def:hy, def:evA, def:model, def:stacks, def:grid, def:ranksA}
Let $\varepsilon\in(0,\eb]$ and $\ell\ge\lA$ with $2\wz\le\varepsilon\Ug_\ell$. On $\Ag_{\ell-1}\cap\Ag_\ell\cap\Vg_\ell^c\cap\Pig_\ell^c$:
\begin{enumerate}
\item[(i)] some optimal arm lies in $\Wg_\ell$, the \term{t:ref}{reference} $\rr_\ell$ is optimal, and $\m^\star-2\varepsilon\le\mh_\ell\le\m^\star+\varepsilon$;
\item[(ii)] every sub-optimal $c$ has $\hm_c(\tl_\ell)<\m^\star-3\varepsilon$, hence $\Kt_\ell(c)\ge\Kh_\ell(c)\ge\varepsilon^2/8$;
\item[(iii)] the \term{t:leader}{leader} $\Ld_\ell$ is optimal.
\end{enumerate}
Moreover, on $\Ag_{\ell-1}$ alone, $\mh_{\ell-1}\le\m^\star+\varepsilon$.
\end{lemma}
\role{used in Lemmas~\ref{lem:countsA}, \ref{lem:goodA}, \ref{lem:decompA}: outside the rare events $\Vg_\ell,\Pig_\ell$ and on accurate means, the
\term{t:leader}{leader} is optimal and the \term{t:ref}{test level} is accurate.}
\noindent(see proof in Subsection~\ref{sec:pf-lem:leaderA}.)

\begin{lemma}[counts of $a$]\label{lem:countsA}
\uses{def:constA, def:evA, def:learner, def:grid, def:ranksA}
Let $\varepsilon\in(0,\eb]$ with $\De_\varepsilon>0$ (so that $\lone$ is defined) and $\ell\ge\lone(\varepsilon,\varepsilon')$ for some $\varepsilon'\in(0,\eb]$. On $\Ag_{\ell-1}$, $\Ng_a(\tl_\ell)\ge\ug_\ell(\varepsilon)$; on $\Ag_{\ell-2}$, $\Ng_a(\tl_{\ell-1})\ge\ug_\ell(\varepsilon)$.
\end{lemma}
\role{used in Lemma~\ref{lem:goodA}: the counts of $a$ are above the \term{s:ranksA}{count threshold}, where its empirical divergence is accurate.}
\noindent(see proof in Subsection~\ref{sec:pf-lem:countsA}.)

%----------------------------------------------------------------------
\subsection{What happens in a good batch?}\label{sec:app-goodA}

\begin{lemma}[good batch]\label{lem:goodA}
\uses{def:constA, def:hy, def:evA, def:grid, def:learner, def:ranksA}
Let $\varepsilon\in(0,\eb]$ with $\De_\varepsilon>0$, let $\ell\ge\lone(\varepsilon,\varepsilon')$ for some $\varepsilon'\in(0,\eb]$, and assume $(1+\om_{\ell-1})\De_\varepsilon\ge\Kp_\varepsilon$. On $\Gg_\ell(\varepsilon)$:
\begin{itemize}
\item $\Ld_{\ell-1}$ and $\Ld_\ell$ are optimal;
\item $a$ is \term{t:verdict}{rejected} at $\tl_\ell$, and $\Kh_\ell(a)\ge\De_\varepsilon$;
\item for every round $t\in(\tl_\ell,\tl_{\ell+1}]$,
\[
\Ng_a(t)\ \le\ \max\Big(\Ng_a(\tl_\ell),\ K+\frac{(1+\om_\ell)\,\yg_\ell\log t}{\De_\varepsilon}\Big).
\]
\end{itemize}
\end{lemma}
\role{used in Lemma~\ref{lem:decompA}: in a good batch, $a$ is \term{t:verdict}{rejected} and pulled at most according to its plan $(1+\om)\yg\log t/\De$.}
\noindent(see proof in Subsection~\ref{sec:pf-lem:goodA}.)

%======================================================================
\subsection{Proof of Lemma~\ref{lem:constA}}\label{sec:pf-lem:constA}
\begin{proof}
\proves{lem:constA}
\uses{lem:single, lem:stacks, lem:pair, lem:theta, lem:kinf, def:model, def:constA, def:kinf, def:stacks, def:hy, def:empkinf}
\stp{(a)}{means}$\varepsilon\le(1-\m^\star)/2$ gives $\m^\star+\varepsilon<1$. For a sub-optimal $c$, $\m_c\le\mtwo=\m^\star-\Dg\le\m^\star-5\varepsilon$, so
$\m_c+\varepsilon\le\m^\star-4\varepsilon<\m^\star-3\varepsilon$; the same with $\mtwo$ gives $\mtwo+\varepsilon<\m^\star-2\varepsilon$.

\stp{(b)}{Lemma~\ref{lem:single}}Apply Lemma~\ref{lem:single} to the rewards of $a^\star$ (i.i.d.\ by Lemma~\ref{lem:stacks}), with $m=\m^\star$, $r:=\mtwo+\varepsilon\in[0,\m^\star)$ and
$\eta=\eg$. Its condition holds:
\[
2\eg r\ \le\ 2\eg(\mtwo+\eb)\ \le\ \Dg-\eb\ \le\ \Dg-\varepsilon\ =\ \m^\star-r\qquad\why{definition of $\eg$}.
\]
Its constant $(4+2\kap(r))e^{2(1+2\eg)}$ is at most $\Ceta$, since $\kap$ is non-decreasing and $r\le\m^\star+\eb$.

\stp{(c)}{Lemma~\ref{lem:pair}}Apply Lemma~\ref{lem:pair} with $\nu=\nu_c$, $r=\m^\star+\varepsilon$, $q=\m^\star-3\varepsilon$, $\theta=\theta_c$, $\eta=\eg$. Its conditions hold:
\begin{itemize}
\item $2\eg\kap(r)\le2\eg\kap(\m^\star+\eb)\le c_{\theta,c}$ \ (definition of $\eg$);
\item $e^{2\eg\kap(r)}\Phis_{\nu_c}(\theta_c,q)\le e^{c_{\theta,c}}\Phis_{\nu_c}(\theta_c,q)\le1$ \ (Lemma~\ref{lem:theta}, as $\varepsilon\le\varepsilon_{\theta,c}$);
\item $e^{2\eg\kap(r)}\Psis_{\nu_c}(\theta_c,q)\le e^{c_{\theta,c}}\Psis_{\nu_c}(\theta_c,q)\le1-r$ \ (same).
\end{itemize}
The event $\{\hm_{c,n'}\ge q\}$ is $\{S_{n'}\ge n'q\}$, and the constant is at most $\Ceta$ as in (b).

\stp{(d)}{the estimates of $\da$}\emph{Lower estimate.} Lemma~\ref{lem:kinf}(v) with $\mu=\m^\star$ and the shift $2\varepsilon\le\m^\star$ (as $\varepsilon\le\Dg/5\le\m^\star/5$) gives
\[
\De_\varepsilon\ \ge\ \da-\frac{2\varepsilon}{1-\m^\star}-\varepsilon\ \ge\ \da-2\cst\varepsilon\qquad\why{$2\cst=2+\tfrac2{1-\m^\star}$}.
\]
\emph{Upper estimate.} Lemma~\ref{lem:kinf}(v) with $\mu=\m^\star+\varepsilon$ and the shift $\varepsilon$ gives
\[
\Ki(\nu_a,\m^\star+\varepsilon)\ \le\ \da+\frac{\varepsilon}{1-\m^\star-\varepsilon}\ \le\ \da+\frac{2\varepsilon}{1-\m^\star}\qquad\why{$1-\m^\star-\varepsilon\ge\tfrac{1-\m^\star}2$},
\]
so $\Kp_\varepsilon\le\da+\frac{2\varepsilon}{1-\m^\star}+2\varepsilon=\da+2\cst\varepsilon$.

\noindent\emph{Consequences of $\varepsilon\le\ew\le\da/(16\cst)$.} First $\De_\varepsilon\ge\da-\da/8=\frac78\da$. Then
\[
\frac1{\De_\varepsilon}-\frac1{\da}\ =\ \frac{\da-\De_\varepsilon}{\De_\varepsilon\da}\ \le\ \frac{2\cst\varepsilon}{\frac78\da^2}\ \le\ \frac{3\cst\varepsilon}{\da^2}.
\]
Finally, with $\varepsilon\le\om\da/(16\cst)$,
\[
(1+\om)\De_\varepsilon-\Kp_\varepsilon\ \ge\ (1+\om)(\da-2\cst\varepsilon)-(\da+2\cst\varepsilon)\ =\ \om\da-2\cst\varepsilon(2+\om)\ \ge\ \om\da-6\cst\varepsilon\ \ge\ \om\da\Big(1-\frac6{16}\Big)\ >\ 0 .
\]
\end{proof}

%----------------------------------------------------------------------
\subsection{Proof of Lemma~\ref{lem:ranksA}}\label{sec:pf-lem:ranksA}
\begin{proof}
\proves{lem:ranksA}
\uses{def:grid, def:constA, def:hy, def:kinf, def:model, def:ranksA}
\stp1{(R1)}$\Lg_\ell\to\infty$, and $\xg_\ell\to x<\xp$.

\stp2{(R2)}Since $\thv_\ell\to0$ and $\xg_\ell\to x>1$, there is $c_0>0$ with $(1+\thv_{\ell'-1})/\xg_{\ell'}\le1-c_0$ for $\ell'$ large. Then
\[
\sA_{\ell'-1}\ \le\ \tl_{\ell'-1}^{1+\thv_{\ell'-1}}\ =\ \tl_{\ell'}^{(1+\thv_{\ell'-1})/\xg_{\ell'}}\ \le\ \tl_{\ell'}^{1-c_0}\qquad\why{$\tl_{\ell'-1}=\tl_{\ell'}^{1/\xg_{\ell'}}$}.
\]
Recall $\Ug_{\ell'}\le\Lg_{\ell'-1}^\beta+1$. For $\ell'$ large, each condition of (R2) holds:
\begin{itemize}
\item activation: $\tl_{\ell'}>\tl_{\ell'}^{1-c_0}\ge\sA_{\ell'-1}$;
\item long paced phase: $\tl_{\ell'}-\sA_{\ell'-1}\ge\tl_{\ell'}/2\ge K(\Lg_{\ell'-1}^\beta+1+K)\ge K(\Ug_{\ell'}+K)$;
\item condition of the queue lemma: $\sA_{\ell'-1}\log\sA_{\ell'-1}\ge\tl_{\ell'-1}\Lg_{\ell'-1}\ge2(K-1)(\Lg_{\ell'-1}^\beta+1)$;
\item $\tl_{\ell'}\ge K(\Lg_{\ell'-1}^\beta+1)\ge K\Ug_{\ell'}$.
\end{itemize}
In each case a power or an exponential of $\Lg_{\ell'-1}$ dominates a polynomial in $\Lg_{\ell'-1}$.

\stp3{(R3)}$\displaystyle\frac{\hb_{\ell'}}{\Lg_{\ell'+2}}=\frac{(1-\gam_{\ell'})\xg_{\ell'-1}^2\Lg_{\ell'}}{\Lg_{\ell'+2}}=\frac{(1-\gam_{\ell'})\xg_{\ell'-1}^2}{\xg_{\ell'+1}\xg_{\ell'+2}}\to1$, and
$\Lg\to\infty$; so the ratio of the left side of (R3) to $\Lg_{\ell'+2}$ tends to $1+2\eg>1+\eg$.

\stp4{(T1)--(T4)}$\Lg\to\infty$, $\beta\ge2$, $\De_\varepsilon>0$, $\kape_\varepsilon\ge\kap(\m^\star)>0$ (as $\m^\star>\m_a\ge0$); for (T4), moreover $\Lg_{\ell'+1}\le\xp\Lg_{\ell'}$ and
$\Ug_{\ell'+1}\ge\Lg_{\ell'}^\beta$.

\stp5{last claim}$\Lg_{k+i}=\xg_{k+1}\cdots\xg_{k+i}\Lg_k$, and $\xg_{k'}\le\xp$ for $k'\ge\lA-3$ by (R1).
\end{proof}

%----------------------------------------------------------------------
\subsection{Proof of Lemma~\ref{lem:H}}\label{sec:pf-lem:H}
\begin{proof}
\proves{lem:H}
\uses{def:hy, def:grid, def:H}
\stp1{lower bound at the ends}$\HH_j\ge\yg_{j-1}\Lg_j=\xg_{j-1}^2\Lg_j$ for $j\ge2$, so $\liminf_j\HH_j/\Lg_j\ge x^2$.

\stp2{upper bound at the ends}Let $\delta'>0$ and $k_{\delta'}$ with $\xg_k\le x(1+\delta')$ for $k\ge k_{\delta'}$. For $k\le j$:
\begin{itemize}
\item if $k-1\ge k_{\delta'}$, then $\yg_{k-1}\Lg_k\le x^2(1+\delta')^2\Lg_k\le x^2(1+\delta')^2\Lg_j$;
\item otherwise $\yg_{k-1}\Lg_k\le C_{\delta'}:=\max_{k'\le k_{\delta'}}\yg_{k'-1}\Lg_{k'}$.
\end{itemize}
So $\HH_j\le\max(C_{\delta'},x^2(1+\delta')^2\Lg_j)$, and $\limsup_j\HH_j/\Lg_j\le x^2(1+\delta')^2$; let $\delta'\to0$.

\stp3{all rounds}For $t\in(\tl_j,\tl_{j+1}]$, $\HH(t)/\log t\ge\yg_j=\xg_j^2$, and, as $\log t>\Lg_j$, $\HH(t)/\log t\le\max(\HH_j/\Lg_j,\yg_j)$. Both bounds tend to $x^2$.
\end{proof}

%----------------------------------------------------------------------
\subsection{Proof of Proposition~\ref{prop:dyA}}\label{sec:pf-prop:dyA}
\noindent\emph{Idea.} The accuracy $\varepsilon_\ell$ decreases with the margin $\ody_\ell$, slowly enough ($\varepsilon_\ell\gtrsim\Lg_\ell^{-1/4}$) for the nets to stay summable,
and it changes only at the few indices where $\ody$ halves. The decomposition of Lemma~\ref{lem:decompA}(b) then has a main term $\approx(1+\ody)\HH(t)/\da$, and remainders
bounded independently of $t$.

\par\smallskip
\Setting \HY$(\wz,\beta,\ody,\thv,\gam)$ with $\thv_\ell\to0$, $\gam_\ell\to0$; a number $\theta'>x^2/\da$.
\Goal $\E[\Ng_a(t)]\le K_0+\varphi\,\HH(t)/\da$ for every $t$ large, with $K_0$ independent of $t$ and $\varphi\,x^2/\da<\theta'$. Since $\HH(t)/\log t\to x^2$
(Lemma~\ref{lem:H}), this gives $\limsup_t\E[\Ng_a(t)]/\log t<\theta'$, for every $\theta'>x^2/\da$.
\Plan Step~1 sets the accuracies; Step~2 a rank from which Lemma~\ref{lem:decompA} applies; Step~3 the decomposition; Steps~4--6 bound the \term{t:net}{nets}; Step~7 concludes.
\par\smallskip

\begin{proof}
\proves{prop:dyA}
\uses{lem:constA, lem:ranksA, lem:decompA, lem:goodA, lem:sumsA, lem:sums, lem:probasA, lem:H, def:hy, def:grid, def:kinf, def:constA, def:model, def:learner, def:evA, def:H, def:ranksA}\dt{s:dy}
\stp0{the dyadic margins}Recall $\ody_\ell=2^{-k_\ell}$ with $k_\ell=\lfloor\frac14\log_2\max(\Lg_\ell,1)\rfloor$. Since $k_\ell$ is non-decreasing:
\begin{itemize}
\item $\ody_\ell\in(0,1]$ is non-increasing, and $\ody_\ell\to0$;
\item $\ody_\ell\ge\max(\Lg_\ell,1)^{-1/4}$;
\item $\ody_\ell=2^{-k}$ with $k\ge1$ implies $\Lg_\ell\ge16^k$.
\end{itemize}

\stp1{the accuracies}Let $\cone:=\da/(16\cst)$, and
\[
\epl_\ell:=\min(\eb,\ \cone\ody_\ell),\qquad\etwo:=\min(\eb,\cone).
\]
Then $\epl_\ell\le\ew$ with $\om=\ody_\ell$; $(\epl_\ell)$ is non-increasing; and $\epl_\ell\ge\etwo\ody_\ell\ge\etwo\Lg_\ell^{-1/4}$ when $\Lg_\ell\ge1$. For every $\ell\ge2$,
Lemma~\ref{lem:constA}(d) with $\om=\ody_\ell$ gives
\[
\De_{\epl_\ell}\ge\tfrac78\da,\qquad\frac1{\De_{\epl_\ell}}\le\frac1{\da}+\frac{3\cst\cone\ody_\ell}{\da^2}=\frac{1+\frac3{16}\ody_\ell}{\da},\qquad
(1+\ody_{\ell-1})\De_{\epl_\ell}\ge(1+\ody_\ell)\De_{\epl_\ell}\ge\Kp_{\epl_\ell} .
\]

\stp2{a uniform rank}Let $\lonep$ be the smallest $\ell\ge\lA+3$ such that every $\ell'\ge\ell$ satisfies
\begin{itemize}
\item $\kap(\m^\star)\Lg_{\ell'-2}^{\beta-1}\ge1$ and $\frac78\da\Lg_{\ell'-1}^{\beta-1}\ge2\xp^3$;
\item (T3) at $\ell'$ with $\varepsilon=\varepsilon'=\epl_{\ell'}$;
\item (T4) at $\ell'$ with $\varepsilon'=\eb$.
\end{itemize}
\emph{It exists}, because $\epl_{\ell'}\Ug_{\ell'-2}\ge\etwo\Lg_{\ell'}^{-1/4}(\Lg_{\ell'}/\xp^3)^\beta\to\infty$ and $\epl_{\ell'}\Lg_{\ell'-2}\ge\etwo\Lg_{\ell'}^{3/4}/\xp^2\to\infty$
(Lemma~\ref{lem:ranksA}).

\noindent\emph{Claim: for $i\ge\lonep$, $i\ge\lone(\epl_i,\eb)$, and $(\epl_i,i)$ satisfies the assumptions of Lemma~\ref{lem:decompA}(a) with $\varepsilon'=\eb$.} For every
$\ell'\ge i$:
\begin{itemize}
\item (T1) and (T2) hold at the accuracy $\epl_i$, since $\kape_{\epl_i}\ge\kap(\m^\star)$ and $\De_{\epl_i}\ge\frac78\da$;
\item (T3) holds with the accuracies $(\epl_i,\eb)$, because it holds with $(\epl_{\ell'},\epl_{\ell'})$, $\epl_{\ell'}\le\epl_i\le\eb$, and (T3) is non-decreasing in the
accuracies;
\item (T4) holds with $\varepsilon'=\eb$.
\end{itemize}
The other assumptions of Lemma~\ref{lem:goodA} are given by Step~1.

\stp3{the decomposition}\emph{A factor $\varphi$ between $1$ and $\theta'\da/x^2$.} Let
\[
q:=\frac{\theta'\da}{x^2}>1,\qquad\xi:=\min\Big(1,\frac{q-1}3\Big),\qquad\varphi:=(1+\xi)\Big(1+\frac\xi4\Big).
\]
Then $\varphi=1+\frac54\xi+\frac14\xi^2\le1+\frac32\xi<q$ (as $\xi\le1$ and $\frac32\xi\le\frac{q-1}2$). Let $i_\xi\ge\lonep$ with $\ody_i\le\xi$ for every $i\ge i_\xi$.

\noindent\emph{The plans.} By Step~1, $\frac1{\De_{\epl_i}}\le\frac{1+\ody_i/4}{\da}$ (as $\frac3{16}\le\frac14$), so
\[
\mt_i(\epl_i,t)\ \le\ K+(1+\ody_i)\Big(1+\frac{\ody_i}4\Big)\frac{\yg_i\log t}{\da}.
\]
The factor $(1+\ody_i)(1+\frac{\ody_i}4)$ is at most $2\cdot\frac54=\frac52$ (as $\ody_i\le1$), and at most $\varphi$ for $i\ge i_\xi$.

\noindent\emph{The decomposition.} Let $j\ge i_\xi$, $t\in(\tl_j,\tl_{j+1}]$, and
\[
\bar M:=K+\frac52\frac{\HH_{i_\xi}}{\da}+\varphi\frac{\HH_j}{\da}.
\]
\emph{$\bar M$ bounds the past plans.} For $\lonep\le i<j$, $\yg_i\Lg_{i+1}\le\HH_{i+1}$; if $i<i_\xi$, the factor is at most $\frac52$ and $\HH_{i+1}\le\HH_{i_\xi}$; if $i\ge i_\xi$,
it is at most $\varphi$ and $\HH_{i+1}\le\HH_j$. So $\mt_i(\epl_i,\tl_{i+1})\le\bar M$.

\noindent\emph{The current plan.} $\max(\bar M,\mt_j(\epl_j,t))\le K+\frac52\HH_{i_\xi}/\da+\varphi\HH(t)/\da$, since $\HH(t)=\max(\HH_j,\yg_j\log t)$.

\noindent Lemma~\ref{lem:decompA}(b) with $i_0=\lonep$, $\varepsilon_i=\epl_i$ and $\varepsilon'=\eb$, taken in expectation, gives
\begin{equation}\label{eq:main2A}
\E[\Ng_a(t)]\ \le\ \tl_{\lonep}+K+\frac52\frac{\HH_{i_\xi}}{\da}+\varphi\frac{\HH(t)}{\da}+2\E[Y'_j]+\sum_{i=\lonep}^j\Bg_{i+1}\PP\big(\Fg_i(\eb)\big),
\qquad Y'_j:=\max_{\lonep\le i\le j}\Ug_{i+1}\1\big(\Gg_i(\epl_i)^c\big).
\end{equation}
By Lemma~\ref{lem:sumsA}(b) at the accuracy $\eb$, the last sum is at most $\cAc(\eb)e^{\SE_{\eb}}+(1+K)\Ceta\sig_{4\beta}(\eg)$. It remains to bound $\E[Y'_j]$ independently of $j$.

\stp4{blocks of constant accuracy}The sequence $(\epl_i)_{\lonep\le i\le j}$ is non-increasing. Let $I_k=[\alpha_k,\alpha'_k]$, $k=1,\dots,r$, be the maximal intervals of
$[\lonep,j]$ on which it is constant, with value $e_k$; so $\alpha_1=\lonep$. Then
\[
Y'_j\ \le\ \sum_kY^{(k)},\qquad Y^{(k)}:=\max_{i\in I_k}\Ug_{i+1}\1\big(\Gg_i(e_k)^c\big).
\]
Fix $k$ and let $v_i:=\Ug_{i+1}$, $v^{(k)}_{\alpha_k-1}:=0$, $v^{(k)}_{i-1}:=v_{i-1}$ for $i>\alpha_k$. By the layer cake (Lemma~\ref{lem:sums}(v)), then the inclusion
\[
\bigcup_{i'=i}^{\alpha'_k}\Gg_{i'}^c\ \subset\ \Ag_{i-2}^c\cup\Lo_i\cup\Up_i\cup\bigcup_{k'=i-1}^{\alpha'_k}\big(\Vg_{k'}\cup\Pig_{k'}\big)
\]
(each $\Gg_{i'}^c$ is a union of nine events, and $\Ag^c$, $\Lo$, $\Up$ are non-increasing in the index), all at the accuracy $e_k$,
\[
\E[Y^{(k)}]\ \le\ \sum_{i\in I_k}\big(v_i-v^{(k)}_{i-1}\big)\Big[\underbrace{\PP(\Ag_{i-2}^c)+\PP(\Lo_i)+\PP(\Up_i)}_{\text{Step 6}}
+\underbrace{\sum_{k'=i-1}^{\alpha'_k}\big(\PP(\Vg_{k'})+\PP(\Pig_{k'})\big)}_{\text{Step 5}}\Big].
\]

\stp5{the $\Vg$ and $\Pig$ terms}Exchanging the sums and telescoping in $i$, the $\Vg$, $\Pig$ part of block $k$ equals
$\sum_{k'=\alpha_k-1}^{\alpha'_k}(\PP(\Vg_{k'})+\PP(\Pig_{k'}))\,v_{\min(\alpha'_k,k'+1)}$, which is at most $\sum_{k'=\alpha_k-1}^{\alpha'_k}\Ug_{k'+2}(\PP(\Vg_{k'})+\PP(\Pig_{k'}))$. The bounds
of Lemma~\ref{lem:probasA}(b) do not depend on the accuracy, and each $k'$ belongs to at most two intervals $[\alpha_k-1,\alpha'_k]$. By Lemma~\ref{lem:sumsA}(a)
($\alpha_k-1\ge\lonep-1\ge\lA$), these parts sum, over $k$, to at most $2(1+K)\Ceta\sig_{5\beta}(\eg)$.

\stp6{the other terms}\emph{A bound in terms of $y:=\Lg_i$.} Let $i\in I_k$, $y:=\Lg_i\ge1$ and $\varepsilon:=e_k=\epl_i\in[\etwo y^{-1/4},\eb]$. Let $\kap_0:=\kap(\m^\star+\eb)$,
$w:=1-\m^\star-\eb\ge\frac{1-\m^\star}2$, and
\begin{gather*}
C_3:=2K\Big(1+\frac1{2\etwo^2}\Big),\quad C_5:=3+\frac{64}{\etwo^2w},\quad C_6:=\kap_0+\Big|\log\frac8{\etwo}\Big|+1,\\
C_{10}:=1+\frac{2C_6^2}{\etwo^2},\quad a_1:=\frac{\etwo^2}{\xp^6},\quad b_1:=\frac{\etwo^2}{2C_6^2\kap_0\xp^2}.
\end{gather*}
We use $\frac1{1-e^{-v}}\le1+\frac1v$ for $v>0$, and $1-e^{-v}\ge\frac v2$ for $v\in[0,1]$.
\begin{itemize}
\item \emph{$\Ag^c$.} $\cAc(\varepsilon)\le2K(1+\frac1{2\varepsilon^2})\le C_3y^{1/2}$. Since $i-3\ge\lA$, $\Lg_{i-3}\ge1$ and $y\le\xp^3\Lg_{i-3}$; with $\beta\ge2$,
$\Ug_{i-2}\ge\Lg_{i-3}^2\ge y^2/\xp^6$, so $\varepsilon^2\Ug_{i-2}\ge a_1y^{3/2}$, and Lemma~\ref{lem:probasA}(a) gives
\[
\PP(\Ag_{i-2}(\varepsilon)^c)\ \le\ C_3\,y^{1/2}e^{-a_1y^{3/2}}.
\]
\item \emph{$\Lo$, $\Up$.} We bound the three factors of Lemma~\ref{lem:probasA}(c) in terms of $y$, using $\varepsilon\ge\etwo y^{-1/4}$.
\begin{itemize}
\item $1+\ngr_\varepsilon\le3+\dfrac{64}{\varepsilon^2w}\le C_5\,y^{1/2}$. \ (Since $1-e^{-\varepsilon/4}\ge\varepsilon/8$.)
\item $\Rt_\varepsilon\le C_6\,y^{1/8}$. \ (Since $-\log(1-e^{-\varepsilon/4})\le\log\frac8\varepsilon\le|\log\frac8{\etwo}|+\frac14\log y$, and $\frac14\log y\le y^{1/8}$.)
\item $\rt_\varepsilon=\dfrac{\varepsilon^2}{2\Rt_\varepsilon^2}\ge\dfrac{\etwo^2}{2C_6^2}\,y^{-3/4}$, hence $\dfrac1{1-e^{-\rt_\varepsilon}}\le1+\dfrac1{\rt_\varepsilon}\le C_{10}\,y^{3/4}$.
\item $\rt_\varepsilon\ug_i(\varepsilon)\ge b_1\,y^{1/4}$. \ (Since $\ug_i(\varepsilon)=\Lg_{i-2}/\kape_\varepsilon$, $\kape_\varepsilon\le\kap_0$ and $y\le\xp^2\Lg_{i-2}$.)
\end{itemize}
Lemma~\ref{lem:probasA}(c) gives
\[
\PP(\Lo_i(\varepsilon))+\PP(\Up_i(\varepsilon))\ \le\ C_5C_{10}\,y^{5/4}e^{-b_1y^{1/4}}.
\]
\end{itemize}
\emph{A non-increasing envelope.} The two bounds above give $\PP(\Ag_{i-2}^c)+\PP(\Lo_i)+\PP(\Up_i)\le\Phi(\Lg_i)$, where
\[
\Phi(y):=C_3\,y^{1/2}e^{-a_1y^{3/2}}+C_5C_{10}\,y^{5/4}e^{-b_1y^{1/4}} .
\]
To apply the layer cake we need a non-increasing bound. Let
\[
C_\Phi:=\sup_{y\ge1}\Phi(y)\,e^{b_1y^{1/4}/2}<\infty,\qquad\Psi(y):=C_\Phi\,e^{-\frac{b_1}2\max(y,1)^{1/4}} .
\]
Then $\Psi$ is non-increasing and $\Phi\le\Psi$ on $[1,\infty)$. Summing over the blocks, the non-$\Vg$ part of $\sum_k\E[Y^{(k)}]$ is at most
\[
\underbrace{\sum_{i=\lonep+1}^j(v_i-v_{i-1})\Psi(\Lg_i)}_{\text{interior}}+\underbrace{\sum_{k=1}^r\Ug_{\alpha_k+1}\Psi(\Lg_{\alpha_k})}_{\text{block starts}} .
\]
\emph{Interior.} Adding the non-negative term $v_{\lonep}\Psi(\Lg_{\lonep})$, the interior is at most $\sum_{i=\lonep}^j(v_i-v_{i-1})\Psi(\Lg_i)$ with $v_{\lonep-1}:=0$. By
Lemma~\ref{lem:sums}(iv) (with $\lambda_i:=\Lg_i^\beta$, $\lambda_{\lonep-1}:=0$) and (iii) (with $p=4\beta$),
\[
\text{interior}\ \le\ \sum_{n\ge0}\Psi(n^{1/\beta})\ \le\ C_\Phi\big(1+\Gamma(4\beta+1)(2/b_1)^{4\beta}\big).
\]
\emph{Block starts.} Let $C_\chi:=C_\Phi\sup_{y\ge1}(1+y^\beta)e^{-b_1y^{1/4}/4}<\infty$, so that $(1+y^\beta)\Psi(y)\le C_\chi e^{-b_1y^{1/4}/4}$. For $k\ge2$,
$\epl_{\alpha_k}<\epl_{\alpha_k-1}$ forces $\ody_{\alpha_k}<\ody_{\alpha_k-1}$; so $p_k:=k_{\alpha_k}\ge1$, $\Lg_{\alpha_k}\ge16^{p_k}$ (Step~0), and the $p_k$ are distinct. With
$\Ug_{\alpha_k+1}\le1+\Lg_{\alpha_k}^\beta$ and $2^m\ge2m$,
\[
\text{block starts}\ \le\ C_\chi+\sum_{m\ge1}\sup_{y\ge16^m}C_\chi e^{-b_1y^{1/4}/4}\ \le\ C_\chi+C_\chi\sum_{m\ge1}e^{-b_1m/2}\ =\ C_\chi\Big(1+\frac1{e^{b_1/2}-1}\Big).
\]
\emph{Total.} Steps~5 and~6 give $\E[Y'_j]\le\mathcal R'$, independent of $j$, with
\[
\mathcal R':=C_\Phi\big(1+\Gamma(4\beta+1)(2/b_1)^{4\beta}\big)+C_\chi\Big(1+\frac1{e^{b_1/2}-1}\Big)+2(1+K)\Ceta\sig_{5\beta}(\eg).
\]

\stp7{conclusion}By~\eqref{eq:main2A}, for $j\ge i_\xi$ and $t\in(\tl_j,\tl_{j+1}]$, $\E[\Ng_a(t)]\le K_0+\varphi\HH(t)/\da$, with
\[
K_0:=\tl_{\lonep}+K+\frac52\frac{\HH_{i_\xi}}{\da}+2\mathcal R'+\cAc(\eb)e^{\SE_{\eb}}+(1+K)\Ceta\sig_{4\beta}(\eg)
\]
independent of $t$. By Lemma~\ref{lem:H}, $\limsup_t\E[\Ng_a(t)]/\log t\le\varphi x^2/\da<qx^2/\da=\theta'$. Since $\theta'>x^2/\da$ is arbitrary, the proposition follows.
\end{proof}
%----------------------------------------------------------------------
\subsection{Proof of Lemma~\ref{lem:decompA}}\label{sec:pf-lem:decompA}
\noindent\emph{Idea.} (a) is a case analysis on the batch: on a good batch, Lemma~\ref{lem:goodA}; otherwise, if the leader is still optimal, $a$ receives at most a net;
if not, the whole batch is charged to the event $\Fg_\ell$. (b) iterates (a) over the batches, keeping only the largest net and the plan of the last batch.

\begin{proof}
\proves{lem:decompA}
\uses{lem:structA, lem:goodA, lem:leaderA, lem:imedround, def:grid, def:evA, def:learner, def:hy, def:constA, def:ranksA}
\stp{(a)}{three cases}Fix $t\in(\tl_\ell,\tl_{\ell+1}]$.

\noindent\emph{Case 1: on $\Gg_\ell(\varepsilon)$.} Lemma~\ref{lem:goodA} gives $\Ng_a(t)\le\max(\Ng_a(\tl_\ell),\mt_\ell(\varepsilon,t))$.

\noindent\emph{Case 2: on $\Gg_\ell(\varepsilon)^c\cap\Fg_\ell(\varepsilon')^c$.} By Lemma~\ref{lem:leaderA} at the accuracy $\varepsilon'$ ((T3) at $\ell+2$ gives $2\wz\le\varepsilon'\Ug_\ell$),
$\Ld_\ell$ is optimal, so $a\ne\Ld_\ell$, and $\Kt_\ell(a)\ge\varepsilon'^2/8$.
\begin{itemize}
\item IMED rounds, $t\le\sA_\ell$:
\begin{align*}
\Ng_a(t)\ &\le\ \max\Big(\Ng_a(\tl_\ell),\ 1+\frac{8\yg_\ell\log t}{\varepsilon'^2}\Big)&&\why{Lemma~\ref{lem:imedround} (\emph{IMED rounds})}\\
&\le\ \max\Big(\Ng_a(\tl_\ell),\ 1+\frac{8\xp^2\Lg_{\ell+1}}{\varepsilon'^2}\Big)&&\why{$\log t\le\Lg_{\ell+1}$, $\yg_\ell\le\xp^2$ (R1)}\\
&\le\ \max\big(\Ng_a(\tl_\ell),\ \Ug_{\ell+1}\big)&&\why{(T4) at $\ell$}.
\end{align*}
\item Paced rounds, $t>\sA_\ell$:
\begin{align*}
\Ng_a(t)\ &\le\ \max\big(\Ng_a(\sA_\ell),\ \Ug_{\ell+1}+K-1\big)\\
&\le\ \max\big(\Ng_a(\tl_\ell),\ 2\Ug_{\ell+1}\big)\qquad\why{(S1), $\nt_{\ell,a}\le\Ug_{\ell+1}+K-1$; $K-1\le\Ug_{\ell+1}$ by (T4)}.
\end{align*}
\end{itemize}

\noindent\emph{Case 3: on $\Fg_\ell(\varepsilon')$.} $\Ng_a(t)\le\Ng_a(\tl_\ell)+(t-\tl_\ell)$.

\noindent In the three cases~\eqref{eq:DA} holds.

\stp{(b)}{induction over the batches}Let $S_k:=\sum_{i=i_0}^k\Bg_{i+1}\1(\Fg_i(\varepsilon'))$, $\Yg_k$ as in the statement, and $S_{i_0-1}:=0$, $\Yg_{i_0-1}:=0$.

\noindent\emph{Claim:} for $k\in[i_0,j]$, $\ \Ng_a(\tl_k)\le\max\big(\Ng_a(\tl_{i_0}),\bar M,\Yg_{k-1}\big)+S_{k-1}$.

\noindent It is trivial for $k=i_0$. If it holds at $k<j$, then, using~\eqref{eq:DA} at $k$ and $t=\tl_{k+1}$, $\mt_k(\varepsilon_k,\tl_{k+1})\le\bar M$, and $\max(p+s,q)\le\max(p,q)+s$ for $s\ge0$,
\begin{align*}
\Ng_a(\tl_{k+1})\ &\le\ \max\big(\Ng_a(\tl_k),\ \bar M,\ 2\Ug_{k+1}\1(\Gg_k(\varepsilon_k)^c)\big)+\Bg_{k+1}\1(\Fg_k(\varepsilon'))&&\why{\eqref{eq:DA}}\\
&\le\ \max\big(\max(\Ng_a(\tl_{i_0}),\bar M,\Yg_{k-1})+S_{k-1},\ \bar M,\ 2\Ug_{k+1}\1(\Gg_k(\varepsilon_k)^c)\big)+\Bg_{k+1}\1(\Fg_k(\varepsilon'))&&\why{induction}\\
&\le\ \max\big(\Ng_a(\tl_{i_0}),\ \bar M,\ \Yg_k\big)+S_k .
\end{align*}

\noindent\emph{Conclusion.} For $t\in(\tl_j,\tl_{j+1}]$,~\eqref{eq:DA} at $j$, the claim at $k=j$ and $t-\tl_j\le\Bg_{j+1}$ give
$\Ng_a(t)\le\max(\Ng_a(\tl_{i_0}),\bar M,\mt_j(\varepsilon_j,t),\Yg_j)+S_j$. Finally $\max(p,q,r)\le p+q+r$ for non-negative numbers, and $\Ng_a(\tl_{i_0})\le\tl_{i_0}$.
\end{proof}

%----------------------------------------------------------------------
\subsection{Proof of Lemma~\ref{lem:probasA}}\label{sec:pf-lem:probasA}
\begin{proof}
\proves{lem:probasA}
\uses{lem:stacks, lem:hoeffding, lem:constA, lem:lowerdev, lem:upperdev, def:hy, def:constA, def:grid, def:evA, def:kinf, def:model, def:ranksA}
All the events are events on the stacks, which are i.i.d.\ in $[0,1]$ (Lemma~\ref{lem:stacks}).

\stp{(a)}{Hoeffding}Lemma~\ref{lem:hoeffding}(b) for each arm, with $\rho=\varepsilon$ and $u=\frac12\Ug_\ell$, gives $2e^{-\varepsilon^2\Ug_\ell}/(1-e^{-2\varepsilon^2})$; a union bound
over the $K$ arms gives $\cAc(\varepsilon)e^{-\varepsilon^2\Ug_\ell}$.

\stp{(b)}{gains}By (R3) at $\ell-1\ge\lA-3$, $(1+2\eg)z_\ell\ge(1+\eg)\Lg_{\ell+1}$. For $\Vg_\ell$, a union bound over the integers $n<\frac12\Ug_\ell$ and
Lemma~\ref{lem:constA}(b) give
\[
\PP(\Vg_\ell)\ \le\ \sum_{n<\Ug_\ell/2}\Ceta n^2e^{-(1+2\eg)z_\ell}\ \le\ \Ceta\,\Ug_\ell^3\,e^{-(1+\eg)\Lg_{\ell+1}}\ =\ \Ceta\,\Ug_\ell^3\,\tl_{\ell+1}^{-(1+\eg)} .
\]
For $\Pig_\ell$: a union bound over the $K-1$ sub-optimal arms and the at most $\Ug_\ell^2$ pairs $(n,n')$, and Lemma~\ref{lem:constA}(c).

\stp{(c)}{deviations of the empirical $\Ki$}\emph{$\Lo_\ell$.} Lemma~\ref{lem:lowerdev} with $\mu=\m^\star-2\varepsilon$ and $\varepsilon'=\varepsilon$ applies; its $R'$ is at most
$\Rt_\varepsilon$ ($\kap$ is non-decreasing), so its $\rho'$ is at least $\rt_\varepsilon$. Since $\De_\varepsilon=\Ki(\nu_a,\m^\star-2\varepsilon)-\varepsilon$, a union bound over the
integers $N\ge\ug_\ell$ gives $\PP(\Lo_\ell)\le e^{-\rt_\varepsilon\ug_\ell}/(1-e^{-\rt_\varepsilon})$.

\noindent\emph{$\Up_\ell$.} Lemma~\ref{lem:upperdev} with $\mu=\m^\star+\varepsilon$ and $\varepsilon'=\varepsilon$, whose constants are exactly $\Rt_\varepsilon$, $\rt_\varepsilon$,
$\ngr_\varepsilon$, and a union bound give $\PP(\Up_\ell)\le\ngr_\varepsilon e^{-\rt_\varepsilon\ug_\ell}/(1-e^{-\rt_\varepsilon})$. Add the two bounds.
\end{proof}

%----------------------------------------------------------------------
\subsection{Proof of Lemma~\ref{lem:sumsA}}\label{sec:pf-lem:sumsA}
\begin{proof}
\proves{lem:sumsA}
\uses{lem:probasA, lem:sums, lem:ranksA, def:evA, def:constA, def:hy, def:grid, def:ranksA}
\stp{(a)}{nets lost on $\Vg$, $\Pig$}Let $k\ge\lA$. By Lemma~\ref{lem:probasA}(b), $\PP(\Vg_k)+\PP(\Pig_k)\le(1+K)\Ceta\,\Ug_k^4\,\tl_{k+1}^{-(1+\eg)}$. Moreover
$\Ug_{k+2}=\lceil\Lg_{k+1}^\beta\rceil\le(1+\Lg_{k+1})^\beta$, $\ \Ug_k\le(1+\Lg_{k+1})^\beta$ and $\Bg_{k+1}\ge1$. So
\[
\Ug_{k+2}\big(\PP(\Vg_k)+\PP(\Pig_k)\big)\ \le\ (1+K)\Ceta\,\Bg_{k+1}(1+\Lg_{k+1})^{5\beta}\,\tl_{k+1}^{-(1+\eg)},
\]
and Lemma~\ref{lem:sums}(ii) bounds the sum over $k$ by $(1+K)\Ceta\,\sig_{5\beta}(\eg)$.

\stp{(b)}{whole batches}Let $i\ge\lA$. Recall $\Fg_i=\Ag_{i-1}^c\cup\Ag_i^c\cup\Vg_i\cup\Pig_i$.

\noindent\emph{The events $\Ag^c$.} By Lemma~\ref{lem:probasA}(a) and $\Ug_i\ge\Ug_{i-1}\ge\Lg_{i-2}^\beta$, $\PP(\Ag_{i-1}^c)+\PP(\Ag_i^c)\le2\cAc(\varepsilon)e^{-\varepsilon^2\Lg_{i-2}^\beta}$. As
$\Lg_{i+1}\le\xp^3\Lg_{i-2}$ (Lemma~\ref{lem:ranksA}),
\[
\Bg_{i+1}e^{-\varepsilon^2\Lg_{i-2}^\beta}\ =\ \Bg_{i+1}\tl_{i+1}^{-2}\,e^{2\Lg_{i+1}-\varepsilon^2\Lg_{i-2}^\beta}\ \le\ \Bg_{i+1}\tl_{i+1}^{-2}\,e^{2\xp^3\Lg_{i-2}-\varepsilon^2\Lg_{i-2}^\beta}\ \le\ e^{\SE_\varepsilon}\,\Bg_{i+1}\tl_{i+1}^{-2},
\]
and $\sum_i\Bg_{i+1}\tl_{i+1}^{-2}\le\tl_1^{-1}\le\frac12$ (Lemma~\ref{lem:sums}(i) with $\theta=1$). This part contributes at most $\cAc(\varepsilon)e^{\SE_\varepsilon}$.

\noindent\emph{The events $\Vg$, $\Pig$.} As in (a), $\Bg_{i+1}(\PP(\Vg_i)+\PP(\Pig_i))\le(1+K)\Ceta\,\Bg_{i+1}(1+\Lg_{i+1})^{4\beta}\tl_{i+1}^{-(1+\eg)}$, whose sum is at most
$(1+K)\Ceta\,\sig_{4\beta}(\eg)$ by Lemma~\ref{lem:sums}(ii).
\end{proof}

%----------------------------------------------------------------------
\subsection{Proof of Lemma~\ref{lem:structA}}\label{sec:pf-lem:structA}
\noindent\emph{Idea.} By (R2) the paced phase is long enough for the queue lemma: every demand is served up to $K-1$ pulls, and the leader keeps the rest. An arm that
ends the batch with fewer than $\frac12\Ug_{\ell+1}$ pulls cannot have been ambiguous (it would have received the net), so it was rejected and served its uncapped plan.

\par\smallskip
\Setting $\ell\ge\lA-4$, so that (R2) holds at $\ell+1$: $\ell\in\Act$, the \term{s:act}{paced phase} has $\tl_{\ell+1}-\sA_\ell\ge K(\Ug_{\ell+1}+K)\ge(K-1)(\Ug_{\ell+1}+K-1)$ rounds,
and the condition $\sA_\ell\log\sA_\ell\ge2(K-1)\Ug_{\ell+1}$ of Lemma~\ref{lem:service}(b) (\emph{queue lemma}) holds.
\Goal (S1)--(S4).
\par\smallskip

\begin{proof}
\proves{lem:structA}
\uses{lem:service, lem:kinf, def:grid, def:hy, def:learner, def:kinf}
\stp{(S1)}{service}This is Lemma~\ref{lem:service}(a) (\emph{queue lemma}), and Lemma~\ref{lem:service}(b) at $t=\tl_{\ell+1}$.

\stp{(S2)}{the leader}By Lemma~\ref{lem:service}(c) (\emph{queue lemma}), the number of pulls of $\Ld_\ell$ in the paced phase is at least
\[
(\tl_{\ell+1}-\sA_\ell)-(K-1)(\Ug_{\ell+1}+K-1)\ \ge\ K(\Ug_{\ell+1}+K)-(K-1)(\Ug_{\ell+1}+K-1)\ =\ \Ug_{\ell+1}+2K-1\ \ge\ \Ug_{\ell+1}.
\]
So $\Ng_{\Ld_\ell}(\tl_{\ell+1})\ge\frac12\Ug_{\ell+1}$, i.e.\ $\Ld_\ell\in\Wg_{\ell+1}$.

\stp{(S3)}{a poorly sampled arm}Let $c\notin\Wg_{\ell+1}$, i.e.\ $\Ng_c(\tl_{\ell+1})<\frac12\Ug_{\ell+1}$. Then $c\ne\Ld_\ell$ by (S2).
\begin{itemize}
\item If $c$ were ambiguous, or rejected with a capped target, then $\nt_{\ell,c}(\tl_{\ell+1})=K-1+\Ug_{\ell+1}$, and (S1) would give $\Ng_c(\tl_{\ell+1})\ge\Ug_{\ell+1}$: a
contradiction.
\item So $c$ is rejected with an uncapped target. If $\Kh_\ell(c)=\infty$ the claim is trivial. Otherwise (S1) gives
\[
\Ng_c(\tl_{\ell+1})\ \ge\ \nt_{\ell,c}(\tl_{\ell+1})-(K-1)\ =\ \Big\lceil\frac{(1+\om_\ell)\yg_\ell\Lg_{\ell+1}}{\Kh_\ell(c)}\Big\rceil\ \ge\ \frac{(1+\om_\ell)\yg_\ell\Lg_{\ell+1}}{\Kh_\ell(c)} .
\]
\end{itemize}
\emph{The reference.} If $\rr_\ell=\Ld_\ell$, it is in $\Wg_{\ell+1}$ by (S2). Otherwise $\Kh_\ell(\rr_\ell)=0$, since the test level $\mh_\ell$ is the mean of $\tnu_{\rr_\ell}(\tl_\ell)$
(Lemma~\ref{lem:kinf}(ii), or the convention at level $1$); so $\rr_\ell$ is not rejected, and by what precedes it is in $\Wg_{\ell+1}$.

\stp{(S4)}{every arm is explored}Let $\mh_\ell<1$.
\begin{itemize}
\item For $\Ld_\ell$: $\Ng\ge\Ug_{\ell+1}$ by (S2).
\item For an ambiguous arm or a rejected arm with a capped target: $\Ng\ge\Ug_{\ell+1}$ by (S1), as in (S3).
\item For a rejected arm with an uncapped target: by the computation of (S3) and Lemma~\ref{lem:kinf}(i) ($\Kh_\ell(c)\le\kap(\mh_\ell)$),
\[
\Ng_c(\tl_{\ell+1})\ \ge\ \frac{(1+\om_\ell)\yg_\ell\Lg_{\ell+1}}{\Kh_\ell(c)}\ \ge\ \frac{\yg_\ell\Lg_{\ell+1}}{\kap(\mh_\ell)} .
\]
\end{itemize}
\end{proof}

%----------------------------------------------------------------------
\subsection{Proof of Lemma~\ref{lem:thrA}}\label{sec:pf-lem:thrA}
\begin{proof}
\proves{lem:thrA}
\uses{lem:structA, def:hy, def:learner, def:grid, def:constA, def:ranksA}
Let $\ell\ge\lA-3$ (so $\ell\ge2$) and $c\in\Qp_\ell$, i.e.\ $c\notin\Wg_\ell$ and $c\ne\Ld_{\ell-1}$. By (S3) at $\ell-1\ge\lA-4$,
\[
\Ng_c(\tl_\ell)\,\Kh_{\ell-1}(c)\ \ge\ (1+\om_{\ell-1})\,\yg_{\ell-1}\Lg_\ell\ \ge\ \yg_{\ell-1}\Lg_\ell .
\]
So the bracket of $c$ in the definition of $\Lh_\ell$ is at least $\Lg_\ell$ (as $\log\Ng_c\ge0$). Hence $\Lh_\ell=\min(\Lg_\ell,\dots)=\Lg_\ell$, and $\hg_\ell=(1-\gam_\ell)\yg_{\ell-1}\Lg_\ell=\hb_\ell$.
\end{proof}

%----------------------------------------------------------------------
\subsection{Proof of Lemma~\ref{lem:leaderA}}\label{sec:pf-lem:leaderA}
\noindent\emph{Idea.} A sub-optimal leader comes either from a test made against a bad level (no optimal arm is well sampled: event $\Vg_\ell$), or from a poorly
sampled sub-optimal arm whose mean climbed after it was rejected at the previous test (event $\Pig_\ell$).

\par\smallskip
\Setting $\varepsilon\in(0,\eb]$, $\ell\ge\lA$ with $2\wz\le\varepsilon\Ug_\ell$, and an outcome in $\Ag_{\ell-1}\cap\Ag_\ell\cap\Vg_\ell^c\cap\Pig_\ell^c$ (only $\Ag_{\ell-1}$
for the last claim).
\Goal (i), (ii), (iii) and the last claim.
\par\smallskip

\begin{proof}
\proves{lem:leaderA}
\uses{lem:structA, lem:pseudo, lem:thrA, lem:kinf, lem:constA, def:hy, def:evA, def:grid, def:stacks, def:model, def:learner, def:constA, def:kinf}
\stp0{the previous test and the rejected arms}\emph{The previous test level.} $\Wg_{\ell-1}$ and $\Wg_\ell$ are non-empty, since they contain $\Ld_{\ell-2}$ and $\Ld_{\ell-1}$
((S2) at $\ell-2$ and $\ell-1$). On $\Ag_{\ell-1}$, every arm of $\Wg_{\ell-1}$ has an $\varepsilon$-accurate mean at $\tl_{\ell-1}$, so
\[
\mh_{\ell-1}\ \le\ \hm_{\rr_{\ell-1}}(\tl_{\ell-1})\ \le\ \m_{\rr_{\ell-1}}+\varepsilon\ \le\ \m^\star+\varepsilon\qquad\why{Lemma~\ref{lem:pseudo}(iii); $\rr_{\ell-1}\in\Wg_{\ell-1}$}.
\]
This is the last claim.

\noindent\emph{A rejected arm.} Let $c$ be rejected at $\tl_{\ell-1}$ and $n:=\Ng_c(\tl_{\ell-1})$. The threshold is $\hg_{\ell-1}=\hb_{\ell-1}$ (Lemma~\ref{lem:thrA}). For every
$r\in[\mh_{\ell-1},1)$,
\begin{equation}\label{eq:rejA}
\begin{aligned}
n\Ki(\hnu_{c,n},r)\ &\ge\ n\Ki(\hnu_{c,n},\mh_{\ell-1})&&\why{Lemma~\ref{lem:kinf}(iii)}\\
&\ge\ n\Ki(\tnu_c(\tl_{\ell-1}),\mh_{\ell-1})-\wz\kap(\mh_{\ell-1})&&\why{Lemma~\ref{lem:pseudo}(ii)}\\
&\ge\ \hb_{\ell-1}-\wz\kap(\m^\star+\eb)\ =\ z_\ell&&\why{$c$ rejected; $\mh_{\ell-1}\le\m^\star+\varepsilon$}.
\end{aligned}
\end{equation}

\stp{(i)}{an optimal arm is well sampled}Assume, by contradiction, that no optimal arm lies in $\Wg_\ell$. By (S2) and (S3) at $\ell-1$, $\Ld_{\ell-1}$ and $\rr_{\ell-1}$ are in
$\Wg_\ell$, so both are sub-optimal, and on $\Ag_{\ell-1}$, $\mh_{\ell-1}\le\m_{\rr_{\ell-1}}+\varepsilon\le\mtwo+\varepsilon$. Now $a^\star\notin\Wg_\ell$ and $a^\star\ne\Ld_{\ell-1}$, so
$a^\star$ was rejected at $\tl_{\ell-1}$ by (S3). Then~\eqref{eq:rejA} with $r=\mtwo+\varepsilon$ gives $n\Ki(\hnu_{a^\star,n},\mtwo+\varepsilon)\ge z_\ell$, with
$1\le n\le\Ng_{a^\star}(\tl_\ell)<\frac12\Ug_\ell$: the outcome is in $\Vg_\ell$, a contradiction.

\noindent So some optimal $a'$ lies in $\Wg_\ell$. On $\Ag_\ell$:
\begin{itemize}
\item $\tm_{a'}(\tl_\ell)\ge\hm_{a'}(\tl_\ell)-2\wz/\Ug_\ell\ge\m^\star-\varepsilon-\varepsilon$ (Lemma~\ref{lem:pseudo}, and $2\wz\le\varepsilon\Ug_\ell$);
\item every sub-optimal $c\in\Wg_\ell$ has $\tm_c(\tl_\ell)\le\hm_c(\tl_\ell)\le\mtwo+\varepsilon<\m^\star-2\varepsilon$ (Lemma~\ref{lem:constA}(a)).
\end{itemize}
Hence $\rr_\ell$ is optimal and $\m^\star-2\varepsilon\le\tm_{a'}(\tl_\ell)\le\mh_\ell\le\hm_{\rr_\ell}(\tl_\ell)\le\m^\star+\varepsilon$.

\stp{(ii)}{the sub-optimal arms}Let $c$ be sub-optimal.
\begin{itemize}
\item If $c\in\Wg_\ell$: $\hm_c(\tl_\ell)\le\m_c+\varepsilon<\m^\star-3\varepsilon$ on $\Ag_\ell$ (Lemma~\ref{lem:constA}(a)).
\item If $c\notin\Wg_\ell$, assume by contradiction $\hm_c(\tl_\ell)\ge\m^\star-3\varepsilon$.
\begin{itemize}
\item $c\ne\Ld_{\ell-1}$, since $\Ld_{\ell-1}\in\Wg_\ell$ by (S2); so $c$ was rejected at $\tl_{\ell-1}$, by (S3).
\item Apply~\eqref{eq:rejA} with $r=\m^\star+\varepsilon$, which is $\ge\mh_{\ell-1}$ by Step~0. With $n:=\Ng_c(\tl_{\ell-1})$ and $n':=\Ng_c(\tl_\ell)$, we have
$1\le n\le n'<\frac12\Ug_\ell$, $n\Ki(\hnu_{c,n},\m^\star+\varepsilon)\ge z_\ell$ and $\hm_{c,n'}\ge\m^\star-3\varepsilon$.
\end{itemize}
So the outcome is in $\Pig_\ell$: a contradiction.
\end{itemize}
So $\hm_c(\tl_\ell)<\m^\star-3\varepsilon$. Then
\begin{align*}
\Kh_\ell(c)\ &\ge\ \Ki(\hnu_c(\tl_\ell),\mh_\ell)&&\why{Lemma~\ref{lem:pseudo}(i)}\\
&\ge\ \Ki(\hnu_c(\tl_\ell),\m^\star-2\varepsilon)&&\why{Lemma~\ref{lem:kinf}(iii), $\mh_\ell\ge\m^\star-2\varepsilon$}\\
&\ge\ \varepsilon^2/8&&\why{Lemma~\ref{lem:kinf}(iv), mean $<\m^\star-3\varepsilon$}.
\end{align*}
Finally $\tms_\ell\ge\mh_\ell$ (the leader maximizes $\tm$ over all arms), so $\Kt_\ell(c)\ge\Kh_\ell(c)$ by Lemma~\ref{lem:kinf}(iii) (or $\Kt_\ell(c)=\infty$ if $\tms_\ell=1$).

\stp{(iii)}{the leader}$\tms_\ell\ge\mh_\ell\ge\m^\star-2\varepsilon$, while every sub-optimal $c$ has $\tm_c(\tl_\ell)\le\hm_c(\tl_\ell)<\m^\star-3\varepsilon$ by (ii). So $\Ld_\ell$ is optimal.
\end{proof}

%----------------------------------------------------------------------
\subsection{Proof of Lemma~\ref{lem:countsA}}\label{sec:pf-lem:countsA}
\begin{proof}
\proves{lem:countsA}
\uses{lem:structA, lem:leaderA, def:grid, def:evA, def:hy, def:model, def:learner, def:kinf, def:constA, def:ranksA}
\stp1{at $\tl_\ell$}On $\Ag_{\ell-1}$, $\mh_{\ell-1}\le\m^\star+\varepsilon<1$ (last claim of Lemma~\ref{lem:leaderA}, applied at the index $\ell\ge\lA$). By (S4) at $\ell-1$,
\begin{align*}
\Ng_a(\tl_\ell)\ &\ge\ \min\Big(\Ug_\ell,\ \frac{\yg_{\ell-1}\Lg_\ell}{\kap(\mh_{\ell-1})}\Big)\ \ge\ \min\Big(\Lg_{\ell-1}^\beta,\ \frac{\Lg_\ell}{\kape_\varepsilon}\Big)&&\why{$\Ug_\ell\ge\Lg_{\ell-1}^\beta$, $\yg\ge1$, $\kap(\mh_{\ell-1})\le\kape_\varepsilon$}\\
&\ge\ \min\Big(\Lg_{\ell-1}^\beta,\ \frac{\Lg_{\ell-1}}{\kape_\varepsilon}\Big)\ =\ \frac{\Lg_{\ell-1}}{\kape_\varepsilon}&&\why{(T1) at $\ell+1$}\\
&\ge\ \frac{\Lg_{\ell-2}}{\kape_\varepsilon}\ =\ \ug_\ell(\varepsilon).
\end{align*}

\stp2{at $\tl_{\ell-1}$}On $\Ag_{\ell-2}$, the same computation at $\ell-2$ (with (T1) at $\ell$) gives $\Ng_a(\tl_{\ell-1})\ge\Lg_{\ell-2}/\kape_\varepsilon=\ug_\ell(\varepsilon)$.
\end{proof}

%----------------------------------------------------------------------
\subsection{Proof of Lemma~\ref{lem:goodA}}\label{sec:pf-lem:goodA}
\noindent\emph{Idea.} The paced plan of batch $\ell$, evaluated at the actual end $\tl_\ell$, aimed above the present threshold; the margin $1+\om_{\ell-1}$ absorbs the change
of the empirical $\Ki$ between the two tests. In the IMED phase, $a$ is explored up to $\yg_\ell\log t$ only, below its paced target.

\par\smallskip
\Setting $\varepsilon\in(0,\eb]$ with $\De_\varepsilon>0$, $\ell\ge\lone(\varepsilon,\varepsilon')$, $(1+\om_{\ell-1})\De_\varepsilon\ge\Kp_\varepsilon$, and an outcome in $\Gg_\ell(\varepsilon)$.
\Goal the three points of the lemma.
\par\smallskip

\begin{proof}
\proves{lem:goodA}
\uses{lem:structA, lem:leaderA, lem:countsA, lem:pseudo, lem:kinf, lem:thrA, lem:imedround, def:hy, def:model, def:learner, def:grid, def:evA, def:kinf, def:stacks, def:constA, def:ranksA}
\stp1{what the good events give}
\begin{itemize}
\item By Lemma~\ref{lem:leaderA} at $\ell-1$ and at $\ell$ ((T3) at $\ell+1$ gives $2\wz\le\varepsilon\Ug_{\ell-1}\le\varepsilon\Ug_\ell$): $\Ld_{\ell-1}$ and $\Ld_\ell$ are optimal,
$\mh_\ell\ge\m^\star-2\varepsilon$ and $\mh_{\ell-1}\le\m^\star+\varepsilon$.
\item By Lemma~\ref{lem:countsA}: $\Ng_a(\tl_\ell)\ge\ug_\ell$ and $\Ng_a(\tl_{\ell-1})\ge\ug_\ell$.
\item Outside $\Lo_\ell$: $\Kh_\ell(a)\ge\Ki(\hnu_a(\tl_\ell),\mh_\ell)\ge\Ki(\hnu_a(\tl_\ell),\m^\star-2\varepsilon)\ge\De_\varepsilon>0$ (Lemma~\ref{lem:pseudo}(i), then
Lemma~\ref{lem:kinf}(iii)).
\end{itemize}

\stp2{$a$ is rejected at $\tl_\ell$}$a\ne\Ld_\ell$, $\Kh_\ell(a)>0$, and the threshold is $\hg_\ell=\hb_\ell\le\yg_{\ell-1}\Lg_\ell$ (Lemma~\ref{lem:thrA}). It remains to show
$\Zg_\ell(a)=\Ng_a(\tl_\ell)\Kh_\ell(a)\ge\yg_{\ell-1}\Lg_\ell$.

\noindent\emph{Case $a\in\Wg_\ell$.} Then $\Ng_a(\tl_\ell)\ge\frac12\Ug_\ell\ge\frac12\Lg_{\ell-1}^\beta$, and
\begin{align*}
\Zg_\ell(a)\ &\ge\ \tfrac12\Lg_{\ell-1}\cdot\Lg_{\ell-1}^{\beta-1}\De_\varepsilon\ \ge\ \xp^3\Lg_{\ell-1}&&\why{(T2) at $\ell$}\\
&\ge\ \xg_{\ell-1}^2\xg_\ell\Lg_{\ell-1}\ =\ \yg_{\ell-1}\Lg_\ell&&\why{(R1)}.
\end{align*}

\noindent\emph{Case $a\notin\Wg_\ell$.} Then $a\ne\Ld_{\ell-1}$, and (S3) at $\ell-1$ gives $\Ng_a(\tl_\ell)\Kh_{\ell-1}(a)\ge(1+\om_{\ell-1})\yg_{\ell-1}\Lg_\ell$ with $\Kh_{\ell-1}(a)>0$.
Outside $\Up_\ell$, the previous divergence of $a$ is at most $\Kp_\varepsilon$:
\begin{align*}
\Kh_{\ell-1}(a)\ &\le\ \Ki(\tnu_a(\tl_{\ell-1}),\m^\star+\varepsilon)&&\why{Lemma~\ref{lem:kinf}(iii), $\mh_{\ell-1}\le\m^\star+\varepsilon$}\\
&\le\ \Ki(\hnu_a(\tl_{\ell-1}),\m^\star+\varepsilon)+\frac{\wz\kap(\m^\star+\varepsilon)}{\Ng_a(\tl_{\ell-1})}&&\why{Lemma~\ref{lem:pseudo}(ii)}\\
&\le\ \big[\Ki(\nu_a,\m^\star+\varepsilon)+\varepsilon\big]+\varepsilon\ =\ \Kp_\varepsilon&&\why{$\Up_\ell^c$; (T3), $\Ng_a(\tl_{\ell-1})\ge\ug_\ell$}.
\end{align*}
Hence
\begin{align*}
\Zg_\ell(a)\ =\ \Ng_a(\tl_\ell)\Kh_\ell(a)\ &\ge\ \frac{(1+\om_{\ell-1})\yg_{\ell-1}\Lg_\ell}{\Kh_{\ell-1}(a)}\cdot\De_\varepsilon&&\why{(S3) at $\ell-1$; Step~1}\\
&\ge\ \frac{(1+\om_{\ell-1})\De_\varepsilon}{\Kp_\varepsilon}\,\yg_{\ell-1}\Lg_\ell\ \ge\ \yg_{\ell-1}\Lg_\ell&&\why{assumption on $\om_{\ell-1}$}.
\end{align*}

\stp3{the counts in the IMED phase}$\Kt_\ell(a)\ge\Kh_\ell(a)\ge\De_\varepsilon$ (as $\tms_\ell\ge\mh_\ell$). Lemma~\ref{lem:imedround} (\emph{IMED rounds}) gives, for every IMED round $t\le\sA_\ell$,
\[
\Ng_a(t)\ \le\ \max\Big(\Ng_a(\tl_\ell),\ 1+\frac{\yg_\ell\log t}{\De_\varepsilon}\Big).
\]

\stp4{the counts in the paced phase}For $t>\sA_\ell$, by (S1), Step~3 at $\sA_\ell$, and the target of the rejected arm $a$ ($\Kh_\ell(a)\ge\De_\varepsilon$),
\begin{align*}
\Ng_a(t)\ &\le\ \max\big(\Ng_a(\sA_\ell),\ \nt_{\ell,a}(t)\big)\\
&\le\ \max\Big(\Ng_a(\tl_\ell),\ 1+\frac{\yg_\ell\log\sA_\ell}{\De_\varepsilon},\ K-1+\Big\lceil\frac{(1+\om_\ell)\yg_\ell\log t}{\De_\varepsilon}\Big\rceil\Big)\\
&\le\ \max\Big(\Ng_a(\tl_\ell),\ K+\frac{(1+\om_\ell)\yg_\ell\log t}{\De_\varepsilon}\Big)&&\why{$\sA_\ell<t$}.
\end{align*}
Steps~3 and~4 give the bound at every round of batch $\ell+1$.
\end{proof}


%======================================================================
\section{Upper bound when \texorpdfstring{$x=1$}{x=1}}\label{sec:app-xone}

\par\smallskip
\Setting (in the whole section) $\xg_\ell\to1$; the margins $\om_\ell\in(0,1]$ are non-increasing with $\om_\ell\to0$ (for instance the dyadic margins $\ody$); $\thv$ and
$\gam$ are the defaults; $\wz\ge0$ and $\beta\ge2$.
\Goal Proposition~\ref{prop:fixedC}: for every accuracy $\varepsilon\in(0,\ez]$ and every $L\ge1$,
\[
\E[\Ng_a(\tl_L)]\ \le\ \frac{\Lg_L}{\da}+\Co\,\varepsilon\Lg_L+\RHo(\varepsilon),\qquad\RHo(\varepsilon)<\infty\ \text{independent of }L.
\]
\Plan the section goes from the top down.
\begin{enumerate}
\item[(1)] Constants and ranks (Subsection~\ref{sec:app-constC}).
\item[(2)] \emph{The skeleton} (Subsection~\ref{sec:app-skelC}): the proposition, the decomposition of $\Ng_a(\tl_L)$ batch by batch according to four \term{s:evC}{classes}, and the
proof of the proposition from it. The lemmas that follow answer the questions raised by the decomposition.
\item[(3)] \emph{Why are many batches high?} (Subsection~\ref{sec:app-water}). A batch $j<L$ is \term{t:high}{high} if, at its end, every arm has been explored almost up to the
level $\log\tl_L$ of a later horizon: either it has many observations, or a strong evidence $N\Kt$ against it. Since the grid grows slowly, many batches are high; this is
deterministic, and holds for \HY{} thanks to Lemma~\ref{lem:W} (\emph{the IMED phase lifts the water level}).
\item[(4)] \emph{Why are the leaders reliable, and why are wrong means rare?} (Subsection~\ref{sec:app-pathC}).
\item[(5)] \emph{Why is an optimal arm well sampled, and how many times is $a$ pulled in a batch?} (Subsection~\ref{sec:app-optW}).
\item[(6)] \emph{Why does $a$ pass the consistency test?} (Subsection~\ref{sec:app-amb}).
\end{enumerate}
\par\smallskip

\noindent\emph{Map.} A line ``$B\leftarrow A_1;\ A_2$'' means that $A_1$ and $A_2$ are used in the proof of $B$; the tools of Section~\ref{sec:app-tools}, the lemmas of Section~\ref{sec:app-det} other than Lemma~\ref{lem:W}, and Lemma~\ref{lem:constC} (used almost
everywhere) are not drawn.
\begin{itemize}
\item Prop.~\ref{prop:fixedC}, fixed accuracy $\leftarrow$ Lemma~\ref{lem:decompC}, decomposition; Lemmas~\ref{lem:controlC}, \ref{lem:sumC}, probabilities, sums; Lemma~\ref{lem:amb}, consistency test.
\item Lemma~\ref{lem:decompC}, decomposition $\leftarrow$ Lemma~\ref{lem:pullsC}, pulls of $a$; Lemma~\ref{lem:fewC}, few observations; Lemma~\ref{lem:logcountC}, logarithmic count; Lemma~\ref{lem:amb}, consistency test.
\item Lemma~\ref{lem:amb}, consistency test $\leftarrow$ Lemmas~\ref{lem:controlC}, \ref{lem:sumC}, probabilities, sums; Lemma~\ref{lem:optW}, optimal arm in $\Wg$; Lemma~\ref{lem:fewC}, few observations; Lemma~\ref{lem:logcountC}, logarithmic count.
\item Lemma~\ref{lem:pullsC}, pulls of $a$ $\leftarrow$ Lemma~\ref{lem:optW}, optimal arm in $\Wg$.
\item Lemma~\ref{lem:fewC}, few observations $\leftarrow$ Lemma~\ref{lem:goodhigh}, good high batches; Lemma~\ref{lem:highC}, evidence in a high batch.
\item Lemma~\ref{lem:logcountC}, logarithmic count $\leftarrow$ Lemma~\ref{lem:goodhigh}, good high batches.
\item Lemma~\ref{lem:optW}, optimal arm in $\Wg$ $\leftarrow$ Lemma~\ref{lem:goodhigh}, good high batches; Lemma~\ref{lem:highC}, evidence in a high batch.
\item Lemma~\ref{lem:goodhigh}, good high batches $\leftarrow$ Lemma~\ref{lem:levelC}, many high batches.
\item Lemma~\ref{lem:levelC}, many high batches $\leftarrow$ Lemma~\ref{lem:lowmass}, low batches; Lemma~\ref{lem:growth}, growth of the grid.
\item Lemma~\ref{lem:lowmass}, low batches $\leftarrow$ Lemma~\ref{lem:W}, IMED phase lifts, the water level.
\end{itemize}


%----------------------------------------------------------------------
\subsection{Constants and ranks}\label{sec:app-constC}

\begin{definition}[constants, $x=1$]\label{def:constC}\dt{s:constC}
\uses{def:model, def:kinf, def:constA, def:hy, def:grid}
\emph{Fixed constants.} They depend on the instance, on $a$ and on the grid, not on $\varepsilon$. Let $\thc>0$, $\ctl>0$ and $\epth>0$ be given by
Lemma~\ref{lem:theta} for $\nu=\nu_a$ (as $\m_a<\m^\star<1$), and
\begin{align*}
\ez&:=\min\Big(\frac{\m^\star-\mtwo}5,\ \epth,\ \frac{\da}{4\cst}\Big)&&\text{(largest accuracy)},\\
\ezz&:=\min\Big(\frac{\ctl}{2\kap(\m^\star+\ez)},\ \frac{\m^\star-\mtwo-\ez}{2(\mtwo+\ez)}\Big)&&\text{(gain)},\\
\dl&:=\frac{\ezz}{6(1+2\ezz)},\qquad\Cz:=\big(4+2\kap(\m^\star+\ez)\big)e^{2(1+2\ezz)},&&\\
\Co&:=\frac2{\da}+\frac{4\cst}{\da^2},\qquad\yb:=\sup_\ell\yg_\ell<\infty\quad(\text{as }\yg_\ell=\xg_\ell^2\to1).&&
\end{align*}
\emph{Constants depending on $\varepsilon\in(0,\ez]$.}
\begin{align*}
\Dp_\varepsilon&:=\Ki(\nu_a,\m^\star-2\varepsilon)-\varepsilon,\qquad
\Me_\varepsilon:=\max\Big(1,\Big\lceil\frac{\log\big((1+2\ezz)\kap(\m^\star+\varepsilon)/(2\varepsilon^2)\big)}{\log\frac1{1-\dl}}\Big\rceil\Big),\\
\Cg_\varepsilon&:=K+\frac{16\yb\,(\kap(\m^\star+\varepsilon)+1)}{(1-\dl)\varepsilon^2},\qquad
\Rz_\varepsilon:=\kap(\m^\star-2\varepsilon)-\log\big(1-e^{-\varepsilon/4}\big),\qquad\rz_\varepsilon:=\frac{\varepsilon^2}{2\Rz_\varepsilon^2},
\end{align*}
and the \emph{count threshold} $\ \up_j(\varepsilon):=\dfrac{(1+\ezz)\Lg_j}{2\varepsilon^2}$ for $j\ge1$, written $\up_j$ when $\varepsilon$ is fixed.
\end{definition}

\begin{lemma}[properties of the constants]\label{lem:constC}
\uses{def:constC, def:grid, def:hy, def:kinf, def:model, def:constA, def:empkinf}
Let $\varepsilon\in(0,\ez]$. The items are ordered by use.
\begin{enumerate}
\item[(a)] \emph{What ``$\ell$ large'' guarantees} (used in Lemmas~\ref{lem:goodhigh}, \ref{lem:fewC}, \ref{lem:optW}). If
\[
\Lg_\ell\ \ge\ \max\Big(\frac4{\dl^2},\ \frac2\dl\log\frac{1+\ezz}{2\varepsilon^2},\ \frac{\wz\kap(\m^\star+\ez)}\dl,\ 2\wz\Big),
\]
then:
\begin{itemize}
\item $1\le\up_\ell\le\tl_\ell^{1-\dl}$;
\item $(1-\dl)\Lg_\ell-\log\up_\ell\ \ge\ (1-2\dl)\Lg_\ell>0$;
\item $(1-\dl)^{-\Me_\varepsilon}\,(1-2\dl)\Lg_\ell/\kap(\m^\star+\varepsilon)\ \ge\ \up_\ell$;
\item $\wz\kap(\m^\star+\varepsilon)\le\dl\Lg_\ell$, and $\wz\le\varepsilon\up_\ell$.
\end{itemize}
\item[(b)] \emph{The final constant} (used in Lemmas~\ref{lem:decompC}, \ref{lem:pullsC}).
\[
\Dp_\varepsilon\ \ge\ \da-2\cst\varepsilon\ \ge\ \frac{\da}2\ \ge\ \frac{\varepsilon^2}8,\qquad\frac{1+\varepsilon}{\Dp_\varepsilon}\ \le\ \frac1{\da}+\Co\varepsilon .
\]
\item[(c)] \emph{Conditions of the deviation inequalities} (used in Lemma~\ref{lem:controlC}; the means also in Lemma~\ref{lem:fewC}).
\begin{itemize}
\item Means: $\varepsilon<1$, $2\varepsilon<\m^\star$, $\m_a+\varepsilon<\m^\star-3\varepsilon$, $0<\mtwo+\varepsilon<\m^\star$, $\m^\star+\varepsilon<1$.
\item Levels: $\kap(\mtwo+\varepsilon)\le\kap(\m^\star+\varepsilon)\le\kap(\m^\star+\ez)$.
\item Gains: $2\ezz(\mtwo+\varepsilon)\le\m^\star-\mtwo-\varepsilon$, and, with $q:=\m^\star-3\varepsilon$,
\[
e^{2\ezz\kap(\m^\star+\varepsilon)}\Phis_{\nu_a}(\thc,q)\le1,\qquad e^{2\ezz\kap(\m^\star+\varepsilon)}\Psis_{\nu_a}(\thc,q)\le1-\m^\star-\varepsilon .
\]
\item Exponent: $(1+2\ezz)(1-4\dl)\ge1+\ezz$.
\end{itemize}
\item[(d)] \emph{Fixed constants} (used in the proofs of (a), (b), (c) and of Lemmas~\ref{lem:fewC}, \ref{lem:optW}). $\da\ge(\m^\star-\m_a)^2/8>0$, $\ \ez>0$, $\ \ezz>0$,
$\ 0<\dl<\frac1{12}$, and $(1-\dl)(1-2\dl)\ge1-3\dl$.
\end{enumerate}
\end{lemma}
\noindent(see proof in Subsection~\ref{sec:pf-lem:constC}.)

\begin{definition}[ranks when $x=1$]\label{def:ranksC}\dt{s:ranksC}
\uses{def:constC, lem:constC, lem:levelC, def:water, def:hy, def:grid, def:high}
Let $\varepsilon\in(0,\ez]$.
\begin{itemize}
\item $\lzp(\varepsilon)$ is the first $\ell$ such that $\Lg_\ell$ satisfies the condition of Lemma~\ref{lem:constC}(a); then every $\ell\ge\lzp(\varepsilon)$ does.
\item $\ly(\varepsilon)$ is the smallest $j\ge2$ with $(1+\om_{j'-1})\yg_{j'-1}\le1+\varepsilon$ for every $j'\ge j$ (it exists since $\om\to0$ and $\yg\to1$).
\item $\lonez(\varepsilon)$ is the rank of Lemma~\ref{lem:levelC} (\emph{many high batches}) for $g=\dl$ and $m'=\Me_\varepsilon+K+2$.
\item $\lz(\varepsilon):=\max\big(\lzp(\varepsilon),\lonez(\varepsilon),\ly(\varepsilon),\lW\big)$, with $\lW$ from Subsection~\ref{sec:app-water-det}.
\end{itemize}
\end{definition}

%----------------------------------------------------------------------
\subsection{The skeleton: events, decomposition and proof of the proposition}\label{sec:app-skelC}

\begin{definition}[events and classes when $x=1$]\label{def:evC}\dt{s:evC}
\uses{def:constC, def:stacks, def:model, def:kinf, def:grid, def:hy}
Let $\varepsilon\in(0,\ez]$ and let $j\ge2$ be a batch; $n,n'$ denote integers.
\begin{itemize}
\item \emph{A deviation of a well-observed mean:} $\ \Dev_j:=\{\exists c,\ \exists n\ge\up_j:\ |\hm_{c,n}-\m_c|>\varepsilon\}$.
\item \emph{An optimal arm explored against a sub-optimal level:} $\ \Vs_j:=\{\exists n\in[1,\up_j]:\ n\Ki(\hnu_{a^\star,n},\mtwo+\varepsilon)\ge(1-4\dl)\Lg_j\}$.
\item \emph{$a$ explored, then its mean climbs:}
$\ \Va_j:=\{\exists\,1\le n\le n'\le\up_j:\ n\Ki(\hnu_{a,n},\m^\star+\varepsilon)\ge(1-4\dl)\Lg_j,\ \hm_{a,n'}\ge\m^\star-3\varepsilon\}$.
\end{itemize}
\emph{Four classes}, exactly one of which holds, defined with the statistics frozen at $\tl_{j-1}$:
\begin{align*}
\cU_j&:=\{\tms_{j-1}<\m^\star-2\varepsilon\}&&\text{(the leader is unreliable)},\\
\cG_j&:=\{\tms_{j-1}\ge\m^\star-2\varepsilon,\ \Ki(\hnu_a(\tl_{j-1}),\m^\star-2\varepsilon)>\Dp_\varepsilon\}&&\text{(good)},\\
\cC_j&:=\{\tms_{j-1}\ge\m^\star-2\varepsilon,\ \Ki(\hnu_a(\tl_{j-1}),\m^\star-2\varepsilon)\le\Dp_\varepsilon,\ \hm_a(\tl_{j-1})\ge\m^\star-3\varepsilon\}&&\text{($a$ looks competitive)},\\
\cO_j&:=\{\tms_{j-1}\ge\m^\star-2\varepsilon,\ \Ki(\hnu_a(\tl_{j-1}),\m^\star-2\varepsilon)\le\Dp_\varepsilon,\ \hm_a(\tl_{j-1})<\m^\star-3\varepsilon\}&&\text{(other)}.
\end{align*}
We write $\Dev_j(\varepsilon)$, \dots, $\cO_j(\varepsilon)$ when the accuracy must be specified.
\end{definition}

\begin{definition}[events of the consistency threshold]\label{def:amb}\dt{s:amb}
\uses{def:constC, def:hy, def:kinf, def:model, def:grid, def:stacks, def:evC, def:ranksC}
Let $\varepsilon\in(0,\ez]$ and $\ell\ge3$. Let
\[
\zet_\ell:=\gam_\ell\da(1-\m^\star)/64,\qquad\zetp_\ell:=\gam_\ell\da/32,\qquad\nl_\ell:=\Big\lceil\frac{(1-\dl)\Lg_\ell}{\kap(\m^\star+\varepsilon)+1}\Big\rceil,
\]
the cost of a \term{t:net}{net} and of the IMED phase in an activated batch, $\ \cell_\ell:=\Ug_{\ell+1}+K+16\yb\Lg_\ell/\varepsilon^2$,
and define the events
\begin{align*}
\Acc_\ell&:=\{\forall c,\ \forall n\ge\tfrac12\Ug_{\ell-1}:\ |\hm_{c,n}-\m_c|\le\zet_\ell\}&&\text{(accurate test levels)},\\
\KC_\ell&:=\{\forall n\ge\nl_\ell:\ |\Ki(\hnu_{a,n},\m^\star)-\da|\le\zetp_\ell\}&&\text{(accurate divergence of $a$)},\\
\Eamb_\ell(\varepsilon)&:=\Dev_\ell\cup\Vs_\ell\cup\Va_\ell\cup\Dev_{\ell+1}\cup\Vs_{\ell+1}\cup\Va_{\ell+1}\cup\Acc_\ell^c\cup\KC_\ell^c&&\text{(failure)}.
\end{align*}
Let $\pi(y):=\big(\frac{2K}{1-e^{-2\varepsilon^2}}+2\Cz(1+(\frac{(1+\ezz)y}{2\varepsilon^2})^4)\big)e^{-(1+\ezz)y}$. The rank $\lfour(\varepsilon)$ is the smallest $\ell\ge\lz(\varepsilon)+1$ such that
every $\ell'\ge\ell$ satisfies:
\begin{itemize}
\item[(a)] the assumptions of Lemma~\ref{lem:optW} (\emph{an optimal arm in $\Wg$}) hold at $j=\ell'$ and at $j=\ell'+1$;
\item[(b)] $\frac12\Ug_{\ell'}\,\varepsilon^2/8\ge\yb\Lg_{\ell'}$;
\item[(c)] $\nl_{\ell'}\ge e$ and $\nl_{\ell'}/\log\nl_{\ell'}\ge16\yb/(\gam_{\ell'}\varepsilon^2)$;
\item[(d)] $\wz\kap(\m^\star+\varepsilon)\le\zetp_{\ell'}\nl_{\ell'}$, $\ 2\wz/\Ug_{\ell'}\le\zet_{\ell'}\ $ and $\ 2\zet_{\ell'}\le\min(\varepsilon,\m^\star,1-\m^\star)$;
\item[(e)] $\Lg_{\ell'}\ge16$ (so $\gam_{\ell'}=\Lg_{\ell'}^{-1/4}$) and $\Lg_{\ell'-2}\ge\Lg_{\ell'}/4$;
\item[(f)] $\pi$ is non-increasing on $[\Lg_{\ell'},\infty)$.
\end{itemize}
It exists, since every condition holds for $\ell'$ large ($\Lg\to\infty$, $\xg\to1$, $\beta\ge2$, $\gam_{\ell'}^{-1}=\Lg_{\ell'}^{1/4}$, $\nl_{\ell'}\asymp\Lg_{\ell'}$).
\end{definition}

\begin{lemma}[decomposition, $x=1$]\label{lem:decompC}\dt{s:decC}
\uses{def:amb, def:constC, def:learner, def:grid, def:kinf, def:evC, def:stacks, def:model, def:hy-vocab}
Let $\varepsilon\in(0,\ez]$, $\lfive(\varepsilon):=\lfour(\varepsilon)+1$, and $n_j(a):=\Ng_a(\tl_j)-\Ng_a(\tl_{j-1})$ the pulls of $a$ in batch $j$. For every $L\ge\lfive(\varepsilon)$,
\begin{align*}
\Ng_a(\tl_L)\ \le\ &\underbrace{\frac{\Lg_L}{\da}+\Co\varepsilon\Lg_L}_{\text{main term}}+\tl_{\lfive-1}+K
+\underbrace{\sum_{j=\lfive}^L\Bg_j\big[\1(\Dev_j)+\1(\Vs_j)+\1(\Va_j)\big]}_{\text{whole batches}}\\
&+\underbrace{\Cg_\varepsilon\sum_{n=1}^{\tl_L}n\,\1\big(\Ki(\hnu_{a,n},\m^\star-2\varepsilon)\le\Dp_\varepsilon\big)}_{\text{class }\cO}
+\underbrace{\sum_{\ell\in\Act,\ \lfive-1\le\ell\le L-1}\cell_\ell\,\1\big(\Eamb_\ell(\varepsilon)\big)}_{\text{nets of an ambiguous }a},
\end{align*}
with $\cell_\ell$ from Definition~\ref{def:amb}.
\end{lemma}
\role{the skeleton of the proof of Proposition~\ref{prop:fixedC}. In each batch exactly one class holds, and each class costs something:
\begin{itemize}
\item $\cU$ or $\cC$: rare, because outside $\Dev_j$ they imply the bad events $\Vs_j$ or $\Va_j$ (Lemma~\ref{lem:fewC}, \emph{few observations}); the whole batch is paid;
\item $\cG$: $a$ is pulled at most $\approx\Lg_L/\da$ times in all, like IMED (Lemma~\ref{lem:pullsC}, \emph{pulls of $a$}): the main term;
\item $\cO$: at most $\Cg_\varepsilon$ times the count, on a deviation of the empirical divergence of $a$ (Lemmas~\ref{lem:pullsC}, \ref{lem:logcountC});
\item an \term{t:verdict}{ambiguous} $a$ in an \term{s:act}{activated} batch: rare (Lemma~\ref{lem:amb}, \emph{consistency test}); a \term{t:net}{net} is paid.
\end{itemize}}
\noindent(see proof in Subsection~\ref{sec:pf-lem:decompC}.)

\begin{lemma}[control lemmas]\label{lem:controlC}
\uses{def:constC, def:evC, def:grid, def:kinf, def:stacks, def:model, def:ranksC}
Let $\varepsilon\in(0,\ez]$ and $j\ge\lz(\varepsilon)$.
\begin{enumerate}
\item[(a)] $\PP(\Dev_j)\le\frac{2K}{1-e^{-2\varepsilon^2}}\,\tl_j^{-(1+\ezz)}$;
\item[(b)] $\PP(\Vs_j)\le\Cz\,\up_j^3\,\tl_j^{-(1+\ezz)}$;
\item[(c)] $\PP(\Va_j)\le\Cz\,\up_j^4\,\tl_j^{-(1+\ezz)}$;
\item[(d)] for every $L$, $\E\big[\Cg_\varepsilon\sum_{n=1}^{\tl_L}n\,\1(\Ki(\hnu_{a,n},\m^\star-2\varepsilon)\le\Dp_\varepsilon)\big]\le\Cg_\varepsilon/\rz_\varepsilon^2$.
\end{enumerate}
\end{lemma}
\role{bounds the expectations of the ``whole batches'' and ``class $\cO$'' terms of Lemma~\ref{lem:decompC}.}
\noindent(see proof in Subsection~\ref{sec:pf-lem:controlC}.)

\begin{lemma}[summability]\label{lem:sumC}\dt{s:sumC}
\uses{def:constC, def:grid}
Let $\varepsilon\in(0,\ez]$ and $\pi_j:=\big(\frac{2K}{1-e^{-2\varepsilon^2}}+2\Cz(1+\up_j^4)\big)\tl_j^{-(1+\ezz)}$. Then for every $L$,
\[
\sum_{j=2}^L\Bg_j\pi_j\ \le\ \Cpi(\varepsilon):=\Big(2K\Big(1+\frac1{2\varepsilon^2}\Big)+2\Cz\Big)\Big(1+\frac2\ezz\Big)+2\Cz\Big(\frac{1+\ezz}{2\varepsilon^2}\Big)^4\Big(\frac8{e\ezz}\Big)^4\Big(1+\frac2\ezz\Big).
\]
\end{lemma}
\role{sums the probabilities of Lemma~\ref{lem:controlC}(a)--(c), weighted by the batch sizes.}
\noindent(see proof in Subsection~\ref{sec:pf-lem:sumC}.)

\begin{proposition}[bound for a fixed accuracy, $x=1$]\label{prop:fixedC}\dt{s:RHo}
\uses{def:constC, def:learner, def:grid, def:kinf, def:amb, lem:decompC, lem:sumC, lem:amb}
For every $\varepsilon\in(0,\ez]$ and every $L\ge1$,
\[
\E[\Ng_a(\tl_L)]\ \le\ \frac{\Lg_L}{\da}+\Co\varepsilon\Lg_L+\RHo(\varepsilon),\qquad
\RHo(\varepsilon):=\tl_{\lfive(\varepsilon)-1}+K+\Cpi(\varepsilon)+\frac{\Cg_\varepsilon}{\rz_\varepsilon^2}+\Cfour(\varepsilon)<\infty,
\]
where $\Cfour(\varepsilon)<\infty$ is the sum of Lemma~\ref{lem:amb} (\emph{consistency test}).
\end{proposition}

\begin{proof}[Proof of Proposition~\ref{prop:fixedC}]
\proves{prop:fixedC}
\uses{lem:decompC, lem:controlC, lem:sumC, lem:amb, def:learner, def:grid, def:constC, def:evC, def:amb}
\stp1{small $L$}If $L<\lfive(\varepsilon)$, then $\Ng_a(\tl_L)\le\tl_L\le\tl_{\lfive-1}$.

\stp2{large $L$}Otherwise take expectations in Lemma~\ref{lem:decompC} (\emph{decomposition}). Each term is bounded independently of $L$:
\begin{itemize}
\item whole batches: at most $\sum_j\Bg_j\pi_j\le\Cpi(\varepsilon)$ (Lemma~\ref{lem:controlC}(a)--(c), then $\up_j^3+\up_j^4\le2(1+\up_j^4)$, then Lemma~\ref{lem:sumC});
\item class $\cO$: at most $\Cg_\varepsilon/\rz_\varepsilon^2$ (Lemma~\ref{lem:controlC}(d));
\item nets of an ambiguous $a$: at most $\Cfour(\varepsilon)$ (Lemma~\ref{lem:amb}, \emph{consistency test}).
\end{itemize}
The sum of these bounds, $\tl_{\lfive-1}$ and $K$ is $\RHo(\varepsilon)$.
\end{proof}

\noindent The rest of the section proves the lemmas used by the decomposition, grouped by the question that they answer.

%----------------------------------------------------------------------
\subsection{Why are many batches high?}\label{sec:app-water}

Recall (W1), (W2) and the \term{s:water}{water level} $\Lam_j$ from Definition~\ref{def:water}.
\begin{definition}[high batches]\label{def:high}\dt{t:high}
\uses{def:water, def:grid}
For $g\in(0,1)$ and $L>j$, batch $j$ is \emph{$(g,L)$-high} if
\[
\Lam_j\ \ge\ (1-g)\,\Lg_L .
\]
\end{definition}
\noindent In words: at the end of batch $j$, every arm has been explored almost up to the level $\log\tl_L$ of the later horizon $\tl_L$, either by many observations or by a strong
evidence $N\Kt$ against it. Two consequences are used later.
\begin{itemize}
\item The \term{t:leader}{leader} of a high batch has at least $\tl_L^{1-g}$ observations, by (W1). So, in all high batches except at most $K$ (the set $\Sgo$ of Lemma~\ref{lem:goodhigh}), the leader was already well observed
at the start of the batch, and its mean is accurate.
\item An arm with few observations has been ``refuted'' at the level $\approx\log\tl_L$ in every high batch; for an optimal arm this is a rare event, $\Vs$
(Lemmas~\ref{lem:highC}, \ref{lem:fewC}).
\end{itemize}
Being high is a deterministic property of the run, not a good event: the next three lemmas show that, on every run, many batches are high.

\begin{lemma}[low batches carry few pulls]\label{lem:lowmass}
\uses{def:hy, def:grid, def:water}
For every run of \HY, every $\ell>\lW$ and every $v\ge0$, there is a set $E$ of at most $K$ batches such that
\[
\sum_{\lW\le i<\ell,\ \Lam_i<v,\ i\notin E}\Bg_i\ \le\ Ke^v .
\]
\end{lemma}
\role{used in Lemma~\ref{lem:levelC}: batches with a low water level carry few pulls.}
\noindent(see proof in Subsection~\ref{sec:pf-lem:lowmass}.)

\begin{lemma}[growth of the grid]\label{lem:growth}
\uses{def:grid}
Let $k_0\ge1$ and $\beta_0\ge0$ with $\xg_j\le1+\beta_0$ for $j>k_0$, let $S$ be a finite set of batches, $n\ge k_0$, and $W'\ge\sum_{k_0<j\le n,\ j\notin S}\Bg_j$. If $h'$
batches of $(k_0,n]$ belong to $S$, then
\[
\log\tl_n\ \le\ (1+\beta_0)^{h'}\big(\log(\tl_{k_0}+W')+h'\log2\big).
\]
\end{lemma}
\role{used in Lemma~\ref{lem:levelC}: if few batches carry the mass, the grid cannot grow enough.}
\noindent(see proof in Subsection~\ref{sec:pf-lem:growth}.)

\begin{lemma}[many high batches]\label{lem:levelC}
\uses{def:hy, def:high}
For every $g\in(0,1)$ and every integer $m'\ge0$ there is a rank $\lonez$, depending only on $g$, $m'$, $K$, $\beta$ and the grid, such that, for every run of \HY{} and
every $L\ge\lonez$, at least $m'$ batches of $\{2,\dots,L-1\}$ are $(g,L)$-high.
\end{lemma}
\role{provides the high batches used by Lemma~\ref{lem:goodhigh}, hence by Lemmas~\ref{lem:fewC}, \ref{lem:logcountC} and~\ref{lem:optW}.}
\noindent(see proof in Subsection~\ref{sec:pf-lem:levelC}.)

%----------------------------------------------------------------------
\subsection{Why are the leaders reliable, and why are wrong means rare?}\label{sec:app-pathC}

\begin{lemma}[good high batches]\label{lem:goodhigh}\dt{s:good}
\uses{def:constC, def:learner, def:hy, def:grid, def:evC, def:model, def:high, def:ranksC}
Let $\varepsilon\in(0,\ez]$ and $j\ge\lz(\varepsilon)$. Let $\Shi_j$ be the set of the $(\dl,j)$-\term{t:high}{high} batches of $\{2,\dots,j-1\}$, and
$\Sgo_j:=\{i\in\Shi_j:\ \Ng_{\Ld_{i-1}}(\tl_{i-1})\ge\up_j\}$ the high batches whose leader was already well observed. Then:
\begin{itemize}
\item $|\Shi_j|\ge\Me_\varepsilon+K+2$, and every $i\in\Shi_j$ has $\Ng_{\Ld_{i-1}}(\tl_i)\ge\tl_j^{1-\dl}\ge\up_j$;
\item $|\Sgo_j|\ge\Me_\varepsilon+2$;
\item outside $\Dev_j$, $\tms_{i-1}\le\m_{\Ld_{i-1}}+\varepsilon$ for every $i\in\Sgo_j$.
\end{itemize}
\end{lemma}
\role{used in Lemmas~\ref{lem:fewC}, \ref{lem:logcountC}, \ref{lem:optW}: in most high batches, the leader is well observed, hence accurate.}
\noindent(see proof in Subsection~\ref{sec:pf-lem:goodhigh}.)

\begin{lemma}[evidence in a high batch]\label{lem:highC}
\uses{def:hy, def:learner, def:grid, def:kinf, def:high, def:water}
Let $i\ge2$, $v\le\Lam_i$, $r'\in[0,1)$ with $\tms_{i-1}\le r'$, and an arm $c$ with $\Ng_c(\tl_i)\le\bar N$ and $Y:=v-\log\bar N>0$. Let $D:=\Ki(\tnu_c(\tl_{i-1}),r')$. Then
\[
0<D\le\kap(r')\qquad\text{and}\qquad\Ng_c(\tl_i)\,D\ \ge\ Y .
\]
\end{lemma}
\role{used in Lemmas~\ref{lem:fewC}, \ref{lem:optW}: in a high batch, an arm with few observations carries a large evidence against the level $r'$.}
\noindent(see proof in Subsection~\ref{sec:pf-lem:highC}.)

\begin{lemma}[few observations]\label{lem:fewC}
\uses{def:constC, def:evC, def:stacks, def:grid, def:model, def:ranksC}
Let $\varepsilon\in(0,\ez]$ and $j\ge\lz(\varepsilon)$. Then
\[
\Dev_j^c\cap\cU_j\ \subset\ \Vs_j\qquad\text{and}\qquad\Dev_j^c\cap\{\hm_a(\tl_{j-1})\ge\m^\star-3\varepsilon\}\ \subset\ \Va_j .
\]
\end{lemma}
\role{used in Lemma~\ref{lem:decompC} (the \term{s:evC}{classes} $\cU$ and $\cC$ are rare) and in Lemma~\ref{lem:amb} (\emph{consistency test}).}
\noindent(see proof in Subsection~\ref{sec:pf-lem:fewC}.)

\begin{lemma}[logarithmic count]\label{lem:logcountC}
\uses{def:constC, def:evC, def:grid, def:kinf, def:model, def:learner, def:ranksC}
Let $\varepsilon\in(0,\ez]$ and $j\ge\lz(\varepsilon)$. Outside $\Dev_j$, every arm $c$ satisfies
\[
(1-\dl)\Lg_j\ \le\ \big(\kap(\m^\star+\varepsilon)+1\big)\,\Ng_c(\tl_{j-1}).
\]
\end{lemma}
\role{used in Lemma~\ref{lem:decompC} (class $\cO$: the pulls of a batch are at most $\Cg_\varepsilon$ times the count) and in Lemma~\ref{lem:amb}, \emph{consistency test}: counts at least $\nl_\ell$.}
\noindent(see proof in Subsection~\ref{sec:pf-lem:logcountC}.)

%----------------------------------------------------------------------
\subsection{Why is an optimal arm well sampled, and how many times is \texorpdfstring{$a$}{a} pulled?}\label{sec:app-optW}

\begin{lemma}[an optimal arm in $\Wg$]\label{lem:optW}
\uses{def:constC, def:hy, def:grid, def:evC, def:model, def:ranksC}
Let $\varepsilon\in(0,\ez]$ and $j\ge\lz(\varepsilon)$ with
\[
\tfrac12\Ug_{j-1}\le\tl_j^{1-\dl},\qquad\up_j\le\tfrac12\Ug_{j-1},\qquad2\wz\le\varepsilon\Ug_{j-1},\qquad\tl_{j-1}\ge K\Ug_{j-1}.
\]
Then, on $\Dev_j^c\cap(\Vs_j)^c$:
\begin{itemize}
\item some optimal arm is \term{t:ref}{well-sampled} at $\tl_{j-1}$ (it lies in $\Wg_{j-1}$), and the \term{t:ref}{reference} $\rr_{j-1}$ is optimal;
\item the \term{t:ref}{test level} is accurate: $\m^\star-2\varepsilon\le\mh_{j-1}\le\m^\star+\varepsilon$;
\item every sub-optimal $c\in\Wg_{j-1}$ has $\tm_c(\tl_{j-1})<\m^\star-2\varepsilon$.
\end{itemize}
\end{lemma}
\role{used in Lemma~\ref{lem:pullsC}(i) ($a$ is not the reference, and the test level is at least $\m^\star-2\varepsilon$) and in Lemma~\ref{lem:amb}, \emph{consistency test}: the test levels are accurate.}
\noindent The last assumption, $\tl_{j-1}\ge K\Ug_{j-1}$, is not used in the proof of this lemma; it is listed here so that (a) of
Definition~\ref{def:amb} provides it to the proof of Lemma~\ref{lem:amb} (Case~(i), $\Wg_{\ell-1}\ne\emptyset$).
\noindent(see proof in Subsection~\ref{sec:pf-lem:optW}.)

\begin{lemma}[pulls of $a$ when $x=1$]\label{lem:pullsC}
\uses{def:constC, def:evC, def:kinf, def:stacks, def:grid, def:model, def:hy, def:learner, def:amb, def:hy-vocab}
Let $\varepsilon\in(0,\ez]$ and let $j$ satisfy the assumptions of Lemma~\ref{lem:optW}. Assume that the outcome is in $(\cG_j\cup\cO_j)\cap\Dev_j^c\cap(\Vs_j)^c$, and let
$D:=\Ki(\hnu_a(\tl_{j-1}),\m^\star-2\varepsilon)$. Then:
\begin{enumerate}
\item[(i)] $a$ is neither the leader $\Ld_{j-1}$ nor the reference $\rr_{j-1}$, and $\Kt_{j-1}(a)\ge\Kh_{j-1}(a)\ge D\ge\varepsilon^2/8$. If moreover the outcome is in the good class $\cG_j$
(the mean of the leader is at least $\m^\star-2\varepsilon$ and the empirical divergence of $a$ is above $\Dp_\varepsilon$), then $D>\Dp_\varepsilon$.
\item[(ii)] If batch $j$ is not \term{s:act}{activated}, or $a$ is \term{t:verdict}{rejected} at $\tl_{j-1}$, and if $a$ is pulled in batch $j$, then
\[
\Ng_a(\tl_j)\ \le\ K+\frac{(1+\om_{j-1})\,\yg_{j-1}\Lg_j}{D}.
\]
\item[(iii)] If batch $j$ is activated and $a$ is \term{t:verdict}{ambiguous} at $\tl_{j-1}$, then $\Ng_a(\tl_j)-\Ng_a(\tl_{j-1})\le\cell_{j-1}$.
\end{enumerate}
\end{lemma}
\role{used in Lemma~\ref{lem:decompC}: (i) shows that $a$ is not the previous leader, so that a non-rejected $a$ is ambiguous; (ii) gives the main term on $\cG$ and the factor $\Cg_\varepsilon$ on $\cO$; (iii) gives the net of an ambiguous $a$.}
\noindent(see proof in Subsection~\ref{sec:pf-lem:pullsC}.)

%----------------------------------------------------------------------
\subsection{Why does \texorpdfstring{$a$}{a} pass the consistency test?}\label{sec:app-amb}

\begin{lemma}[a sub-optimal arm passes the consistency test]\label{lem:amb}
\uses{def:constC, def:amb, def:hy, def:grid, def:hy-vocab}
Let $\varepsilon\in(0,\ez]$ and $\ell\ge\lfour(\varepsilon)$. Then, on $\Eamb_\ell(\varepsilon)^c$: $a\notin\{\Ld_\ell,\rr_\ell\}$, $\Kh_\ell(a)\ge\varepsilon^2/8$, and $a$ is
\term{t:verdict}{rejected} at $\tl_\ell$. Moreover,
\[
\Cfour(\varepsilon):=\sum_{\ell\in\Act,\ \ell\ge\lfour(\varepsilon)}\cell_\ell\,\PP\big(\Eamb_\ell(\varepsilon)\big)\ <\ \infty .
\]
\end{lemma}
\role{used in Lemma~\ref{lem:decompC} and Proposition~\ref{prop:fixedC}: an ambiguous $a$ in an activated batch happens only on $\Eamb$, and the nets paid there sum to
$\Cfour(\varepsilon)<\infty$.}
\noindent(see proof in Subsection~\ref{sec:pf-lem:amb}.)

%======================================================================
\subsection{Proof of Lemma~\ref{lem:constC}}\label{sec:pf-lem:constC}
\noindent We prove the items in the order (d), (c), (b), (a), since each uses the previous ones. Recall $\dl:=\ezz/(6(1+2\ezz))$ and
$\cst=\frac{2-\m^\star}{1-\m^\star}=1+\frac1{1-\m^\star}$.
\begin{proof}
\proves{lem:constC}
\uses{lem:kinf, lem:theta, def:kinf, def:model, def:constC, def:constA, def:grid, def:hy, def:empkinf}
\stp{(d)}{fixed constants}
\begin{itemize}
\item $\da\ge(\m^\star-\m_a)^2/8>0$. \ (Lemma~\ref{lem:kinf}(iv) with $p=q=\m_a<\mu=\m^\star$.)
\item $\ez>0$. \ ($\m^\star>\mtwo$, $\epth>0$, $\da>0$ and $\cst>0$.)
\item $\ezz>0$. \ (Its first term is positive, and $\m^\star-\mtwo-\ez\ge\frac45(\m^\star-\mtwo)>0$ since $\ez\le\frac{\m^\star-\mtwo}5$.)
\item $0<\dl<\frac1{12}$. \ ($\dl=\frac16\cdot\frac{\ezz}{1+2\ezz}$, $\ezz>0$ and $\frac{\ezz}{1+2\ezz}<\frac12$.)
\item $(1-\dl)(1-2\dl)\ge1-3\dl$. \ (The difference is $2\dl^2\ge0$.)
\end{itemize}

\stp{(c)}{conditions of the deviation inequalities}
\begin{itemize}
\item $2\varepsilon<\m^\star$, $\m_a+\varepsilon<\m^\star-3\varepsilon$ and $0<\mtwo+\varepsilon<\m^\star$. \ ($5\varepsilon\le5\ez\le\m^\star-\mtwo\le\m^\star-\m_a$.)
\item $\varepsilon<1$ and $\m^\star+\varepsilon<1$. \ ($\varepsilon\le\epth<1-\m^\star$ by Lemma~\ref{lem:theta}.)
\item The order of the levels. \ ($\kap$ is non-decreasing, and $\mtwo+\varepsilon\le\m^\star+\varepsilon\le\m^\star+\ez$.)
\item $2\ezz(\mtwo+\varepsilon)\le\m^\star-\mtwo-\varepsilon$. \ (The map $\varepsilon\mapsto2\ezz(\mtwo+\varepsilon)+\mtwo+\varepsilon-\m^\star$ is increasing, and non-positive at $\ez$
by the second term in the definition of $\ezz$.)
\item The two inequalities on $\Phis$ and $\Psis$. \ (Lemma~\ref{lem:theta} with $\varepsilon\le\epth$ gives them with the factor $e^{\ctl}$, and
$2\ezz\kap(\m^\star+\varepsilon)\le2\ezz\kap(\m^\star+\ez)\le\ctl$ by the first term in the definition of $\ezz$.)
\item $(1+2\ezz)(1-4\dl)=1+\frac43\ezz\ge1+\ezz$. \ (By definition of $\dl$, $4\dl(1+2\ezz)=\frac23\ezz$, so $(1+2\ezz)(1-4\dl)=1+2\ezz-\frac23\ezz$.)
\end{itemize}

\stp{(b)}{the final constant}\emph{Lower bound on $\Dp_\varepsilon$.} Lemma~\ref{lem:kinf}(v) with $\mu=\m^\star$ and the shift $2\varepsilon\le\m^\star$ gives
\begin{align*}
\Dp_\varepsilon\ &\ge\ \da-\frac{2\varepsilon}{1-\m^\star}-\varepsilon\ \ge\ \da-2\cst\varepsilon&&\why{$2\cst\varepsilon=2\varepsilon+\tfrac{2\varepsilon}{1-\m^\star}$}\\
&\ge\ \frac{\da}2&&\why{$4\cst\varepsilon\le4\cst\ez\le\da$}\\
&\ge\ \frac{(5\varepsilon)^2}{16}\ \ge\ \frac{\varepsilon^2}8&&\why{(d): $\da\ge(\m^\star-\m_a)^2/8\ge(5\varepsilon)^2/8$}.
\end{align*}
\emph{Upper bound on $(1+\varepsilon)/\Dp_\varepsilon$.} First,
\[
\frac1{\Dp_\varepsilon}-\frac1{\da}\ =\ \frac{\da-\Dp_\varepsilon}{\Dp_\varepsilon\,\da}\ \le\ \frac{2\cst\varepsilon}{(\da/2)\,\da}\ =\ \frac{4\cst\varepsilon}{\da^2}.
\]
Then, expanding the product and using $\frac{4\cst\varepsilon^2}{\da^2}\le\frac\varepsilon{\da}$ (as $4\cst\varepsilon\le\da$),
\[
\frac{1+\varepsilon}{\Dp_\varepsilon}\ \le\ (1+\varepsilon)\Big(\frac1{\da}+\frac{4\cst\varepsilon}{\da^2}\Big)\ \le\ \frac1{\da}+\varepsilon\Big(\frac2{\da}+\frac{4\cst}{\da^2}\Big)\ =\ \frac1{\da}+\Co\varepsilon .
\]

\stp{(a)}{what ``$\ell$ large'' guarantees}Let $y:=\Lg_\ell$, which satisfies the four lower bounds of (a).
\begin{itemize}
\item \emph{$\log\up_\ell\le\dl y$.} For $y\ge4/\dl^2$, $\log y\le\sqrt y\le\frac\dl2y$. With $y\ge\frac2\dl\log\frac{1+\ezz}{2\varepsilon^2}$,
\[
\log\up_\ell\ =\ \log\frac{1+\ezz}{2\varepsilon^2}+\log y\ \le\ \frac\dl2y+\frac\dl2y\ =\ \dl y .
\]
\item \emph{$1\le\up_\ell\le\tl_\ell^{1-\dl}$.} $\up_\ell\ge y/2\ge1$ (as $\varepsilon<1$ and $y\ge4/\dl^2\ge2$), and $\up_\ell\le e^{\dl y}=\tl_\ell^{\dl}\le\tl_\ell^{1-\dl}$ (as $\dl<\frac12$).
\item \emph{$(1-\dl)y-\log\up_\ell\ge(1-2\dl)y>0$.} By the first point.
\item \emph{$(1-\dl)^{-\Me_\varepsilon}(1-2\dl)y/\kap(\m^\star+\varepsilon)\ge\up_\ell$.} By definition of $\Me_\varepsilon$,
$(1-\dl)^{-\Me_\varepsilon}\ge(1+2\ezz)\kap(\m^\star+\varepsilon)/(2\varepsilon^2)$ (if the logarithm in its definition is negative, the right side is below $1$). Hence
\[
(1-\dl)^{-\Me_\varepsilon}\frac{(1-2\dl)y}{\kap(\m^\star+\varepsilon)}\ \ge\ \frac{(1+2\ezz)(1-2\dl)y}{2\varepsilon^2}\ \ge\ \frac{(1+\ezz)y}{2\varepsilon^2}\ =\ \up_\ell,
\]
since $(1+2\ezz)(1-2\dl)\ge(1+2\ezz)(1-4\dl)\ge1+\ezz$ by (c).
\item \emph{$\wz\kap(\m^\star+\varepsilon)\le\dl y$.} $\wz\kap(\m^\star+\varepsilon)\le\wz\kap(\m^\star+\ez)\le\dl y$ by the third lower bound on $y$.
\item \emph{$\wz\le\varepsilon\up_\ell$.} $\varepsilon\up_\ell=\frac{(1+\ezz)y}{2\varepsilon}\ge\frac y2\ge\wz$ by the fourth lower bound on $y$.
\end{itemize}
\end{proof}

\subsection{Proof of Lemma~\ref{lem:decompC}}\label{sec:pf-lem:decompC}
\noindent\emph{Idea.} In each batch, exactly one of the four classes holds. The classes $\cU$ and $\cC$, and the deviations, cost a whole batch on rare events. On $\cG$
the pulls of $a$ are bounded by the IMED inequality or by the paced plan (Lemma~\ref{lem:pullsC}(ii)), and counted once over the batches. On $\cO$ they are bounded by
$\Cg_\varepsilon$ times the count, and summed over distinct counts. The only new cost is the net of an ambiguous $a$ in an activated batch, charged to $\Eamb$.

\par\smallskip
\Setting $\varepsilon\in(0,\ez]$, $L\ge\lfive:=\lfive(\varepsilon)$, and $G'_j:=\Dev_j^c\cap(\Vs_j)^c$.
\Goal the bound of the lemma, by splitting $\Ng_a(\tl_L)=\Ng_a(\tl_{\lfive-1})+\sum_{j=\lfive}^Ln_j(a)$ according to the classes.
\par\smallskip

\begin{proof}
\proves{lem:decompC}
\uses{lem:fewC, lem:pullsC, def:amb, lem:amb, lem:constC, lem:distinct, lem:logcountC, def:evC, def:grid, def:learner, def:hy-vocab, def:hy, def:constC, def:kinf, def:model, def:stacks, def:ranksC}
\stp1{splitting by classes}For every batch $j\ge\lfive$ exactly one of $\cU_j,\cG_j,\cC_j,\cO_j$ holds, and $\Dev_j^c\cap\cU_j\subset\Vs_j$, $\Dev_j^c\cap\cC_j\subset\Va_j$
(Lemma~\ref{lem:fewC}). So outside $\Dev_j\cup\Vs_j\cup\Va_j$ the class is $\cG_j$ or $\cO_j$. With $n_j(a)\le\Bg_j$ and $\Ng_a(\tl_{\lfive-1})\le\tl_{\lfive-1}$,
\[
\Ng_a(\tl_L)\ \le\ \tl_{\lfive-1}+\sum_{j=\lfive}^L\Bg_j\big[\1(\Dev_j)+\1(\Vs_j)+\1(\Va_j)\big]+\underbrace{\sum_{j=\lfive}^Ln_j(a)\1(\cG_j\cap G'_j)}_{\Sigma_{\mathcal G}}
+\underbrace{\sum_{j=\lfive}^Ln_j(a)\1(\cO_j\cap G'_j)}_{\Sigma_{\mathcal O}} .
\]

\stp2{splitting by the verdict}Let $P^1_j:=\{j-1\notin\Act\text{, or }a\text{ is rejected at }\tl_{j-1}\}$ and $P^2_j:=(P^1_j)^c$. On $(\cG_j\cup\cO_j)\cap G'_j$,
Lemma~\ref{lem:pullsC}(i) (whose assumptions hold by (a) of Definition~\ref{def:amb}) gives $a\ne\Ld_{j-1}$. So there $P^2_j$ means: $j-1\in\Act$ and $a$ is ambiguous at
$\tl_{j-1}$. Then:
\begin{itemize}
\item Lemma~\ref{lem:amb} (\emph{consistency test}), applied at $j-1\ge\lfour$, puts the outcome in $\Eamb_{j-1}(\varepsilon)$;
\item Lemma~\ref{lem:pullsC}(iii) gives $n_j(a)\le\cell_{j-1}$.
\end{itemize}
So the $P^2$ parts of $\Sigma_{\mathcal G}$ and $\Sigma_{\mathcal O}$ are at most the last sum of the lemma.

\stp3{the $P^1$ part of $\Sigma_{\mathcal G}$}If $a$ is pulled in such a batch $j$, Lemma~\ref{lem:pullsC}(ii) with $D>\Dp_\varepsilon$, $(1+\om_{j-1})\yg_{j-1}\le1+\varepsilon$
($j\ge\ly(\varepsilon)$) and $\Lg_j\le\Lg_L$ gives
\[
\Ng_a(\tl_j)\ \le\ K+\frac{(1+\varepsilon)\Lg_L}{\Dp_\varepsilon}\ \le\ K+\frac{\Lg_L}{\da}+\Co\varepsilon\Lg_L\qquad\why{Lemma~\ref{lem:constC}(b)}.
\]
By Lemma~\ref{lem:distinct}(ii), the $P^1$ part of $\Sigma_{\mathcal G}$ is at most this bound: this is the main term, plus $K$.

\stp4{the $P^1$ part of $\Sigma_{\mathcal O}$}If $a$ is pulled in such a batch $j$, Lemma~\ref{lem:pullsC}(ii) with $D\ge\varepsilon^2/8$, $1+\om\le2$ and $\yg\le\yb$, then
Lemma~\ref{lem:logcountC}, give
\[
n_j(a)\ \le\ \Ng_a(\tl_j)\ \le\ K+\frac{16\yb\Lg_j}{\varepsilon^2}\ \le\ \Big(K+\frac{16\yb(\kap(\m^\star+\varepsilon)+1)}{(1-\dl)\varepsilon^2}\Big)\Ng_a(\tl_{j-1})\ =\ \Cg_\varepsilon\Ng_a(\tl_{j-1}).
\]
On $\cO_j$, $\Ki(\hnu_{a,\Ng_a(\tl_{j-1})},\m^\star-2\varepsilon)\le\Dp_\varepsilon$. By Lemma~\ref{lem:distinct}(i), whose hypothesis $\Ng_a(\tl_1)\ge1$ holds for \HY{} (Definition~\ref{def:hy-vocab}), this part is at most the
$\Cg_\varepsilon$-sum of the lemma.
\end{proof}

%----------------------------------------------------------------------
\subsection{Proof of Lemma~\ref{lem:controlC}}\label{sec:pf-lem:controlC}
\begin{proof}
\proves{lem:controlC}
\uses{lem:hoeffding, lem:single, lem:constC, lem:pair, lem:lowerdev, def:evC, def:constC, def:grid, def:model, def:kinf, def:stacks}
\stp{(a)}{$\Dev_j$}Lemma~\ref{lem:hoeffding}(b) for each arm with $\rho=\varepsilon$ and $u=\up_j$, a union bound over the $K$ arms, and $2\varepsilon^2\up_j=(1+\ezz)\Lg_j$.

\stp{(b)}{$\Vs_j$}A union bound over the at most $\up_j$ integers $n\in[1,\up_j]$, and Lemma~\ref{lem:single} with $\nu=\nu_{a^\star}$, $m=\m^\star$, $r=\mtwo+\varepsilon$, $\eta=\ezz$
(its conditions are Lemma~\ref{lem:constC}(c)) and $z=(1-4\dl)\Lg_j$. Its constant is at most $\Cz$, and $(1+2\ezz)(1-4\dl)\ge1+\ezz$ gives the exponent.

\stp{(c)}{$\Va_j$}A union bound over the at most $\up_j^2$ pairs, and Lemma~\ref{lem:pair} with $\nu=\nu_a$, $r=\m^\star+\varepsilon$, $q=\m^\star-3\varepsilon$, $\theta=\thc$, $\eta=\ezz$
(its conditions are Lemma~\ref{lem:constC}(c)).

\stp{(d)}{the class $\cO$}Lemma~\ref{lem:lowerdev} with $\mu=\m^\star-2\varepsilon$ and $\varepsilon'=\varepsilon$, whose $\rho'$ is $\rz_\varepsilon$, gives
$\PP(\Ki(\hnu_{a,n},\m^\star-2\varepsilon)\le\Dp_\varepsilon)\le e^{-\rz_\varepsilon n}$; and $\sum_{n\ge1}ne^{-\rho n}=\frac{e^{-\rho}}{(1-e^{-\rho})^2}\le\frac1{\rho^2}$, as $(1-e^{-\rho})^2\ge e^{-\rho}\rho^2$.
\end{proof}

%----------------------------------------------------------------------
\subsection{Proof of Lemma~\ref{lem:sumC}}\label{sec:pf-lem:sumC}
\begin{proof}
\proves{lem:sumC}
\uses{lem:sums, def:constC, def:grid}
Use $\frac1{1-e^{-2\varepsilon^2}}\le1+\frac1{2\varepsilon^2}$, $\up_j^4=(\frac{1+\ezz}{2\varepsilon^2})^4\Lg_j^4$, and Lemma~\ref{lem:sums}(vi) with $\theta=\ezz$ for the two sums
$\sum_j\Bg_j\tl_j^{-(1+\ezz)}$ and $\sum_j\Bg_j\Lg_j^4\tl_j^{-(1+\ezz)}$.
\end{proof}

%----------------------------------------------------------------------
\subsection{Proof of Lemma~\ref{lem:lowmass}}\label{sec:pf-lem:lowmass}
\noindent\emph{Idea.} By (W2), an arm pulled in a batch of level $<v$ had fewer than $e^v$ observations at its start. So each arm is pulled in such batches only while it
has fewer than $e^v$ observations, except during its last such batch.

\begin{proof}
\proves{lem:lowmass}
\uses{lem:W, def:learner, def:grid, def:water}
\stp1{the exceptional set}Let $\mathrm{Lo}$ be the set of the batches $i\in[\lW,\ell)$ with $\Lam_i<v$. For each arm $c$ pulled in some batch of $\mathrm{Lo}$, let $e(c)$ be the
last one, and $E:=\{e(c)\}$, a set of at most $K$ batches.

\stp2{pulls of one arm}By (W2) at $e(c)\ge\lW$ (Lemma~\ref{lem:W}, \emph{the IMED phase lifts the water level}), $\log\Ng_c(\tl_{e(c)-1})\le\Lam_{e(c)}<v$. The batches of $\mathrm{Lo}\setminus E$ in which $c$ is pulled lie in
$[\lW,e(c))$, so $c$ is pulled in them at most $\Ng_c(\tl_{e(c)-1})-\Ng_c(\tl_{\lW-1})<e^v$ times.

\stp3{conclusion}$\sum_{i\in\mathrm{Lo}\setminus E}\Bg_i$ is the number of pulls in these batches; summing Step~2 over the at most $K$ arms gives $Ke^v$.
\end{proof}

%----------------------------------------------------------------------
\subsection{Proof of Lemma~\ref{lem:growth}}\label{sec:pf-lem:growth}
\par\smallskip
\Setting for $n'\in[k_0,n]$, let $w(n')$ be the mass of the batches of $(k_0,n']\setminus S$, $h(n')$ the number of batches of $(k_0,n']$ in $S$, and $A:=\log(\tl_{k_0}+W')\ge0$.
\Goal by induction on $n'$: $\ \log\big(\tl_{n'}+W'-w(n')\big)\le(1+\beta_0)^{h(n')}\big(A+h(n')\log2\big)$. At $n'=n$, $\log\tl_n\le\log(\tl_n+W'-w(n))$ since $W'\ge w(n)$, which
gives the lemma.
\par\smallskip
\begin{proof}
\proves{lem:growth}
\uses{def:grid}
The case $n'=k_0$ is an equality. Let $Z:=\tl_{n'}+W'-w(n')\ge1$.
\begin{itemize}
\item If $n'+1\notin S$: $\tl_{n'+1}=\tl_{n'}+\Bg_{n'+1}$ and $w(n'+1)=w(n')+\Bg_{n'+1}$, so the quantity is unchanged, and so is $h$.
\item If $n'+1\in S$: $\tl_{n'+1}=\tl_{n'}^{\xg_{n'+1}}\le Z^{1+\beta_0}$ and $W'-w(n'+1)=W'-w(n')\le Z\le Z^{1+\beta_0}$. So
\begin{align*}
\log\big(\tl_{n'+1}+W'-w(n'+1)\big)\ &\le\ \log2+(1+\beta_0)\log Z\ \le\ \log2+(1+\beta_0)^{h+1}(A+h\log2)\\
&\le\ (1+\beta_0)^{h+1}\big(A+(h+1)\log2\big),
\end{align*}
with $h:=h(n')$ and $h(n'+1)=h+1$.
\end{itemize}
\end{proof}

%----------------------------------------------------------------------
\subsection{Proof of Lemma~\ref{lem:levelC}}\label{sec:pf-lem:levelC}
\noindent\emph{Idea.} If fewer than $m'$ batches were high, all the other batches except at most $K$ would carry few pulls (Lemma~\ref{lem:lowmass}). So $\tl$ could only
grow through at most $K+m'$ batches, each raising $\log\tl$ by a factor close to $1$: $\log\tl_{L-1}\le(1-g/2)\log\tl_L+O(1)$. But $\xg_L\to1$.

\par\smallskip
\Setting the following quantities, which depend only on $g$, $m'$, $K$, $\beta$ and the grid.
\begin{itemize}
\item $\bar h:=K+m'$, the number of exceptional batches allowed.
\item $\beta_0:=\log\frac{1-g/2}{1-g}\big/(\bar h+1)>0$, chosen so that $(1+\beta_0)^{\bar h}(1-g)\le1-g/2$.
\item A rank $k_0\ge\lW$ with $\xg_j\le1+\beta_0$ for every $j>k_0$.
\item $\omega':=(1+\beta_0)^{\bar h}\big(\log(\tl_{k_0}+K)+\bar h\log2\big)$, a constant.
\item A rank $\lonez\ge k_0+2$ such that every $L\ge\lonez$ has $\xg_L\le(1-g/4)^{-1}$ and $\Lg_L>4\omega'/g$.
\end{itemize}
\Goal for $L\ge\lonez$, the set $\mathrm{Hi}$ of the $(g,L)$-high batches of $[2,L)$ has at least $m'$ elements.
\par\smallskip
\begin{proof}
\proves{lem:levelC}
\uses{lem:lowmass, lem:growth, def:grid, def:water, def:high}
\stp1{assume the contrary}Let $L\ge\lonez$, $v:=(1-g)\Lg_L$, and assume $|\mathrm{Hi}|<m'$. Let $E$ be given by Lemma~\ref{lem:lowmass} with $\ell:=L$, and $S:=\mathrm{Hi}\cup E$, so
$|S|\le\bar h$.

\stp2{little mass outside $S$}A batch $j\in(k_0,L-1]\setminus S$ satisfies $\lW\le j<L$, $\Lam_j<v$ (it is not high) and $j\notin E$. So $W':=Ke^v$ bounds the mass of
$(k_0,L-1]\setminus S$.

\stp3{the grid cannot grow enough}Lemma~\ref{lem:growth} with $n=L-1$ (its right-hand side is non-decreasing in $h'$), and $\tl_{k_0}+Ke^v\le(\tl_{k_0}+K)e^v$, give
\begin{align*}
\log\tl_{L-1}\ &\le\ (1+\beta_0)^{\bar h}\big(\log(\tl_{k_0}+K)+v+\bar h\log2\big)\\
&=\ (1+\beta_0)^{\bar h}(1-g)\Lg_L+\omega'\ \le\ (1-g/2)\Lg_L+\omega'&&\why{choice of $\beta_0$}.
\end{align*}
But $\log\tl_{L-1}=\Lg_L/\xg_L\ge(1-g/4)\Lg_L$. So $(g/4)\Lg_L\le\omega'$, which contradicts $\Lg_L>4\omega'/g$.
\end{proof}

%----------------------------------------------------------------------
\subsection{Proof of Lemma~\ref{lem:goodhigh}}\label{sec:pf-lem:goodhigh}
\noindent\emph{Idea.} The leader of a high batch has many observations at the end of its batch (W1). If it had few at the start, it cannot be the poorly observed leader
of a later high batch.

\begin{proof}
\proves{lem:goodhigh}
\uses{lem:levelC, lem:constC, lem:pseudo, def:constC, def:learner, def:hy, def:grid, def:evC, def:stacks, def:model, def:high, def:ranksC, def:water}
\stp1{many high batches with a well-observed leader}Lemma~\ref{lem:levelC} ($j\ge\lonez(\varepsilon)$) gives $|\Shi_j|\ge\Me_\varepsilon+K+2$. For $i\in\Shi_j$, (W1) gives
$\log\Ng_{\Ld_{i-1}}(\tl_i)\ge\Lam_i\ge(1-\dl)\Lg_j$, i.e.\ $\Ng_{\Ld_{i-1}}(\tl_i)\ge\tl_j^{1-\dl}$, and $\tl_j^{1-\dl}\ge\up_j$ by Lemma~\ref{lem:constC}(a).

\stp2{at most $K$ high batches outside $\Sgo_j$}Let $i<i'$ be two elements of $\Shi_j\setminus\Sgo_j$ with the same leader $\Ld_{i-1}=\Ld_{i'-1}=:c$. Then
$\Ng_c(\tl_{i'-1})\ge\Ng_c(\tl_i)\ge\up_j$ by Step~1, so $i'\in\Sgo_j$: a contradiction. So $i\mapsto\Ld_{i-1}$ is injective on $\Shi_j\setminus\Sgo_j$, which has at most $K$
elements, and $|\Sgo_j|\ge\Me_\varepsilon+2$.

\stp3{accurate leaders}For $i\in\Sgo_j$ and $n:=\Ng_{\Ld_{i-1}}(\tl_{i-1})\ge\up_j$, outside $\Dev_j$,
$\tms_{i-1}=\tm_{\Ld_{i-1}}(\tl_{i-1})\le\hm_{\Ld_{i-1},n}\le\m_{\Ld_{i-1}}+\varepsilon$ (Lemma~\ref{lem:pseudo}(iii)).
\end{proof}

%----------------------------------------------------------------------
\subsection{Proof of Lemma~\ref{lem:highC}}\label{sec:pf-lem:highC}
\begin{proof}
\proves{lem:highC}
\uses{lem:kinf, def:hy, def:learner, def:grid, def:kinf, def:water}
Since $\tms_{i-1}\le r'<1$ and $\yg_{i-1}\ge1$, Lemma~\ref{lem:kinf}(iii) gives $\Kt_{i-1}(c)/\yg_{i-1}\le\Kt_{i-1}(c)\le D$. By (W1) at batch $i$,
\[
v\ \le\ \Lam_i\ \le\ \frac{\Ng_c(\tl_i)\Kt_{i-1}(c)}{\yg_{i-1}}+\log\Ng_c(\tl_i)\ \le\ \Ng_c(\tl_i)\,D+\log\bar N .
\]
So $\Ng_c(\tl_i)D\ge v-\log\bar N=Y>0$, hence $D>0$; and $D\le\kap(r')$ by Lemma~\ref{lem:kinf}(i).
\end{proof}

%----------------------------------------------------------------------
\subsection{Proof of Lemma~\ref{lem:fewC}}\label{sec:pf-lem:fewC}
\noindent\emph{Idea.} An arm whose mean is wrong either has many observations (event $\Dev_j$) or few. With few observations, it was explored in every good high batch,
and the chain lemma turns this into a large $n\Ki$ at some count $n\le\up_j$.

\begin{proof}
\proves{lem:fewC}
\uses{lem:goodhigh, lem:highC, lem:constC, lem:chain, lem:pseudo, def:evC, def:learner, def:grid, def:constC, def:kinf, def:model, def:hy, def:stacks, def:high}
Assume $\Dev_j^c$, and let $\Sgo_j$ be as in Lemma~\ref{lem:goodhigh}.

\stp1{common step}
\par\smallskip
\Setting an arm $c$ with $\Ng_c(\tl_{j-1})<\up_j$, and $r'\in[0,1)$ with $\kap(r')\le\kap(\m^\star+\varepsilon)$ and $\tms_{i-1}\le r'$ for every $i\in\Sgo_j$.
\Goal some $i\in\Sgo_j$ has $n:=\Ng_c(\tl_{i-1})\in[1,\up_j)$ and $n\Ki(\hnu_{c,n},r')\ge(1-4\dl)\Lg_j$.
\par\smallskip
\emph{Each good high batch gives evidence at its end.} Let $i\in\Sgo_j$ and $D_i:=\Ki(\tnu_c(\tl_{i-1}),r')$.
\begin{itemize}
\item $\Ng_c(\tl_i)\le\Ng_c(\tl_{j-1})<\up_j$, since $i\le j-1$.
\item Lemma~\ref{lem:highC} (\emph{evidence in a high batch}) applies with $v=(1-\dl)\Lg_j$ (batch $i$ is high) and $\bar N=\up_j$. It gives $D_i\in(0,\kap(r')]$ and
\[
\Ng_c(\tl_i)\,D_i\ \ge\ (1-\dl)\Lg_j-\log\up_j\ \ge\ (1-2\dl)\Lg_j\qquad\why{Lemma~\ref{lem:constC}(a)}.
\]
\end{itemize}
\emph{The chain turns evidence at the end into evidence at the start.} We want the evidence at the \emph{start} of a batch, $\Ng_c(\tl_{i-1})D_i$. Apply
Lemma~\ref{lem:chain} with
\[
w_i:=\Ng_c(\tl_{i-1}),\qquad Y:=(1-2\dl)\Lg_j,\qquad\delta':=\dl,\qquad D_{\max}:=\kap(r'),\qquad S:=\Sgo_j,\qquad j':=j-1\ge\max S .
\]
It gives one of two cases.
\begin{itemize}
\item \emph{Some $i\in\Sgo_j$ has $\Ng_c(\tl_{i-1})D_i\ge(1-\dl)(1-2\dl)\Lg_j$}, which is at least $(1-3\dl)\Lg_j$ by Lemma~\ref{lem:constC}(d).
\item \emph{Or the count at $\tl_{j-1}$ is large.} Since $|\Sgo_j|-1\ge\Me_\varepsilon$ and $\kap(r')\le\kap(\m^\star+\varepsilon)$,
\[
\Ng_c(\tl_{j-1})\ \ge\ (1-\dl)^{-\Me_\varepsilon}\frac{(1-2\dl)\Lg_j}{\kap(r')}\ \ge\ (1-\dl)^{-\Me_\varepsilon}\frac{(1-2\dl)\Lg_j}{\kap(\m^\star+\varepsilon)}\ \ge\ \up_j
\qquad\why{Lemma~\ref{lem:constC}(a)}.
\]
This contradicts $\Ng_c(\tl_{j-1})<\up_j$.
\end{itemize}
\emph{Conclusion.} So the first case holds. With $n:=\Ng_c(\tl_{i-1})$,
\begin{align*}
n\Ki(\hnu_{c,n},r')\ &\ge\ n\Ki(\tnu_c(\tl_{i-1}),r')-\wz\kap(r')&&\why{Lemma~\ref{lem:pseudo}(ii)}\\
&\ge\ (1-3\dl)\Lg_j-\wz\kap(\m^\star+\varepsilon)\ \ge\ (1-4\dl)\Lg_j&&\why{Lemma~\ref{lem:constC}(a)}.
\end{align*}

\stp2{first inclusion}Assume $\cU_j$: $\tms_{j-1}<\m^\star-2\varepsilon$.
\begin{itemize}
\item \emph{Every optimal arm has fewer than $\up_j$ observations at $\tl_{j-1}$.} Otherwise an optimal $a'$ with $n:=\Ng_{a'}(\tl_{j-1})\ge\up_j$ would have
$\hm_{a',n}\ge\m^\star-\varepsilon$ (outside $\Dev_j$), and $\tm_{a'}(\tl_{j-1})\ge\m^\star-2\varepsilon>\tms_{j-1}$ (Lemma~\ref{lem:pseudo}(iii), $\wz\le\varepsilon\up_j$), which contradicts
the definition of $\tms_{j-1}$ as a maximum.
\item \emph{The leaders of the good high batches are sub-optimal.} For $i\in\Sgo_j$, $\Ld_{i-1}$ has at least $\up_j$ observations at $\tl_{i-1}$, so it is sub-optimal, and
$\tms_{i-1}\le\mtwo+\varepsilon$ (Lemma~\ref{lem:goodhigh}).
\end{itemize}
The common step with $c=a^\star$ and $r'=\mtwo+\varepsilon$ gives $\Vs_j$.

\stp3{second inclusion}Assume $\hm_a(\tl_{j-1})\ge\m^\star-3\varepsilon$.
\begin{itemize}
\item \emph{$a$ has few observations.} $\m^\star-3\varepsilon>\m_a+\varepsilon$ (Lemma~\ref{lem:constC}(c)), so $\Ng_a(\tl_{j-1})\ge\up_j$ would put the outcome in $\Dev_j$. Hence
$\Ng_a(\tl_{j-1})<\up_j$.
\item \emph{The levels of the good high batches are below $\m^\star+\varepsilon$.} For $i\in\Sgo_j$, $\tms_{i-1}\le\m_{\Ld_{i-1}}+\varepsilon\le\m^\star+\varepsilon<1$
(Lemma~\ref{lem:goodhigh}).
\end{itemize}
The common step with $c=a$ and $r'=\m^\star+\varepsilon$ gives $n\le n':=\Ng_a(\tl_{j-1})<\up_j$ with
\[
n\Ki(\hnu_{a,n},\m^\star+\varepsilon)\ \ge\ (1-4\dl)\Lg_j\qquad\text{and}\qquad\hm_{a,n'}=\hm_a(\tl_{j-1})\ \ge\ \m^\star-3\varepsilon:
\]
$\Va_j$ occurs.
\end{proof}

%----------------------------------------------------------------------
\subsection{Proof of Lemma~\ref{lem:logcountC}}\label{sec:pf-lem:logcountC}
\begin{proof}
\proves{lem:logcountC}
\uses{lem:goodhigh, lem:kinf, def:hy, def:model, def:kinf, def:learner, def:grid, def:constC, def:water}
Take $i\in\Sgo_j$ (non-empty by Lemma~\ref{lem:goodhigh}). Then $\tms_{i-1}\le\m^\star+\varepsilon<1$, so $\Kt_{i-1}(c)/\yg_{i-1}\le\kap(\m^\star+\varepsilon)$ (Lemma~\ref{lem:kinf}(i),(iii)).
By (W1) at $i$, $\Ng_c(\tl_i)\le\Ng_c(\tl_{j-1})$ and $\log N\le N$,
\[
(1-\dl)\Lg_j\ \le\ \Lam_i\ \le\ \frac{\Ng_c(\tl_i)\Kt_{i-1}(c)}{\yg_{i-1}}+\log\Ng_c(\tl_i)\ \le\ \kap(\m^\star+\varepsilon)\Ng_c(\tl_{j-1})+\Ng_c(\tl_{j-1}).
\]
\end{proof}

%----------------------------------------------------------------------
\subsection{Proof of Lemma~\ref{lem:optW}}\label{sec:pf-lem:optW}
\noindent\emph{Idea.} The leaders of the high batches have $\tl^{1-\dl}$ observations. If no optimal arm had $\frac12\Ug$ of them, these leaders would all be
sub-optimal, the levels of the good high batches would be below $\mtwo+\varepsilon$, and the chain argument would turn the exploration of an optimal arm into the event
$\Vs_j$.

\begin{proof}
\proves{lem:optW}
\uses{lem:goodhigh, lem:pseudo, lem:highC, lem:chain, lem:constC, def:hy, def:grid, def:learner, def:constC, def:evC, def:model, def:kinf, def:stacks, def:high}
Let $\bar u:=\frac12\Ug_{j-1}$, and assume by contradiction that every optimal arm has fewer than $\bar u$ observations at $\tl_{j-1}$.

\stp1{the leaders of the high batches are sub-optimal}For $i\in\Shi_j$, Lemma~\ref{lem:goodhigh} gives $\Ng_{\Ld_{i-1}}(\tl_{j-1})\ge\Ng_{\Ld_{i-1}}(\tl_i)\ge\tl_j^{1-\dl}\ge\bar u$,
so $\Ld_{i-1}$ is sub-optimal. Outside $\Dev_j$, $\tms_{i-1}\le\mtwo+\varepsilon$ for $i\in\Sgo_j$.

\stp2{$a^\star$ is poorly observed at the good high batches}For every $i\in\Sgo_j$, $\Ng_{a^\star}(\tl_{i-1})<\up_j$. Otherwise, outside $\Dev_j$,
$\tm_{a^\star}(\tl_{i-1})\ge\m^\star-2\varepsilon>\mtwo+\varepsilon\ge\tms_{i-1}$ (Lemma~\ref{lem:pseudo}(iii), $\wz\le\varepsilon\up_j$), which contradicts the definition of $\tms_{i-1}$.

\stp3{the chain}This is the common step of Lemma~\ref{lem:fewC} for $c=a^\star$, with one batch removed, because here we only know that $a^\star$ is poorly observed at
the good high batches (Step~2), not at $\tl_{j-1}$.
\begin{itemize}
\item Let $i_m:=\max\Sgo_j$ and $S':=\Sgo_j\setminus\{i_m\}$, so $|S'|\ge\Me_\varepsilon+1$.
\item \emph{Evidence at the end of each batch of $S'$.} For $i\in S'$, $\Ng_{a^\star}(\tl_i)\le\Ng_{a^\star}(\tl_{i_m-1})<\up_j$ (Step~2). Lemma~\ref{lem:highC} with $v=(1-\dl)\Lg_j$,
$r'=\mtwo+\varepsilon$ and $\bar N=\up_j$ gives, with $D_i:=\Ki(\tnu_{a^\star}(\tl_{i-1}),\mtwo+\varepsilon)\in(0,\kap(\mtwo+\varepsilon)]$,
\[
\Ng_{a^\star}(\tl_i)\,D_i\ \ge\ (1-2\dl)\Lg_j\qquad(i\in S').
\]
\item \emph{The chain.} Lemma~\ref{lem:chain} with $w_i:=\Ng_{a^\star}(\tl_{i-1})$, $Y:=(1-2\dl)\Lg_j$ and $j'=i_m-1$ gives one of two cases:
\begin{itemize}
\item some $i\in S'$ has $\Ng_{a^\star}(\tl_{i-1})D_i\ge(1-3\dl)\Lg_j$;
\item or $\Ng_{a^\star}(\tl_{i_m-1})\ge(1-\dl)^{-\Me_\varepsilon}(1-2\dl)\Lg_j/\kap(\m^\star+\varepsilon)\ge\up_j$ (Lemma~\ref{lem:constC}(a)), which contradicts Step~2.
\end{itemize}
\item \emph{The first case is $\Vs_j$.} With $n:=\Ng_{a^\star}(\tl_{i-1})\in[1,\up_j)$, Lemma~\ref{lem:pseudo}(ii) and Lemma~\ref{lem:constC}(a) give
$n\Ki(\hnu_{a^\star,n},\mtwo+\varepsilon)\ge(1-4\dl)\Lg_j$. So $\Vs_j$ occurs, a contradiction.
\end{itemize}

\stp4{conclusion}So some optimal $a'$ is in $\Wg_{j-1}$. Every arm of $\Wg_{j-1}$ has at least $\bar u\ge\up_j$ observations, hence an $\varepsilon$-accurate mean outside $\Dev_j$:
\begin{itemize}
\item $\tm_{a'}(\tl_{j-1})\ge\m^\star-\varepsilon-2\wz/\Ug_{j-1}\ge\m^\star-2\varepsilon$;
\item every sub-optimal $c\in\Wg_{j-1}$ has $\tm_c(\tl_{j-1})\le\mtwo+\varepsilon<\m^\star-2\varepsilon$.
\end{itemize}
Hence $\rr_{j-1}$ is optimal, and $\m^\star-2\varepsilon\le\mh_{j-1}\le\hm_{\rr_{j-1}}(\tl_{j-1})\le\m^\star+\varepsilon$.
\end{proof}

%----------------------------------------------------------------------
\subsection{Proof of Lemma~\ref{lem:pullsC}}\label{sec:pf-lem:pullsC}
\begin{proof}
\proves{lem:pullsC}
\uses{lem:kinf, lem:constC, lem:optW, lem:pseudo, lem:imedround, lem:service, def:evC, def:stacks, def:grid, def:model, def:constC, def:hy, def:kinf, def:learner, def:amb}
\stp{(i)}{$a$ is neither the leader nor the reference}
\emph{The divergence $D$ and the mean of $a$}, in each class.
\begin{itemize}
\item On $\cO_j$: $\hm_a(\tl_{j-1})<\m^\star-3\varepsilon$ by definition, so $D\ge\varepsilon^2/8$ (Lemma~\ref{lem:kinf}(iv)).
\item On $\cG_j$: $D>\Dp_\varepsilon$ by definition, and $\Dp_\varepsilon\ge\varepsilon^2/8$ (Lemma~\ref{lem:constC}(b)). Since $D>0$, $\hm_a(\tl_{j-1})<\m^\star-2\varepsilon$
(Lemma~\ref{lem:kinf}(ii)).
\end{itemize}
\emph{$a$ is not the leader.} In both classes, $\tm_a(\tl_{j-1})\le\hm_a(\tl_{j-1})<\m^\star-2\varepsilon\le\tms_{j-1}$.

\noindent\emph{$a$ is not the reference.} $\rr_{j-1}$ is optimal (Lemma~\ref{lem:optW}).

\noindent\emph{The IMED and test divergences are at least $D$.} The levels satisfy $\tms_{j-1}\ge\mh_{j-1}\ge\m^\star-2\varepsilon$ (Lemma~\ref{lem:optW}). So
\[
\Kt_{j-1}(a)\ \ge\ \Kh_{j-1}(a)\ \ge\ \Ki\big(\hnu_a(\tl_{j-1}),\mh_{j-1}\big)\ \ge\ D\qquad\why{Lemma~\ref{lem:kinf}(iii); Lemma~\ref{lem:pseudo}(i); Lemma~\ref{lem:kinf}(iii)},
\]
with $\Kt_{j-1}(a)=\infty$ if $\tms_{j-1}=1$.

\stp{(ii)}{batch not activated, or $a$ rejected}Let $t$ be the last round of batch $j$ at which $a$ is pulled, so $\Ng_a(\tl_j)=\Ng_a(t)$.
\begin{itemize}
\item If $t$ is an IMED round, Lemma~\ref{lem:imedround} (\emph{IMED rounds}) and (i) give
\[
\Ng_a(t)\ =\ \Ng_a(t-1)+1\ <\ 1+\frac{\yg_{j-1}\log t}{\Kt_{j-1}(a)}\ \le\ 1+\frac{\yg_{j-1}\Lg_j}{D}.
\]
\item If $t$ is a paced round, the batch is activated, so $a$ is rejected by assumption, and
\[
\Ng_a(t)\ \le\ \nt_{j-1,a}(t)\ \le\ K-1+\Big\lceil\frac{(1+\om_{j-1})\yg_{j-1}\log t}{\Kh_{j-1}(a)}\Big\rceil\ <\ K+\frac{(1+\om_{j-1})\yg_{j-1}\Lg_j}{D}.
\]
\end{itemize}

\stp{(iii)}{batch activated, $a$ ambiguous}\emph{IMED phase.} By Lemma~\ref{lem:imedround} (\emph{IMED rounds}) with (i), $\log\sA_{j-1}\le2\Lg_{j-1}$ ($\thv\le1$) and $\yg_{j-1}\le\yb$,
\[
\Ng_a(\sA_{j-1})\ \le\ \max\Big(\Ng_a(\tl_{j-1}),\ 1+\frac{8\yb\cdot2\Lg_{j-1}}{\varepsilon^2}\Big).
\]
\emph{Paced phase.} By Lemma~\ref{lem:service}(a) (\emph{queue lemma}) and $\nt_{j-1,a}\le\Ug_j+K-1$,
\[
\Ng_a(\tl_j)\ \le\ \max\big(\Ng_a(\sA_{j-1}),\ \Ug_j+K-1\big)\ \le\ \Ng_a(\tl_{j-1})+\max\Big(1+\frac{16\yb\Lg_{j-1}}{\varepsilon^2},\ \Ug_j+K-1\Big)\ \le\ \Ng_a(\tl_{j-1})+\cell_{j-1}.
\]
\end{proof}

%----------------------------------------------------------------------
\subsection{Proof of Lemma~\ref{lem:amb}}\label{sec:pf-lem:amb}
\noindent\emph{Idea.} The arm $a$ is in $\Qp_\ell$ unless it is well sampled (then its evidence is huge) or it was the previous leader (a bad event). In $\Qp_\ell$, its old
evidence bounds the threshold from above; its new evidence is at least $1-\gam_\ell/2$ times the old one, because both test levels are $\zet_\ell$-accurate and
$\Ki(\hnu_{a,n},\m^\star)$ is $\zetp_\ell$-accurate at both counts. The failure probability is $e^{-c\sqrt{\Lg_\ell}/\log^2\Lg_\ell}$, and the activated batches are sparse
(Lemma~\ref{lem:sparse}, \emph{sparse activations}).

\par\smallskip
\Setting $\varepsilon\in(0,\ez]$, $\ell\ge\lfour(\varepsilon)$, and (for Steps~1--4) an outcome in $\Eamb_\ell(\varepsilon)^c$. Write $N:=\Ng_a(\tl_\ell)$.
\Goal Steps~1--4: $a\notin\{\Ld_\ell,\rr_\ell\}$, $\Kh_\ell(a)\ge\varepsilon^2/8$, and $\Zg_\ell(a)\ge\hg_\ell$ (so $a$ is rejected). Steps~5--7: $\Cfour(\varepsilon)<\infty$.
\par\smallskip

\begin{proof}
\proves{lem:amb}
\uses{lem:optW, lem:fewC, lem:pseudo, lem:kinf, lem:logcountC, def:amb, lem:controlC, lem:hoeffding, lem:lowerdev, lem:upperdev, lem:sparse, def:grid, def:hy, def:model, def:stacks, def:evC, def:kinf, def:constC, def:learner, def:ranksC, def:water}
\stp1{statistics at $\tl_\ell$}
\begin{itemize}
\item Lemma~\ref{lem:optW} at $j=\ell+1$: some optimal $a'$ lies in $\Wg_\ell$, $\rr_\ell$ is optimal and $\mh_\ell\ge\m^\star-2\varepsilon$.
\item Lemma~\ref{lem:fewC} at $j=\ell+1$: $\hm_a(\tl_\ell)<\m^\star-3\varepsilon$ (outside $\Va_{\ell+1}$). So $\Kh_\ell(a)\ge\Ki(\hnu_a(\tl_\ell),\m^\star-2\varepsilon)\ge\varepsilon^2/8$
(Lemmas~\ref{lem:pseudo}(i), \ref{lem:kinf}(iii),(iv)); and $\tm_a(\tl_\ell)<\m^\star-2\varepsilon\le\mh_\ell\le\tms_\ell$, so $a\ne\Ld_\ell$; and $a\ne\rr_\ell$.
\item Lemma~\ref{lem:logcountC} at $j=\ell+1$: $N\ge(1-\dl)\Lg_{\ell+1}/(\kap(\m^\star+\varepsilon)+1)$, so $N\ge\nl_\ell$.
\end{itemize}

\stp2{statistics at $\tl_{\ell-1}$}\emph{$a\ne\Ld_{\ell-1}$.} Otherwise $\hm_a(\tl_{\ell-1})\ge\tms_{\ell-1}$. If $\tms_{\ell-1}<\m^\star-2\varepsilon$, Lemma~\ref{lem:fewC} at $j=\ell$
puts the outcome in $\Dev_\ell\cup\Vs_\ell$; otherwise $\hm_a(\tl_{\ell-1})\ge\m^\star-3\varepsilon$ puts it in $\Dev_\ell\cup\Va_\ell$. Both are excluded.

\noindent\emph{A logarithmic count.} Lemma~\ref{lem:logcountC} at $j=\ell$ gives $\Ng_a(\tl_{\ell-1})\ge\nl_\ell$.

\stp3{if $a\in\Wg_\ell$, $a$ is rejected}Then $N\ge\frac12\Ug_\ell$ and
\[
\Zg_\ell(a)\ \ge\ \tfrac12\Ug_\ell\cdot\varepsilon^2/8\ \ge\ \yb\Lg_\ell\ \ge\ \yg_{\ell-1}\Lh_\ell\ \ge\ \hg_\ell\qquad\why{(b) of Definition~\ref{def:amb}; $\Lh_\ell\le\Lg_\ell$},
\]
and $\Kh_\ell(a)>0$.

\stp4{if $a\notin\Wg_\ell$, $a$ is rejected}Then $a\in\Qp_\ell$ (Step~2), so, by definition of $\Lh_\ell$ as a minimum,
\begin{equation}\label{eq:old}
N\,\Kh_{\ell-1}(a)\ \ge\ \yg_{\ell-1}\big(\Lh_\ell-\log N\big).
\end{equation}
Write $\gamma:=\gam_\ell$, $\zeta:=\zet_\ell$, $\zeta':=\zetp_\ell$, so that $\frac{2\zeta}{1-\m^\star}=\zeta'$.

\noindent\emph{Case (ii): $\log N>\frac\gamma2\Lh_\ell$.} By (c) ($u\mapsto u/\log u$ is non-decreasing on $[e,\infty)$, and $N\ge\nl_\ell\ge e$) and Step~1,
\[
\hg_\ell\ \le\ \yg_{\ell-1}\Lh_\ell\ <\ \frac{2\yb}\gamma\log N\ \le\ \frac{\varepsilon^2}8N\ \le\ N\Kh_\ell(a)\ =\ \Zg_\ell(a).
\]

\noindent\emph{Case (i): $\log N\le\frac\gamma2\Lh_\ell$.}

\noindent\emph{The two test levels are $\zeta$-accurate.} $\Wg_{\ell-1}\ne\emptyset$ (as $\tl_{\ell-1}\ge K\Ug_{\ell-1}$: the last assumption of Lemma~\ref{lem:optW} at $j=\ell$, which (a) of
Definition~\ref{def:amb} gives), and every arm of $\Wg_{\ell-1}\cup\Wg_\ell$ has at least
$\frac12\Ug_{\ell-1}$ observations. So on $\Acc_\ell$, by Lemma~\ref{lem:pseudo}(iii) and (d),
\[
\mh_{\ell-1}\ \le\ \hm_{\rr_{\ell-1}}(\tl_{\ell-1})\ \le\ \m^\star+\zeta,\qquad\mh_\ell\ \ge\ \tm_{a'}(\tl_\ell)\ \ge\ \m^\star-\zeta-\frac{2\wz}{\Ug_\ell}\ \ge\ \m^\star-2\zeta .
\]
\emph{The old divergence is at most $\da(1+\frac{3\gamma}{32})$.}
\begin{align*}
\Kh_{\ell-1}(a)\ &\le\ \Ki\big(\tnu_a(\tl_{\ell-1}),\m^\star+\zeta\big)&&\why{Lemma~\ref{lem:kinf}(iii)}\\
&\le\ \Ki\big(\hnu_a(\tl_{\ell-1}),\m^\star+\zeta\big)+\frac{\wz\kap(\m^\star+\varepsilon)}{\nl_\ell}&&\why{Lemma~\ref{lem:pseudo}(ii), $\Ng_a(\tl_{\ell-1})\ge\nl_\ell$, $\zeta\le\varepsilon$}\\
&\le\ \Ki\big(\hnu_a(\tl_{\ell-1}),\m^\star\big)+\frac{2\zeta}{1-\m^\star}+\zeta'&&\why{Lemma~\ref{lem:kinf}(v) at $\mu=\m^\star+\zeta$; (d)}\\
&\le\ \da+3\zeta'\ =\ \da\Big(1+\frac{3\gamma}{32}\Big)&&\why{$\KC_\ell$ at $\Ng_a(\tl_{\ell-1})\ge\nl_\ell$}.
\end{align*}
(In the third line, $1-\m^\star-\zeta\ge\frac{1-\m^\star}2$ by (d).)

\noindent\emph{The new divergence is at least $\da(1-\frac\gamma{16})$.}
\begin{align*}
\Kh_\ell(a)\ &\ge\ \Ki\big(\hnu_a(\tl_\ell),\m^\star-2\zeta\big)&&\why{Lemmas~\ref{lem:pseudo}(i), \ref{lem:kinf}(iii)}\\
&\ge\ \Ki\big(\hnu_a(\tl_\ell),\m^\star\big)-\frac{2\zeta}{1-\m^\star}&&\why{Lemma~\ref{lem:kinf}(v), shift $2\zeta\le\m^\star$}\\
&\ge\ \da-2\zeta'\ =\ \da\Big(1-\frac\gamma{16}\Big)&&\why{$\KC_\ell$ at $N\ge\nl_\ell$}.
\end{align*}
\emph{Conclusion of case (i).} Since $\gamma\le\frac12$, $(1-\frac\gamma2)(1+\frac{3\gamma}{32})=1-\frac{13\gamma}{32}-\frac{3\gamma^2}{64}\le1-\frac\gamma{16}$, so
$\Kh_\ell(a)\ge(1-\frac\gamma2)\Kh_{\ell-1}(a)$. Then
\begin{align*}
\Zg_\ell(a)\ =\ N\Kh_\ell(a)\ &\ge\ \Big(1-\frac\gamma2\Big)N\Kh_{\ell-1}(a)\\
&\ge\ \Big(1-\frac\gamma2\Big)\yg_{\ell-1}\big(\Lh_\ell-\log N\big)&&\why{\eqref{eq:old}}\\
&\ge\ \Big(1-\frac\gamma2\Big)^2\yg_{\ell-1}\Lh_\ell\ \ge\ (1-\gamma)\yg_{\ell-1}\Lh_\ell\ =\ \hg_\ell&&\why{$\log N\le\tfrac\gamma2\Lh_\ell$}.
\end{align*}
In all cases $\Kh_\ell(a)\ge\varepsilon^2/8>0$ and $\Zg_\ell(a)\ge\hg_\ell$, so $a$ is rejected.

\stp5{probabilities of the eight events}\emph{The events $\Dev,\Vs,\Va$.} By Lemma~\ref{lem:controlC} and $\up^3+\up^4\le2(1+\up^4)$, $\PP(\Dev_j)+\PP(\Vs_j)+\PP(\Va_j)\le\pi(\Lg_j)$ for
$j\ge\lz(\varepsilon)$; by (f), $\pi(\Lg_{\ell+1})\le\pi(\Lg_\ell)$. Together, at most $2\pi(\Lg_\ell)$.

\noindent\emph{The event $\Acc_\ell^c$.} Let $c_\zeta:=\da(1-\m^\star)/64$.
\begin{itemize}
\item By (e) of Definition~\ref{def:amb}, $\gam_\ell=\Lg_\ell^{-1/4}$, so $\zet_\ell=c_\zeta\Lg_\ell^{-1/4}$.
\item $\Ug_{\ell-1}\ge\Lg_{\ell-2}^\beta$ by definition of $\Ug$, and $\Lg_{\ell-2}^\beta\ge(\Lg_\ell/4)^\beta$ by (e) again.
\item Lemma~\ref{lem:hoeffding}(b) with $u=\frac12\Ug_{\ell-1}$, a union bound over the arms, and $\frac1{1-e^{-v}}\le1+\frac1v$.
\end{itemize}
Together they give
\[
\PP(\Acc_\ell^c)\ \le\ 2K\Big(1+\frac1{2\zet_\ell^2}\Big)e^{-\zet_\ell^2\Ug_{\ell-1}}\ \le\ p_A(\Lg_\ell):=2K\Big(1+\frac{\Lg_\ell^{1/2}}{2c_\zeta^2}\Big)\exp\big(-c_\zeta^2\,4^{-\beta}\,\Lg_\ell^{\beta-1/2}\big).
\]
\emph{The event $\KC_\ell^c$.} Lemmas~\ref{lem:lowerdev} and~\ref{lem:upperdev} with $\nu=\nu_a$, $\mu=\m^\star$, $\varepsilon'=\zetp_\ell$, and a union bound over $n\ge\nl_\ell$, give
\[
\PP(\KC_\ell^c)\ \le\ p_K(\Lg_\ell):=\frac{(1+n^{\mathrm g}_\ell)\,e^{-\rho'_\ell\nl_\ell}}{1-e^{-\rho'_\ell}},
\]
where $R'_\ell:=\kap(\m^\star)-\log(1-e^{-\zetp_\ell/4})$, $\rho'_\ell:=\zetp_\ell^{\,2}/(2R_\ell'^2)$ and
$n^{\mathrm g}_\ell:=\lceil8/(\zetp_\ell(1-e^{-\zetp_\ell/4})(1-\m^\star))\rceil+1$; by (e) it is a function of $\Lg_\ell$.

\stp6{an explicit bound on $p_K$}Let $y:=\Lg_\ell\ge16$, so $\zetp_\ell=\frac{\da}{32}y^{-1/4}$, and
\begin{gather*}
a_0:=\min\Big(\frac{\da}{128},1\Big),\quad A_1:=\frac14+\frac{\kap(\m^\star)+\log(2/a_0)}{\log16},\quad c_1:=\frac{\da^2}{2048A_1^2},\\
C_2:=3+\frac{512}{\da a_0(1-\m^\star)},\quad c_2:=\frac{1-\dl}{\kap(\m^\star+\varepsilon)+1}.
\end{gather*}
Since $1-e^{-v}\ge\min(v,1)/2$ for $v\ge0$ and $y^{-1/4}\le1$, $1-e^{-\zetp_\ell/4}\ge\frac{a_0}2y^{-1/4}$. Hence, using $\log y\ge\log16$,
\begin{align*}
R'_\ell\ &\le\ \kap(\m^\star)+\log\frac2{a_0}+\frac14\log y\ \le\ A_1\log y,\\
\rho'_\ell\ &=\ \frac{\zetp_\ell^{\,2}}{2R_\ell'^2}\ \ge\ \frac{\da^2y^{-1/2}}{2048A_1^2\log^2y}\ =\ \frac{c_1y^{-1/2}}{\log^2y},\\
\frac1{1-e^{-\rho'_\ell}}\ &\le\ 1+\frac1{\rho'_\ell}\ \le\ \Big(1+\frac1{c_1}\Big)y^{1/2}\log^2y,\\
1+n^{\mathrm g}_\ell\ &\le\ 3+\frac8{\frac{\da}{32}y^{-1/4}\cdot\frac{a_0}2y^{-1/4}(1-\m^\star)}\ \le\ C_2\,y^{1/2},\\
\nl_\ell\ &\ge\ c_2\,y .
\end{align*}
Therefore
\[
p_K(y)\ \le\ \bar p_K(y):=C_2\Big(1+\frac1{c_1}\Big)\,y\log^2y\ \exp\Big(-\frac{c_1c_2\sqrt y}{\log^2y}\Big)\qquad(y\ge16).
\]

\stp7{the series $\Cfour(\varepsilon)$}Let $\cell(y):=y^\beta+1+K+16\yb y/\varepsilon^2$, so that $\cell_\ell\le\cell(\Lg_\ell)$, and let
\[
f(y):=\cell(y)\big(2\pi(y)+p_A(y)+\bar p_K(y)\big)\ \text{ for }y\ge16,\qquad f(y):=0\ \text{ for }y\in[1,16).
\]
By Step~5, $\cell_\ell\PP(\Eamb_\ell)\le f(\Lg_\ell)$ for $\ell\ge\lfour$. Since $\Lg_\ell\ge16$ for $\ell\ge\lfour$ and $\tl_\ell\ge e^e$ for $\ell\ge\lW-1$, Lemma~\ref{lem:sparse} (\emph{sparse activations}) gives
\[
\Cfour(\varepsilon)\ \le\ \sum_{k\ge4}(k+1)\sup_{y\in[2^k,2^{k+1})}f(y).
\]
On $[2^k,2^{k+1})$, the polynomial prefactors are at most their values at $2^{k+1}$, $\sqrt y\ge2^{k/2}$, and $\log y\le(k+1)\log2$. So the $k$-th term is at most a polynomial
in $2^k$ times
\[
e^{-(1+\ezz)2^k}+e^{-c_\zeta^24^{-\beta}2^{k(\beta-1/2)}}+\exp\Big(-\frac{c_1c_2\,2^{k/2}}{((k+1)\log2)^2}\Big),
\]
and the series converges.
\end{proof}

%----------------------------------------------------------------------
\section{End of the proof of Theorem~\ref{thm:upper}}\label{sec:app-end}

\subsection{Proof of Theorem~\ref{thm:upper}}\label{sec:app-proofupper}
\begin{proof}
\proves{thm:upper}
\uses{thm:lower, lem:diag, prop:fixedC, prop:dyA, def:hy, def:grid, def:model, def:learner, def:kinf, def:constC}
\par\smallskip
\Setting \HY{} with the default parameters (the proof works for every $\wz\ge0$ and $\beta\ge2$), a grid with $\xg_\ell\to x\in[1,\infty)$, $\nu\in\cP$, a sub-optimal arm $a$.
\Goal consistency, and $\lim_\ell\E[\Ng_a(\tl_\ell)]/\Lg_\ell=x^2/\da$. The upper bound comes from Section~\ref{sec:app-xone} ($x=1$) or~\ref{sec:app-xsup} ($x>1$), the lower bound
from Theorem~\ref{thm:lower}, which needs consistency.
\par\smallskip

\paragraph{Case $x=1$.} \emph{Upper bound.} By Lemma~\ref{lem:diag} with $\bar\varepsilon=\ez$ and $f=\RHo$, there is a sequence $\varepsilon_L\to0$ with $\RHo(\varepsilon_L)/\Lg_L\to0$.
Proposition~\ref{prop:fixedC} with $\varepsilon=\varepsilon_L$ gives
\[
\frac{\E[\Ng_a(\tl_L)]}{\Lg_L}\ \le\ \frac1{\da}+\Co\varepsilon_L+\frac{\RHo(\varepsilon_L)}{\Lg_L}\ \xrightarrow[L\to\infty]{}\ \frac1{\da}.
\]
\emph{Consistency.} With a fixed $\varepsilon$, Proposition~\ref{prop:fixedC} gives $\E[\Ng_c(\tl_L)]=O(\Lg_L)$ for every instance of $\cP$ and every sub-optimal arm $c$; so \HY{}
is consistent on $\cP$.

\noindent\emph{Lower bound.} Theorem~\ref{thm:lower}(i) gives $\liminf_L\E[\Ng_a(\tl_L)]/\Lg_L\ge1/\da=x^2/\da$.

\paragraph{Case $x>1$.} \emph{Upper bound and consistency.} Proposition~\ref{prop:dyA} gives $\limsup_t\E[\Ng_a(t)]/\log t\le x^2/\da$, hence along $(\tl_\ell)$, and
$\E[\Ng_c(\tl_\ell)]=O(\Lg_\ell)$ for every instance and every sub-optimal $c$: \HY{} is consistent.

\noindent\emph{Lower bound.} The counts are logarithmic, so Theorem~\ref{thm:lower}(ii) gives $\liminf_\ell\E[\Ng_a(\tl_\ell)]/\Lg_\ell\ge x^2/\da$.

\paragraph{Regret.} Sum the limits over the finitely many sub-optimal arms, weighted by $\Dg_c$.
\end{proof}

%======================================================================
\section{Sanity checks}\label{sec:sanity}

\begin{itemize}
\item \emph{Batches of size one ($x=1$, the sequential case).} Take $\Bg_1:=\max(2,K)$ and $\Bg_\ell:=1$ for $\ell\ge2$. Then
$\tl_\ell=\Bg_1+\ell-1$ and $\xg_\ell\to1$, and a batched learner observes every reward before the next round: the setting
is the sequential one. Theorem~\ref{thm:main} then gives the constant $1/\Ki(\nu_a,\m^\star)$, the sequential optimal
constant for distributions on $[0,1]$ \cite{BK96} (to be checked), \cite[(2) and Theorem~8]{HT10}. On this grid no batch of \HY{} is
activated: $\sA_\ell=\lfloor\tl_\ell^{1+\thv_\ell}\rfloor\ge\tl_\ell+1=\tl_{\ell+1}$, since $\tl^{1+\vartheta}\ge2t\ge t+1$ when
$\tl_\ell<e^e$ ($\vartheta=1$), and $\tl^{1+\vartheta}=t\,e^{\log t/\log\log t}\ge e\,t\ge t+1$ when $\tl_\ell\ge e^e$. So every
round is an IMED round with statistics refreshed at every round: \HY{} is IMED \cite{HT15}, with the weight
$\yg_\ell=\xg_\ell^2\to1$ on the logarithm and the pseudo-observation of weight $\wz$.
\item \emph{An explicit example, $K=2$, Bernoulli rewards.} Let $\nu_1:=\Ber(1/2)$ and $\nu_2:=\Ber(3/10)$, so that
$\nu\in\cP$ ($\m^\star=1/2<1$) and $a=2$. For $0<p<\mu<1$ the dual formula gives $\Ki(\Ber(p),\mu)=\kl(p,\mu)$ (the maximizer
is $\lambda^\star=(\mu-p)/(\mu(1-\mu))\in\Lset{\mu}$, where $1-\lambda^\star(1-\mu)=p/\mu$ and $1+\lambda^\star\mu=(1-p)/(1-\mu)$),
so $\da=\kl(0.3,0.5)\approx0.0823$. On the grid $\tl_\ell:=2^{2^\ell}$ ($\tl_1=4\ge\max(2,K)$, and $\xg_\ell=2$ for $\ell\ge2$),
Theorem~\ref{thm:main} gives $\E[\Ng_2(\tl_\ell)]\sim4\log\tl_\ell/\da\approx48.6\log\tl_\ell$; on the grid
$\tl_\ell:=2^{\ell+1}$ ($\xg_\ell=(\ell+1)/\ell\to1$) it gives $\approx12.2\log\tl_\ell$. All the hypotheses hold at once:
rewards in $[0,1]$, $\m^\star<1$, $\Bg_1\ge K$, $\xg_\ell\to x$.
\item \emph{The constant is finite and positive.} By Lemma~\ref{lem:kinf}(i) and~(iv) (with $q:=\m_a$),
$\Dg_a^2/8\le\da\le\kap(\m^\star)<\infty$ for every sub-optimal arm $a$ of an instance of $\cP$.
\item \emph{Boundary of the class.} If $\m^\star=1$, then $\Ki(\nu_a,1)=+\infty$ for every arm of mean smaller than $1$, and
the first-order constant would be $0$: such instances are excluded from $\cP$.
\item \emph{Several optimal arms and ties.} $\cP$ allows several optimal arms, and the tie rules of \HY{} are arbitrary; the
proofs use neither a unique optimal arm nor a particular tie rule (Subsection~\ref{sec:algo}, Properties~(ii)).
\item \emph{Every arm observed.} Batch $1$ has $\Bg_1\ge K$ rounds, so every count is at least $1$ at every end
(Definition~\ref{def:hy-vocab}): no $\log0$ and no empty empirical distribution occurs.
\item \emph{Consistency with the roadmap.} KNOWN.md of BatchedBandits has no \texttt{[Lean]} entry yet. The statements
agree with Appendix~A, items~1--2, of the roadmap, and with the author's private Lean statements
(Subsection~\ref{sec:lean}).
\end{itemize}

%======================================================================
\section{Weak points}\label{sec:weak}

The independent verifications (by AI agents) found no error; these are the steps where an error would be the hardest to
see.
\begin{enumerate}
\item[(1)] Lemma~\ref{lem:twobatches}, Steps 3--4: the order of the limits ($\ell\to\infty$, then $b\to\infty$, then
$C\to\infty$), with constants that depend on $C$ through the tilt, and the reason why the two instances $\nu^{\rhom}$ and
$\nu^{\rhop}$ are used instead of $\nu$ itself. This is also where the logarithmic counts of Theorem~\ref{thm:lower}(ii)
enter.
\item[(2)] Lemma~\ref{lem:amb} (a sub-optimal arm passes the consistency test): the bound on the probability that $a$ is
ambiguous at the start of an activated batch, of order $e^{-c\sqrt{\Lg_\ell}/\log^2\Lg_\ell}$, and its summability through
Lemma~\ref{lem:sparse}; the constants of its Steps 5--7.
\item[(3)] Lemma~\ref{lem:W} (the IMED phase lifts the water level) and the chain argument of the case $x=1$
(Lemmas~\ref{lem:levelC}, \ref{lem:fewC} and~\ref{lem:optW}), which must hold from some rank on for every grid with
$\xg_\ell\to1$, including short batches alternating with long ones.
\end{enumerate}

%======================================================================
\section{Open questions}\label{sec:open}

\begin{openquestion}[explicit remainders]\label{oq:remainders}
Is there a learner of $\Mz$ such that, for every $x\in[1,\infty)$, every grid with $\xg_\ell\to x$, every $\nu\in\cP$ and every
sub-optimal arm $a$, and for every $\ell$,
$\E[\Ng_a(\tl_\ell)]\le\frac{x^2}{\Ki(\nu_a,\m^\star)}\log\tl_\ell+C(\nu)\,(\log\tl_\ell)^{\alpha}+r_\ell$, with an exponent $\alpha<1$,
a constant $C(\nu)$ polynomial in $1/\Ki(\nu_a,\m^\star)$, $1/(1-\m^\star)$ and $K$, and a term $r_\ell=o(\log\tl_\ell)$ given
explicitly by the grid (through $\xg_\ell-x$)?
\end{openquestion}
\noindent\emph{Why it is open.} The remainders of \HY{} come from ranks after which the margins, the nets and the
deviation inequalities with a gain take effect; they are existential in the Lean proof and astronomical on paper. It
extends roadmap directions D1 and D2.

\begin{openquestion}[the liminf without logarithmic counts]\label{oq:liminf}
Let $x>1$, a grid with $\xg_\ell\to x$, a batched learner consistent on $\cP$, $\nu\in\cP$ and a sub-optimal arm $a$. Does
$\liminf_\ell\E_\nu[\Ng_a(\tl_\ell)]/\log\tl_\ell\ge x^2/\Ki(\nu_a,\m^\star)$ hold without the assumption
$\E_\nu[\Ng_a(\tl_\ell)]=O(\log\tl_\ell)$ of Theorem~\ref{thm:lower}(ii), or is there a consistent learner that violates it?
\end{openquestion}
\noindent\emph{Why it is open.} Lemma~\ref{lem:twobatches} needs the number of pulls of $a$ in batch $\ell$ to be tight,
which is where the logarithmic counts are used. It extends item~1 of Appendix~A of the roadmap.

\begin{openquestion}[random batch ends]\label{oq:random}
Let the batch ends $\tl_1<\tl_2<\cdots$ be random, independent of the rewards and of the learner's randomization, with
$\log\tl_\ell/\log\tl_{\ell-1}\to x\in[1,\infty)$ almost surely, each end being revealed when it happens. Is
$x^2/\Ki(\nu_a,\m^\star)$ still the optimal first-order constant of $\E[\Ng_a(\tl_\ell)]/\log\tl_\ell$: a lower bound for every
consistent batched learner that knows the law of the grid, attained by \HY?
\end{openquestion}
\noindent\emph{Why it is open.} Theorem~\ref{thm:lower} fixes a deterministic grid, and the proof of
Theorem~\ref{thm:upper} uses deterministic ranks of the grid; with random ends both must be conditioned on the grid, and a
learner that knows the law of the grid may anticipate the end of a batch. It extends the setting of the roadmap (a new
direction).

%======================================================================
\section{Attribution}\label{sec:attribution}

\begin{center}
\begin{tabular}{p{0.40\textwidth}p{0.52\textwidth}}
\hline
Technique (where) & Earliest source found \\
\hline
Change of measure at the ends of the batches (Lemma~\ref{lem:onebatch}) & \cite{LR85} (to be checked); the form with
events at a stopping time \cite[Lemma~1]{KCG16}; the form with a $[0,1]$-valued variable \cite[Lemma~1 and (6)]{GMS19}.
Batched form, in which the observations available in a batch are those of the rounds before its start (``a product
distribution of at most $t_{J(t)-1}$ marginals''): \cite[Appendix~A]{PRCS16} (proof of its lower bound, Theorem~4, for a
fixed grid) and \cite[Section~3.1]{GHRZ19}; for static grids and anytime learners, \cite[Appendix~C]{JTXHXG21} \\
Tilting only the rewards of one batch towards an intermediate distribution, then the risk under the two neighbouring
alternatives (Lemmas~\ref{lem:com}, \ref{lem:typ} and~\ref{lem:twobatches}) & the expected sample size at a parameter
between the two hypotheses \cite{H60}; for multi-stage tests, the sample size after the penultimate stage forced by an
intermediate density, so that the last stage is negligible \cite[Theorem~3 and its proof]{L83}; a reference measure for
the observations before a batch \cite[Section~3.2, Lemmas~4--5 and Appendices~B.4--B.5]{GHRZ19}. New here, as far as we
found: consistency on both sides at the batch ends (an instance where $a$ is barely sub-optimal and one where it is barely
optimal), the boundary tilt of $\Ki$ mixed with atoms at $0$ and $1$ (for the mixing step itself, see two rows below), and the tightness of the number of pulls of $a$ in
batch $\ell$ \\
Boundary tilt from the optimality condition of the dual formula (Lemma~\ref{lem:tilts}(a)) & the optimal distribution of
the proof of \cite[Theorem~8]{HT10} \\
Mixing with a small atom at $0$ or at $1$ (Lemma~\ref{lem:tilts}(b)) & the mixture $(1-\alpha)\nu'+\alpha\delta_x$ with $x$
close to $1$, which raises the mean, in the proof of the regularity lemma of $\Ki$ \cite[Appendix~B.1]{GHMS22} (``a variation
on the proofs'' of Honda and Takemura); an atom at the right end also appears in the optimal distribution of the boundary case
of the proof of \cite[Theorem~8]{HT10}. New here: the use of both mixtures, with $\delta_0$ and with $\delta_1$, to obtain an
instance where $a$ is barely sub-optimal and one where it is barely optimal \\
Dual formula of $\Ki$ and its properties (Lemma~\ref{lem:kinf}) & \cite[(5) and Theorems~5--8]{HT10}; also
\cite{GHMS22} \\
Deviations of $n\Ki$ with a gain (Lemmas~\ref{lem:grid}--\ref{lem:upperdev}) & variants of \cite[Section~5,
Theorem~2]{CGMMS13}, \cite{HT15} and \cite[Section~2.3]{GHMS22}; the gain is this paper's \\
IMED index (IMED phase of \HY) & \cite[Section~2, Algorithm~1]{HT15} \\
Reward stacks (Lemma~\ref{lem:stacks}) & \cite[Section~4.6]{LS} (to be checked) \\
Paced targets $(1+\omega)\,c\log t$ tracked at every round (Definition~\ref{def:hy}) & \cite[Section~5,
Algorithm~1]{CMP17} (an arm is explored while its count is below $c(1+\gamma)\log t$); also \cite{LS17} (to be checked) \\
The queue lemma (Lemma~\ref{lem:service}) & closest known form: the tracking lemma \cite[Appendix~B.1, Lemma~15]{GK16}
(the rule that pulls an arm maximising the cumulated proportion minus its count keeps every count within $K-1$ of its
cumulated proportion), from which \cite[Lemma~7]{GK16} derives the forced exploration of C-tracking. Lemma~\ref{lem:service}
differs in three ways: its rule pulls the first arm below its target in a fixed order, its demands are integer targets that
grow like $\log t$ and are capped, and it has an initial burst of size $O(K\Ug_{\ell+1})$ \\
The plan made from the statistics frozen at $\tl_\ell$ towards the test two batches later (Definition~\ref{def:hy}) & source
not known \\
Water level of IMED rounds (Lemma~\ref{lem:W}) and high batches & the companion manuscript \cite{P26} (to be checked;
unpublished); source not known in the published literature \\
Consistency threshold (Lemma~\ref{lem:amb}) & source not known \\
\hline
\end{tabular}
\end{center}

\paragraph{Searches.} The conversion of 2026-11-14 ran no web search. The revision of 2026-11-14 read the sources above on
the web (the PDFs of HT10, GMS19 arXiv v3, KCG16, JTXHXG21 and its supplement, BK18, GHRZ19 arXiv v3, PRCS16 arXiv v3,
KO21b, JGS13 arXiv v2; the abstract of H60; the account of L83, Theorem~3, in Bartroff and Tartakovsky, arXiv:2403.18782,
Section~2.5), and reuses the queries of the scope report of 2026-11-14, all run on that date, with what they found:
``batched bandits imposed batch sizes fixed grid asymptotic lower bound Lai Robbins constant'' (PRCS16, GHRZ19, no constant
on imposed grids); ``stochastic bandit feedback revealed at exogenous batch ends unknown in advance asymptotically optimal
regret log t'' (EKMM21, nothing on unknown imposed ends); ``Kalkanli Ozgur asymptotic performance Thompson sampling batched
multi-armed bandits batch size condition'' (KO21b); ``bandits delayed feedback asymptotic lower bound regret constant
depends on delay growth polynomial delays Lai Robbins'' and ``bandit with delayed feedback delay proportional to time
t\^{}beta regret lower bound increases factor'' (constant or i.i.d.\ delays only); ``multi-armed bandit unknown horizon
doubly exponential checkpoints regret constant factor lower bound anytime doubling trick'' (BK18; JTXHXG21, Theorem~4);
``batched bandit exponential time grid regret constant b\^{}2 lower bound K\_inf IMED batched'' (no $b^2$ constant found);
``explore two batches ahead bandit lower bound grid'' (no hit); ``Prevost batched bandits imposed grid IMED K\_inf 2026
arXiv'' (no public version); ``Lorden 1983 asymptotic efficiency three-stage hypothesis tests last stage negligible'' (L83,
and its account in arXiv:2403.18782); ``Hoeffding 1960 lower bounds for the expected sample size and the average risk of a
sequential procedure intermediate distribution'' (H60); ``lower bound batched bandit change of measure only rewards of last
batch intermediate distribution'' (GHRZ19, Lemmas~2--5); ``Honda Takemura 2010 COLT bounded support models K\_inf dual
Theorem 8 Lemma 2'' (HT10; the numbering of HT15 and of arXiv:2111.09724); ``OSSB Combes Magureanu Proutiere exploration
rule (1+gamma) log t tracking'' (CMP17, Algorithm~1); ``adaptive treatment allocation grouped responses group sequential
multi-armed bandit asymptotic regret'' (nothing with a sharp constant); ``non-asymptotic regret bound batched bandit unknown
grid second order term explicit K\_inf'' (AJK21; nothing on unknown grids); ``batched bandits random batch sizes revealed
online'' and ``bandits feedback at random times renewal process batch ends'' (only KO21b); ``lower bound liminf expected
number of pulls consistent policy without assumption logarithmic regret batched or delayed'' (nothing on the batched
$\liminf$). The second revision (2026-11-14) reuses two queries of the second scope report, with what they found:
``K\_inf regularity proof mixture Dirac mass point x close to 1 increase mean continuity Honda Takemura'' (GHMS22,
Appendix~B.1, read on ar5iv) and ``round-robin scheduling lemma deficits of several arms C-tracking lemma Garivier
Kaufmann'' (GK16, Lemmas~7 and~15, read on ar5iv); it read the PMLR page of GK16 for its metadata.


%======================================================================
\section{Notation}\label{sec:notation}
%----------------------------------------------------------------------
\subsection{Vocabulary}\label{sec:app-vocab}
Every word below is a link to its definition, and so are its main uses in the statements and in the paragraphs ``Setting/Goal''.

\noindent \term{t:leader}{leader} $\Ld_\ell$\quad arm of largest $\tm_c(\tl_\ell)$; in batch $\ell+1$, preferred among the IMED minimizers, and pulled in the paced phase when no other arm has a demand\hfill{\small Def.~\ref{def:hy}}\par
\noindent \term{t:ref}{well-sampled}, \term{t:ref}{reference}, \term{t:ref}{test level}\quad $N_c\ge\frac12\Ug_\ell$; best well-sampled arm $\rr_\ell$; its mean $\mh_\ell$, against which the arms are tested\hfill{\small Def.~\ref{def:hy}}\par
\noindent \term{t:verdict}{rejected} / \term{t:verdict}{ambiguous}\quad verdict of the test at $\tl_\ell$: evidence $\Zg_\ell(c)$ above the threshold $\hg_\ell$, or not\hfill{\small Def.~\ref{def:hy}}\par
\noindent \term{t:net}{net}\quad the target $K-1+\Ug_{\ell+1}$ of an ambiguous arm, polylogarithmic\hfill{\small Def.~\ref{def:hy}}\par
\noindent \term{s:act}{IMED phase} / \term{s:act}{paced phase}\quad rounds $\le\sA_\ell$ of batch $\ell+1$ / rounds $>\sA_\ell$\hfill{\small Sec.~\ref{sec:algo}}\par
\noindent \term{s:act}{activated} batch\quad batch $\ell+1$ with paced rounds, $\tl_{\ell+1}>\sA_\ell$ ($\ell\in\Act$)\hfill{\small Sec.~\ref{sec:algo}}\par
\noindent \term{s:water}{water level} $\Lam_j$\quad smallest IMED index (divided by $\yg_{j-1}$) at the end of batch $j$\hfill{\small Sec.~\ref{sec:app-water-det}}\par
\noindent \term{s:evA}{good batch} ($x>1$)\quad batch $\ell$ on the event $\Gg_\ell$: accurate means, no wrong rejection, typical divergence of $a$\hfill{\small Def.~\ref{def:evA}}\par
\noindent \term{t:high}{$(g,L)$-high} batch ($x=1$)\quad batch $j<L$ whose water level is at least $(1-g)\log\tl_L$: every arm explored almost up to the horizon $\tl_L$\hfill{\small Sec.~\ref{sec:app-water}}\par
\noindent \term{s:evC}{classes} $\cU,\cG,\cC,\cO$ ($x=1$)\quad unreliable leader; good; $a$ looks competitive; other (at the end of the previous batch)\hfill{\small Def.~\ref{def:evC}}\par
\noindent count threshold, \term{s:ranksA}{$x>1$} / \term{s:constC}{$x=1$}\quad $\ug_\ell(\varepsilon)$ / $\up_j(\varepsilon)$: a count above which the means are accurate\hfill{\small Sec.~\ref{sec:app-constA}, Def.~\ref{def:constC}}\par

\subsection{Notation}\label{sec:app-notation}
Every symbol is a link to its definition.

\medskip\noindent \emph{Model}\par\smallskip
\noindent $\tl_\ell$, $\Bg_\ell$, $\Lg_\ell$, $\xg_\ell$\quad end and size of batch $\ell$, $\Lg_\ell:=\log\tl_\ell$, $\xg_\ell:=\Lg_\ell/\Lg_{\ell-1}$\hfill{\small Sec.~\ref{sec:setting}}\par
\noindent $\nu_c$, $\m_c$, $\m^\star$, $\Dg_c$, $\mtwo$, $\Dg$, $\cP$\quad arm distributions and means, best mean, gaps, second best mean, smallest gap, class of instances\hfill{\small Sec.~\ref{sec:setting}}\par
\noindent $a$, $a^\star$\quad the sub-optimal arm studied, a fixed optimal arm\hfill{\small Sec.~\ref{sec:setting}}\par
\noindent $\Ki$, $\kap$, $\Lset{\mu}$, $\da$\quad divergence (dual formula), $\kap(\mu):=\log\frac1{1-\mu}$, $[0,\frac1{1-\mu})$, $\da:=\Ki(\nu_a,\m^\star)$\hfill{\small Sec.~\ref{sec:setting}}\par
\noindent $\Ng_c(t)$, $\Reg(t)$\quad pulls of $c$ in rounds $1,\dots,t$ (included), regret\hfill{\small Sec.~\ref{sec:setting}}\par
\noindent $\Xs_{c,k}$, $\hnu_{c,n}$, $\hm_{c,n}$, $\hnu_c(t)$, $\hm_c(t)$\quad reward stacks, empirical distribution and mean of the first $n$ rewards, and at round $t$\hfill{\small Sec.~\ref{sec:setting}}\par
\noindent $\Mz$\quad information model: at round $t$ of batch $\ell+1$, $\tl_1,\dots,\tl_\ell$ and $t$ are known\hfill{\small Sec.~\ref{sec:setting}}\par
\medskip\noindent \emph{Algorithm \HY}\par\smallskip
\noindent $\wz$, $\beta$, $\om_\ell$, $\ody_\ell$, $\thv_\ell$, $\gam_\ell$\quad pseudo-count, net exponent, margins, dyadic margins, activation exponents, consistency margins\hfill{\small Def.~\ref{def:hy}}\par
\noindent $\yg_\ell$, $\Ug_{\ell+1}$, $\tauA_\ell$, $\sA_\ell$\quad weight $\xg_\ell^2$, net $\lceil\Lg_\ell^\beta\rceil$, activation round $\tl_\ell^{1+\thv_\ell}$, last IMED round $\lfloor\tauA_\ell\rfloor$\hfill{\small Def.~\ref{def:hy}}\par
\noindent $\tnu_c$, $\tm_c$, $\Ld_\ell$, $\tms_\ell$, $\Kt_\ell$, $\Jg_\ell$\quad statistics with the pseudo-observation, IMED leader and its mean, IMED divergences and index\hfill{\small Def.~\ref{def:hy}}\par
\noindent $\Wg_\ell$, $\rr_\ell$, $\mh_\ell$, $\Kh_\ell$\quad well-sampled arms, reference, test level, test divergences\hfill{\small Def.~\ref{def:hy}}\par
\noindent $\Qp_\ell$, $\Lh_\ell$, $\hg_\ell$, $\Zg_\ell$\quad poorly sampled arms except the previous leader, consistency level, threshold, evidence\hfill{\small Def.~\ref{def:hy}}\par
\noindent $\nt_{\ell,c}(t)$, $\Act$\quad paced target, set of the $\ell$ such that batch $\ell+1$ is activated\hfill{\small Def.~\ref{def:hy}, Sec.~\ref{sec:algo}}\par
\noindent $\HH_j$, $\HH(t)$\quad largest planned level up to batch $j$, resp.\ up to round $t$\hfill{\small Sec.~\ref{sec:app-constA}}\par
\medskip\noindent \emph{Tools and lower bound}\par\smallskip
\noindent $\Phis_\nu$, $\Psis_\nu$; $\taun_n$, $\lamk^{(k)}$, $\kn_n$\quad exponential moments of the tilt; grid of the dual formula\hfill{\small Sec.~\ref{sec:app-tools}}\par
\noindent $\sig_p$\quad elementary sums\hfill{\small Lemma~\ref{lem:sums}}\par
\noindent $\rhom$, $\rhop$, $\QQ$, $\nzero$, $\Mb$, $\Sl$, $\cB$\quad boundary instances, tilted measure and its log-likelihood ratios\hfill{\small Sec.~\ref{sec:app-lb-two}}\par
\medskip\noindent \emph{Deterministic properties}\par\smallskip
\noindent $\Lam_j$, $\lW$\quad water level of batch $j$, rank of Lemma~\ref{lem:W}\hfill{\small Sec.~\ref{sec:app-water-det}}\par
\medskip\noindent \emph{Case $x>1$}\par\smallskip
\noindent $\eb$, $\eg$, $\Ceta$, $\cst$, $\ew$, $\xp$, $\hb_\ell$\quad fixed constants: accuracy bound, gain, its constant, $\frac{2-\m^\star}{1-\m^\star}$, accuracy for the margin $\om$, $2x$, deterministic threshold\hfill{\small Def.~\ref{def:constA}}\par
\noindent $\De_\varepsilon$, $\Kp_\varepsilon$, $\kape_\varepsilon$, $\Rt_\varepsilon$, $\rt_\varepsilon$, $\ngr_\varepsilon$, $\cAc(\varepsilon)$, $\CLU(\varepsilon)$\quad constants depending on $\varepsilon$, of the deviation inequalities\hfill{\small Def.~\ref{def:constA}}\par
\noindent $\lA$, $\lone(\varepsilon,\varepsilon')$, $\ug_\ell(\varepsilon)$\quad ranks and count threshold\hfill{\small Sec.~\ref{sec:app-constA}}\par
\noindent $\Ag$, $\Vg$, $\Pig$, $\Lo$, $\Up$, $\Gg$, $\Fg$\quad events\hfill{\small Def.~\ref{def:evA}}\par
\noindent $\mt_\ell$, $\Yg_j$\quad plan of $a$ and nets of the decomposition\hfill{\small Lemma~\ref{lem:decompA}}\par
\noindent $\SE_\varepsilon$\quad $\sup_y(2\xp^3y-\varepsilon^2y^\beta)$\hfill{\small Lemma~\ref{lem:sumsA}}\par
\noindent $\epl_\ell$, $\cone$, $\etwo$, $\lonep$\quad accuracies and rank of the dyadic margins\hfill{\small Prop.~\ref{prop:dyA}}\par
\medskip\noindent \emph{Case $x=1$}\par\smallskip
\noindent $\ez$, $\ezz$, $\dl$, $\Cz$, $\Co$, $\yb$, $\thc$, $\ctl$, $\epth$\quad fixed constants\hfill{\small Def.~\ref{def:constC}}\par
\noindent $\Dp_\varepsilon$, $\Me_\varepsilon$, $\Cg_\varepsilon$, $\Rz_\varepsilon$, $\rz_\varepsilon$, $\up_j(\varepsilon)$\quad constants depending on $\varepsilon$\hfill{\small Def.~\ref{def:constC}}\par
\noindent $\lzp$, $\ly$, $\lonez$, $\lz$, $\lfour$, $\lfive$\quad ranks (functions of $\varepsilon$)\hfill{\small Sec.~\ref{sec:app-constC}, Def.~\ref{def:amb}}\par
\noindent $\Dev$, $\Vs$, $\Va$, $\cU$, $\cG$, $\cC$, $\cO$\quad events and classes\hfill{\small Def.~\ref{def:evC}}\par
\noindent $\Shi_j$, $\Sgo_j$\quad high batches, those with a well-sampled leader\hfill{\small Lemma~\ref{lem:goodhigh}}\par
\noindent $\zet_\ell$, $\zetp_\ell$, $\nl_\ell$, $\Acc_\ell$, $\KC_\ell$, $\Eamb_\ell$, $\cell_\ell$, $\Cfour(\varepsilon)$\quad objects of the consistency threshold\hfill{\small Def.~\ref{def:amb}, Lemma~\ref{lem:amb}}\par
\noindent $\Cpi(\varepsilon)$, $\RHo(\varepsilon)$\quad summed probabilities, remainder\hfill{\small Lemma~\ref{lem:sumC}, Prop.~\ref{prop:fixedC}}\par

\bibliographystyle{alpha}
% refs: the works cited; ../../papers/library: the library's papers (generated; keys INC-P<nnn>).
\bibliography{refs,../../papers/library}

\end{document}
